% Modernised reading — fonds Alexandre Grothendieck, shelfmark 12
% (pages 1-167, the whole folder).
%
% This is an INTERPRETATION, not a transcription. It re-reads the pages in
% today's notation and vocabulary, states the hypotheses the manuscript leaves
% implicit, and drops the critical apparatus so that the mathematics reads
% straight through. Everything the brackets and underlines carried moves into a
% footnote. It is held to being mathematically correct as it stands; where the
% manuscript is loose or mistaken, the true statement is what appears here and
% the footnote says what the page has. In any dispute of substance the
% transcription governs: whoever cites Grothendieck must cite batch-01.fr.tex
% (pp. 1-20) ... batch-09.fr.tex (pp. 161-167).
%
% Pass: Opus 5.5 (claude-opus-5-5), 2026-09-23 — first pass, unchecked against the pages by a human.
% The folder was read whole, all nine batches together; this is one document
% for the shelfmark. It was made on Opus 5.5 and has not been measured against
% the reference reading of folder 115 (made on Opus 5).
%
% Scope. Of the 167 leaves, about 125 carry his hand or a typescript of his.
% The rest (EGA IV galleys, a Bourbaki exposé, an English typescript on Brauer
% groups, an administrative note and a printer's cover dated March 1968) are
% versos he wrote on and are not mentioned. The March 1968 dates on the versos
% of pp. 116 and 118 are recorded here only; they are not a dating of the
% folder, whose \dating is Montpellier's and is copied from the batches.
%
% Spine. One programme, written out three times at three altitudes, plus the
% dossiers it spawned:
%   pp. 1-18    « Esquisse d'une théorie des motifs », his §§1-14 (no §11 on
%               the leaves): motives over a field, the Tannakian formalism,
%               the motivic Galois group, pi_1, Frobenius, the Hodge cocharacters,
%               Mumford-Tate, the Hodge conjecture as G = G_0.
%   p. 20       the 19-point plan « Motifs ».
%   pp. 22-79   the plan carried out: points 1-10 (pp. 22-40), then §§11-16
%               (pp. 41-79): coniveau, L(K) -> M^+(K) and birational invariants,
%               Galois characterisation of weight and level, commutation
%               theorems, the Lie-algebra motive, absolute cohomology and
%               absolute duality, abelian schemes and BSD.
%   pp. 81-94   loose leaves, then §17 positive forms (85-88) and §18
%               L-functions (89-94), the last two points of the plan but one.
%   pp. 95-121  « Poids et niveaux »: his typescript introduction (96-97), its
%               plan (98), and the working notes behind it: weights of the
%               l-adic cohomology of singular / open varieties (104-108),
%               admissible l-adic sheaves (110-119), the relative case (121).
%   pp. 120-122 two leaves of a typed report on Deligne's work; read here in
%               the order 122 -> 120 (see below).
%   pp. 124-136 « Divers types de Q_l-faisceaux » (Lemmes 1-3, Prop. 1).
%   pp. 138-141 smooth motives versus smooth isomotifs over a base.
%   pp. 142-166 the abstract construction: pseudo-abelian envelope, grading,
%               Tate object, level, duals, involutions, positive forms; logical
%               order 160 before 158.
%
% Moutures kept distinct: « niveau » three times (p. 2; D_i pp. 34-35;
% p. 154); positivity twice (§17 pp. 85-88, cones L^+(M) in fixed weight over
% an algebraically closed field; pp. 157-166, systems Phi_X of fundamental
% forms); Tate's conjecture twice (as an algebraic hull, pp. 11-12; as a
% double commutant, pp. 55-56).
%
% What only a folder-wide reading can say (ours, and said so in the body):
%   - the two typed leaves 120 and 122 are one text: p. 122 ends « qui se »,
%     p. 120 begins « une sorte de synthèse géométrico-arithmétique » — read
%     122 -> 120; the last word of 122 is overwritten, so this stays a reading.
%   - p. 77 (H^2(S, T_l(A)) = End(J) (x) Q_l =/= 0 for S a proper curve over
%     F_p) contradicts the « H^i(C,A) = 0 pour i >= 2 » of p. 81; p. 81's
%     conclusion survives without it.
%   - p. 92's « cf. N° 12 » is §12 of pp. 42-51; p. 121's « 2) » continues the
%     « (1) » of p. 115 and the « (2°) » of p. 119.
%   - the D_i of pp. 34-35 is the level of p. 2 extended to mixed motives, and
%     the xi of p. 36 is the zeta of pp. 43-57.
%
% NOTATION fixed for the folder:
%   - weights: Q(-1) is effective of weight 2 EVERYWHERE; M(n) = M (x) Q(n).
%     In the abstract part the Tate object is written Lambda(-1), of degree p
%     (p = 2 in geometry); pp. 152-154 write it Lambda(1), pp. 155-166
%     Lambda(-i): the first is renormalised (footnoted at p. 152).
%   - xi = cl Q(-1) = cl of the affine line, the Lefschetz class, for his xi
%     (p. 36) and his zeta (pp. 43-57); zeta is kept for zeta functions (p. 93).
%   - categories \mathcal{M}^+(X), \mathcal{M}(X); their Grothendieck rings
%     \mathbf{M}^+(X), \mathbf{M}(X) (plain M on the pages).
%   - H^i_c, Rf_! for his H^i_!, R^i_! f.
%   - y : G_m -> G the weight cocharacter (folder 13 writes i), i_1, i_2 the
%     Hodge cocharacters, eps : G -> G_m the Tate character (folder 10: j/tau).
%   - Z_G the centre of G (his figure-3 letter).
%   - « niveau » = level; « type dimensionnel » = level of a mixed motive.
%
% Substantive departures, each footnoted where it occurs:
%   p. 1    the equivalence M^{+1}(k) ~ Isab(k): with the page's two formulas
%           the functor is covariant; the exponent read ° is left open.
%   p. 2    M'_i((n-i)/2) -> M'_i((i-n)/2) (Q(-1) of weight 2).
%   p. 3    graded pieces « purs de poids i-2j » -> of level i-2j.
%   p. 8    (*) boxed as an isomorphism before the argument: canonical epi,
%           iso under Tate/Hodge.
%   p. 15   conj(i_1(lambda)) = i_2(conj lambda), i.e. i_2 = conj(i_1).
%   p. 26   R^{2d}f_*(1_X) = Z_S(-d) -> Q_S(-d); p. 27 purity needs the shift
%           [-2d].
%   p. 30   traces: duals exist only in M(X) (supplied).
%   p. 33   G_i ~ G_0 holds for i even only; G_i ~ G_1 for i odd.
%   pp. 38-39 xi has degree 2, not 1; M^pair = M_0[xi, xi^{-1}]; the table
%           D_i cap Sigma^+_j and the quotient D_i/D_{i-1} recomputed.
%   p. 36   « Ext^i = 0 » (doubtful reading) cannot hold: Kummer extensions.
%   p. 45   « dual des H^j »: by Lefschetz beta'_j = cl H^j.
%   pp. 48-49 the questions (isogenous jacobians; Gamma_i = 0 => unirational;
%           Gamma_n = 0 => ruled) have negative answers in char 0 too
%           (Riemann-Hurwitz; Enriques surfaces). « Cast. ok ! » overclaims.
%   p. 50   Gr_i -> Gr_alpha.
%   p. 52   effectivity by integrality: only => is known.
%   pp. 63-64, 70 absolute cohomology over F_p-bar: p. 63's vanishing plus
%           Leray contradicts the non-zero H^{2j+1} of pp. 64/70; left open.
%   p. 63   R^i -> R^{2i}.
%   p. 68   the reduction from Spec Z to its closed points is not formal.
%   p. 72   case i = 1: -2 <= rho <= 0, not <= 1.
%   p. 73   T_mot(A) = R^1f_*Q(1) has weight -1, not 1.
%   p. 81   H^1(C,A) finite (conj.), not 0; H^2(C,A) =/= 0 in general;
%           H^1(K_x,A) finite in its prime-to-p part only.
%   p. 86   B = Hom(M'(x)M'', Q(-p)), not Q(p).
%   p. 88   (iv) and (vii), not (vi); orthogonality of isotypic parts needs
%           M_alpha self-dual up to twist (supplied).
%   p. 90   L : R(N,k) -> 1 + t k[[t]] is not injective: nilpotent classes
%           die; injective on automorphisms.
%   p. 100  the answer to « inj ? » is no (computed).
%   p. 106  descent: proper surjective maps suffice (SGA 4 Vbis).
%   p. 108  C) 2): the page's bounds lack j; + 2j supplied (B's argument);
%           p. 119's « + j ? » left as his open question.
%   p. 111  Prop. 3.3 d): 2n - P = {2n - i}, twist (-n); p. 112 Prop. 4.2 d)
%           twist (-n).
%   p. 113  R^i f_!(F) is a direct summand (Deligne's decomposition), not
%           only a subquotient.
%   p. 132  Prop. 1 with « d+d'-1 » fails for F = Q_l; d+d' stated.
%   p. 139  « isomotif sur S » -> sur K; u_l -> alpha_l.
%   p. 147  (ii) is the normal form of every morphism, A (+) B -> B (+) C.
%   p. 154  K_{i,k} -> K_{i+j,k}.
%   p. 160  u^{(i)} : h^i(Y) -> h^i(X); p. 162 phi_{2i}; p. 164 h^i(Y).
%
% Announced and never established: point 19 of the plan (Hodge cohomology);
% §11 of the Esquisse; §14 « Structure de positivité » (title only, taken up
% pp. 85, 157); the explicit sigma^0, sigma^1 (p. 40); H(alpha) (p. 16);
% whether Ext^* is the Ext of motives (p. 62); the finiteness of H^1(S,A)(l)
% (box p. 74) and the Brauer surjectivity (p. 78); the programme of p. 94; §§1-9
% of the paper planned p. 98; the composition of p. 101; §6 of p. 114; Prop. 1
% of p. 132 in the case R^n j_*(Q_l) (« à traiter…!! »); the purity question
% and the « Pb » of p. 141; the verification of p. 166; the « partie
% algébrique » of the fundamental pro-motive (p. 58); the injectivity question
% of p. 51.

\documentclass[11pt,a4paper]{article}
% Le préambule commun à toutes les transcriptions du fonds Grothendieck.
%
% Il tient en une idée : **l'incertitude se compose, elle ne se résout pas.**
% Une transcription qui lisse un mot douteux a détruit la seule chose pour
% laquelle on la faisait — un lecteur doit pouvoir savoir, à chaque endroit,
% ce qui était sur le feuillet et ce qui a été ajouté. Les six macros qui
% suivent sont donc l'appareil critique, et non de la décoration.
%
% Les mêmes commandes sont reconnues par `scripts/render.mjs`, qui produit la
% vue de lecture du site. Une macro ajoutée ici doit l'être aussi là-bas,
% faute de quoi le rendu échouera bruyamment — ce qui est voulu : un rendu
% silencieusement faux est pire qu'un rendu absent.

\ProvidesPackage{grothendieck}[2026/08/08 Transcriptions du fonds Grothendieck]

\usepackage{amsmath,amssymb,amsthm}
\usepackage{mathrsfs}
\usepackage{stmaryrd}
% Les diagrammes commutatifs sont partout dans ce fonds : c'est la langue
% même de la géométrie algébrique de Grothendieck, et une transcription qui
% les rendrait en prose ne transcrirait rien.
\usepackage{tikz-cd}
% Français par défaut — c'est la langue des feuillets — mais l'anglais est
% chargé aussi : la traduction et le résumé sont des documents anglais, et la
% typographie française (espaces avant les deux-points) y serait une faute.
% Ces documents basculent par \selectlanguage{english} en tête de corps.
\usepackage[english,french]{babel}
\usepackage{fontspec}
\usepackage{geometry}
\geometry{a4paper,margin=25mm,marginparwidth=28mm}
\usepackage{marginnote}
\usepackage{xcolor}
% ulem, not soul. `soul`'s \st and \ul are built on a character-by-character
% scan that XeLaTeX's font machinery breaks: the document fails with an
% undefined control sequence rather than degrading. ulem does the same two jobs
% and is Unicode-engine safe. `normalem` stops it hijacking \emph, which the
% transcriptions use for Grothendieck's own emphasis.
\usepackage[normalem]{ulem}
\usepackage{hyperref}
\hypersetup{colorlinks=true,linkcolor=black,urlcolor=black,pdfborder={0 0 0}}

% --- Métadonnées du lot -----------------------------------------------------
% Reprises telles quelles par le script de rendu, qui les affiche en tête de
% la vue de lecture. Elles fixent ce que le fichier prétend couvrir : sans
% elles, une transcription flotte sans point d'attache dans un fonds de seize
% mille feuillets.

\newcommand{\folder}[1]{\gdef\grothFolder{#1}}
\newcommand{\batch}[1]{\gdef\grothBatch{#1}}
\newcommand{\leaves}[2]{\gdef\grothFirst{#1}\gdef\grothLast{#2}}
\newcommand{\foldertitle}[1]{\gdef\grothFoldertitle{#1}}
\gdef\grothFolder{?}\gdef\grothBatch{?}\gdef\grothFirst{?}\gdef\grothLast{?}\gdef\grothFoldertitle{}

% --- Le résumé --------------------------------------------------------------
%
% Il ouvre la lecture modernisée, sous le titre « Résumé », et il est composé
% en petit. Ce n'est pas un ornement : c'est le seul passage du document qui
% ne soit pas des mathématiques, et un lecteur doit pouvoir le reconnaître
% avant de l'avoir lu — puis le sauter s'il connaît déjà le sujet, ou s'y
% arrêter s'il ne le connaît pas. Le filet qui le suit marque la reprise du
% corps.

\newenvironment{resume}
  {\small\setlength{\parskip}{0.45em}}
  {\par\medskip\hrule\medskip}

% --- Mots-clés ---------------------------------------------------------------
%
% En anglais, en clôture du résumé de la lecture modernisée : le vocabulaire
% moderne sous lequel chercher ce que les feuillets construisent. Le manifeste
% du site (`npm run manifest`) les extrait pour étiqueter la cote entière —
% cette ligne est la source unique des tags, il n'y a pas de fichier à côté.
\newcommand{\keywords}[1]{\par\smallskip\noindent
  {\footnotesize\textsc{Keywords} --- \textit{#1}}\par}

% --- L'appareil critique ----------------------------------------------------

% Le numéro de feuillet, en marge. C'est l'ancre qui permet de régler un
% désaccord sur une formule en regardant *un* feuillet plutôt qu'une chemise
% de six cents. Le site s'en sert aussi pour tourner les pages du fac-similé
% quand on fait défiler la transcription.
\newcommand{\leaf}[1]{%
  \marginnote{\footnotesize\color{gray!70}\textsf{#1}}%
  \label{leaf:#1}\ignorespaces}

% Illisible : on ne devine pas. Un mot inventé qui se lit comme les autres est
% le pire résultat possible.
\newcommand{\ill}{\textcolor{red!60!black}{[\,\ldots\,]}}

% Lecture douteuse : le mot est donné, et signalé comme incertain.
\newcommand{\uncertain}[1]{\uline{#1}}

% Ajout de l'éditeur — une lettre manquante, un mot sous-entendu.
\newcommand{\add}[1]{\textcolor{blue!50!black}{[#1]}}

% Biffé par Grothendieck lui-même. Ce qu'il a rayé fait partie du document :
% une rature dit souvent où la pensée a changé de direction.
\newcommand{\struck}[1]{\sout{#1}}

% Note du transcripteur, sur l'état matériel du feuillet.
\newcommand{\note}[1]{\par\smallskip{\small\color{gray!60!black}\textit{[#1]}}\par\smallskip}

% Note marginale de Grothendieck, distincte du corps du texte.
\newcommand{\marginal}[1]{\par\smallskip\noindent\hspace{1em}%
  {\small\color{gray!60!black}#1}\par\smallskip}

% --- Titre ------------------------------------------------------------------

\AtBeginDocument{%
  \begin{center}
    {\large\bfseries Fonds Alexandre Grothendieck --- cote n\textsuperscript{o}~\grothFolder}\\[2pt]
    {\itshape\grothFoldertitle}\\[4pt]
    {\small feuillets \grothFirst--\grothLast\ (lot \grothBatch)}\\[2pt]
    {\footnotesize\color{gray!60!black}Transcription établie d'après le fac-similé numérisé
     par l'Université de Montpellier.\\
     Les lectures incertaines sont soulignées, les illisibles notées [\,\ldots\,],
     les ajouts entre crochets.}
  \end{center}
  \vspace{1em}}

\endinput


\folder{12}
\pages{1}{167}
\dating{1964-[vers 1971]}
\watermark{Édition de démonstration\\interprétation personnelle de l'œuvre}
\foldertitle{Généralités (sous-groupes de Galois motiviques) (1964 ?) : notes manuscrites (s.d.), tapuscrit (s.d.) — lecture modernisée du dossier entier}

\begin{document}

\section*{Résumé}

\begin{resume}
Ces feuillets cherchent ce qu'une variété algébrique a de commun avec toutes
ses cohomologies à la fois, et proposent de lui donner un nom : son
\emph{motif}.

Voici le problème. À une variété définie par des équations, on sait associer
plusieurs espaces vectoriels qui en mesurent la forme : la cohomologie de Betti
si l'on travaille sur les nombres complexes, la cohomologie de De Rham faite de
formes différentielles, et, pour chaque nombre premier $\ell$, une cohomologie
$\ell$-adique qui a l'avantage de fonctionner sur n'importe quel corps et de
porter une action du groupe de Galois. Ces théories donnent les mêmes nombres
--- les mêmes dimensions, les mêmes traces des endomorphismes --- et pourtant
aucune ne se déduit des autres : elles ne vivent pas sur le même corps de
coefficients. Tout se passe comme si une même chose était photographiée sous
plusieurs éclairages, sans qu'on ait jamais la chose elle-même.

L'idée qui ouvre le dossier est de construire cette chose. Un motif serait
un objet d'une catégorie où l'on sait additionner, prendre des facteurs
directs, faire des produits tensoriels ; chaque cohomologie deviendrait un
foncteur, une « réalisation », de cette catégorie vers des espaces vectoriels.
L'analogie est celle d'un groupe et de ses représentations : de même qu'un
groupe se reconstitue à partir de la catégorie de ses représentations et du
foncteur qui oublie l'action, la catégorie des motifs munie d'une réalisation
se trouve être la catégorie des représentations d'un groupe. Ce groupe, le
\emph{groupe de Galois motivique}, généralise le groupe de Galois ordinaire,
qui doit en être le quotient
par la composante neutre, et le premier cahier du dossier montre comment
les grandes conjectures du moment --- celle de Tate, celle de Hodge --- se
disent toutes deux comme une seule phrase sur la taille de ce groupe.

Mais le dossier ne s'arrête pas à la définition. Il fait travailler l'objet
qu'il n'a pas encore. Sur un corps fini, la fonction zêta d'une variété
se lit dans les valeurs propres du Frobenius, et leurs valeurs absolues ---
les \emph{poids} --- découpent la cohomologie en morceaux. Grothendieck
s'aperçoit que ces poids, et une notion voisine qu'il appelle le
\emph{niveau}, ordonnent tout : quels morceaux de la cohomologie d'une
variété singulière ou non compacte peuvent apparaître, quels invariants
restent quand on remplace la variété par une autre qui lui est
birationnelle, ce que la fonction $L$ sait de la cohomologie. Il en tire
un article, dont l'introduction dactylographiée est ici : il y dit que
l'idée de ces raisonnements par les poids « remonte à 1964 », qu'ils ne sont
pour l'instant qu'heuristiques, et qu'ils méritent pourtant « de ne pas être
dédaignés ». On voit à l'œuvre, sur cent soixante pages, une méthode qui consiste
à raisonner avec rigueur sur un objet dont l'existence est conjecturale, pour
découvrir ce qui doit être vrai des objets qui existent.

Le dossier contient enfin trois choses qui ont eu une postérité précise. Une
construction abstraite d'une catégorie de motifs --- on part des variétés,
on remplace les applications par les correspondances algébriques, puis on
ajoute formellement les images des projecteurs --- qui est exactement la
construction qu'on enseigne aujourd'hui. Une théorie des formes positives
sur les motifs, où les algèbres d'endomorphismes deviennent des algèbres
munies d'une involution positive, comme celles des variétés abéliennes.
Et, sur deux feuillets dactylographiés, un rapport sur les premiers travaux
de Deligne, où Grothendieck raconte qu'il lui avait communiqué comme une
conjecture l'énoncé que Deligne a démontré.

Presque tout ici est conditionnel : aux conjectures standard, à celles de
Weil, de Tate, de Hodge, à la résolution des singularités. Le dossier le sait
et le dit, et c'est ce qui le rend lisible aujourd'hui, où une partie de ces
conditions a été levée et une autre non.

Les noms sous lesquels on cherchera la suite : motifs purs et mixtes,
catégories tannakiennes et groupe de Galois motivique, conjectures de Tate et
de Hodge, groupe de Mumford--Tate, filtration par le poids et filtration par
le coniveau, cohomologie motivique, anneau de Grothendieck des variétés,
conjectures standard, conjecture de Birch et Swinnerton-Dyer.

\keywords{pure motive, effective motive, mixed motive, Tate motive, Lefschetz motive, realisation functor, l-adic realisation, Tannakian category, fibre functor, motivic Galois group, torsor of fibre functors, gerbe, Artin motive, etale fundamental group, Frobenius element, Tate conjecture, semisimplicity conjecture, Hodge conjecture, Mumford-Tate group, weight cocharacter, Hodge cocharacters, filtered fibre functor, supersingular elliptic curve, level of a motive, coniveau filtration, generalized Hodge conjecture, coniveau spectral sequence, six functor formalism, projection formula, independence of l, lambda-ring, Grothendieck ring of motives, Grothendieck ring of varieties, motivic Euler characteristic, virtual Betti numbers, birational invariants, stable rationality, unirationality, ruled varieties, double commutant theorem, Lie algebra of the motivic Galois group, motivic cohomology, absolute cohomology, Ext of mixed motives, arithmetic duality, Hochschild-Serre spectral sequence, Mordell-Weil theorem, Kummer sequence, Tate-Shafarevich group, Lang's theorem, Brauer group, Artin's conjecture on the Brauer group, Birch and Swinnerton-Dyer conjecture, positive cone, Weil form, L-function, big Witt vectors, Lefschetz trace formula, Weil conjectures, weight filtration, weights of l-adic cohomology, simple normal crossings divisor, cohomological descent, relative hard Lefschetz, Deligne's degeneration theorem, lambda-structure on K-theory of perfect complexes, constructible l-adic sheaf, absolute purity, Artin's elementary fibrations, smooth motive, integral structure on motives, good reduction, Neron-Ogg-Shafarevich criterion, pseudo-abelian envelope, Karoubi envelope, semisimple abelian category, category of correspondences, homological equivalence, numerical equivalence, standard conjecture D, Kunneth projectors, Koszul sign, primitive decomposition, rigid tensor category, internal Hom, Rosati involution, positive involution, Riemann algebra, Hodge standard conjecture}
\end{resume}

\pagerange{1}{167}
\section*{Le fil du dossier, et les conventions}

\subsection*{Les stations}

Le dossier écrit trois fois le même programme, à trois altitudes, puis ouvre
les chantiers que ce programme appelle.

\begin{enumerate}
  \item \emph{Pages 1 à 18}, « Esquisse d'une théorie des motifs », ses
  paragraphes 1 à 14 paginés par lui de 1 à 17 (la page 10 est un feuillet
  blanc, il n'y a pas de paragraphe 11) : les motifs sur un corps, leurs
  réalisations, le formalisme tannakien, le groupe de Galois motivique, son
  lien avec le groupe fondamental, le Frobenius, les structures de Hodge, le
  groupe de Mumford--Tate et la conjecture de Hodge lue comme une égalité de
  groupes. Le paragraphe 14 n'a que son titre.
  \item \emph{Page 20}, un plan en dix-neuf points, « Motifs ».
  \item \emph{Pages 22 à 40}, les points 1 à 10 de ce plan, un à un : motifs
  sur une base, opérations, réalisation $\ell$-adique, objet de Tate, traces,
  filtration par le poids et par le « type dimensionnel », motifs constants
  tordus et anneaux de Grothendieck.
  \item \emph{Pages 41 à 79}, ses paragraphes 11 à 16 : le coniveau,
  l'homomorphisme $L(K) \to \mathbf{M}^{+}(K)$ et les invariants birationnels,
  la caractérisation galoisienne des filtrations, les théorèmes de commutation,
  le motif de Lie, la cohomologie absolue et sa dualité, les schémas abéliens
  et la conjecture de Birch et Swinnerton-Dyer.
  \item \emph{Pages 81 à 94} : deux feuillets isolés, puis ses paragraphes 17
  (formes positives) et 18 (fonctions $L$).
  \item \emph{Pages 95 à 121}, « Poids et niveaux d'espaces de cohomologie
  $\ell$-adiques » : l'introduction dactylographiée d'un article, son plan, et
  les notes de travail qui le portent --- poids de la cohomologie des variétés
  singulières ou non propres, faisceaux $\ell$-adiques admissibles, cas
  relatif.
  \item \emph{Pages 120 et 122}, deux feuillets d'un rapport dactylographié
  sur les travaux de Deligne, qu'on lit ici dans l'ordre 122 puis 120.
  \item \emph{Pages 124 à 136}, « Divers types de $\mathbb{Q}_{\ell}$-faisceaux ».
  \item \emph{Pages 138 à 141}, motifs lisses et isomotifs lisses sur une base.
  \item \emph{Pages 142 à 166}, la construction abstraite des catégories de
  motifs : enveloppe pseudo-abélienne, graduation, objet de Tate, niveau,
  duaux, involutions, formes définies positives.
\end{enumerate}

\subsection*{Les moutures, gardées distinctes}

Trois notions reviennent sous plusieurs formes, et la différence est du
contenu. Le \emph{niveau} est défini trois fois : pour les motifs semi-simples
sur un corps (page 2), comme filtration par le « type dimensionnel » d'une
catégorie qui n'est plus semi-simple (pages 34 et 35), et dans le cadre abstrait
(page 154). La \emph{positivité} est traitée deux fois : par des cônes de formes
en poids fixé sur un corps algébriquement clos (pages 85 à 88), puis par des
systèmes de « formes fondamentales » sur toutes les variétés à la fois (pages
157 à 166). La \emph{conjecture de Tate} est dite deux fois : comme l'égalité
du groupe de Galois motivique avec l'enveloppe algébrique de l'image de Galois
(pages 11 et 12), et comme un théorème du double commutant (pages 55 et 56).

\subsection*{Ce que seule une lecture d'ensemble peut dire}

Quatre liens, qu'aucun feuillet n'énonce, et qui sont de nous.

Les deux feuillets dactylographiés 120 et 122 sont un seul texte. La page 122
s'arrête sur « la théorie des motifs, qui se » ; la page 120 commence par « une
sorte de synthèse géométrico-arithmétique des nombreuses théories
cohomologiques ». La phrase s'enchaîne, et la page 122, qui est le point 4 du
rapport, précède la page 120, qui en est le point 5. Le dernier mot de la page
122 est surchargé à la main, de sorte que ce raccord reste une lecture.

La page 77 contredit la page 81. La première établit que, pour une courbe
propre et lisse $S$ sur $\mathbb{F}_{p}$ et $A = J \times S$, le groupe
$H^{2}(S, T_{\ell}(A))$ est $\operatorname{End}(J) \otimes \mathbb{Q}_{\ell}
\neq 0$ ; la seconde affirme que $H^{i}(C, A) = 0$ pour $i \geqslant 2$ et
toute courbe propre $C$ sur un corps fini. C'est la page 77 qui a raison, et la
conclusion de la page 81 survit sans l'énoncé fautif.

La page 92 renvoie à son « N° 12 » : c'est le paragraphe 12 des pages 42 à 51,
les invariants birationnels fondamentaux, que la page 94 nomme d'ailleurs
comme « envisagés dans les notes sur les \emph{Motifs} ». Et le « 2) » de la
page 121 est la suite du « (1) » de la page 115 et du « (2°) » de la page 119.

Enfin, la filtration $D_{i}$ des pages 34 et 35 est le niveau de la page 2,
étendu aux motifs qui ne sont plus semi-simples ; et la lettre $\xi$ de la page
36 et la lettre $\zeta$ des pages 43 à 57 désignent le même élément, la classe
du motif de Lefschetz.

\subsection*{Les conventions, valables pour tout le dossier}

\begin{itemize}
  \item Le motif de Tate $\mathbb{Q}(-1)$ est effectif et de poids $2$ ; on
  écrit $M(n) = M \otimes \mathbb{Q}(n)$, de sorte que $M(-1)$ a le poids de
  $M$ augmenté de $2$. C'est la convention des pages 15, 27 et 31. Dans le cadre
  abstrait des pages 143 à 166, l'objet de Tate est noté $\Lambda(-1)$, de degré
  $p$ ($p = 2$ en géométrie) ; les pages 152 à 154 l'écrivent $\Lambda(1)$ et les
  pages 155 à 166 $\Lambda(-i)$ : on a renormalisé les premières.
  \item $\xi = \operatorname{cl}\mathbb{Q}(-1)$ est la classe de Lefschetz, qui
  est aussi l'image de la droite affine ; on l'écrit $\xi$ partout, là où les
  pages 43 à 57 écrivent $\zeta$. La lettre $\zeta$ est réservée aux fonctions
  zêta (page 93).
  \item Les catégories de motifs sont $\mathcal{M}^{+}(X)$ (effectifs) et
  $\mathcal{M}(X)$ ; leurs anneaux de Grothendieck sont notés
  $\mathbf{M}^{+}(X)$ et $\mathbf{M}(X)$, là où les pages écrivent un $M$ droit.
  \item On écrit $H^{i}_{c}$ et $Rf_{!}$ la cohomologie à supports propres, là
  où il écrit $H^{i}_{!}$ et $R^{i}_{!}f$.
  \item $G$ est le groupe de Galois motivique d'une sous-catégorie tensorielle
  munie d'un foncteur fibre, $Z_{G}$ son centre, $y \colon \mathbb{G}_{m} \to G$
  le cocaractère des poids, $i_{1}$, $i_{2}$ les cocaractères de Hodge,
  $\varepsilon \colon G \to \mathbb{G}_{m}$ le caractère défini par
  $\mathbb{Q}(-1)$.\footnote{Le dossier 13, de la même époque, écrit $i$ pour ce
  que ce dossier note $y$ ; le dossier 10 écrit $j$ ou $\tau$ pour
  $\varepsilon$. Les lettres sont celles de chaque dossier ; les objets sont les
  mêmes.}
  \item « Niveau » et « type dimensionnel » désignent la même notion ; on dit
  \emph{niveau}. Un motif est \emph{bipur} de type $i$ s'il est pur de poids $i$
  et de niveau $i$.
\end{itemize}

\subsection*{Ce que le dossier suppose sans l'établir, et ce qu'il annonce sans
le faire}

Il suppose, presque partout, l'existence d'une catégorie de motifs qui soit
tannakienne et semi-simple sur les corps --- ce que les conjectures standard
donneraient --- et, selon les endroits, les conjectures de Weil, de Tate (avec
la semi-simplicité de l'action de Galois), de Hodge, et la résolution des
singularités. On le rappelle à chaque énoncé qui en dépend.

Il annonce et ne fait pas : le point 19 du plan (la relation avec la
cohomologie de Hodge) ; le paragraphe 14 de l'« Esquisse », dont il n'y a que
le titre, repris plus tard par les pages 85 et 157 ; la détermination explicite
des ensembles $\sigma_{0}$ et $\sigma_{1}$ promise page 40 ; la construction du
groupe $H(\alpha)$ (page 16) ; la question de savoir si les $\operatorname{Ext}$
de la cohomologie absolue sont ceux de la catégorie des motifs (page 62) ; la
finitude de $H^{1}(S, A)(\ell)$ (page 74) et la surjectivité sur les groupes de
Brauer (page 78) ; le programme de la page 94 ; les paragraphes 1 à 9 de
l'article planifié page 98 ; la composition annoncée page 101 ; le paragraphe 6
de la page 114 ; le cas de la Proposition 1 de la page 132 qu'il laisse « à
traiter…!! » ; la question de pureté et le « Pb » de la page 141 ; la
vérification de la page 166 ; la « partie algébrique » du pro-motif
fondamental (page 58).

\pagerange{1}{18}
\section*{I. Esquisse d'une théorie des motifs (pages 1 à 18)}

Dix-sept feuillets de sa main, sur papier à en-tête de l'IHES, paginés par lui
de 1 à 17 --- la page 10 est un feuillet blanc, de sorte que ses pages 10 à 17
sont les pages 11 à 18 du dossier --- sous le titre souligné « Esquisse d'une
théorie des motifs ». Rien n'y est daté. C'est le texte le plus ramassé du
dossier : un programme dont chaque paragraphe tient en une page.

\pagerange{1}{1}
\subsection*{Le programme, les motifs effectifs, le motif de Tate (page 1)}

Le motif doit remplacer, pour les schémas projectifs et lisses sur un corps
$k$, la cohomologie tout entière : il est à la cohomologie ce que la théorie
des modules de type fini sur $\mathbb{Z}$ est à la théorie des espaces
vectoriels sur chacun des corps premiers. Le résumé qu'il en fait tient en
trois mots : isomotifs, motifs semi-simples, base ponctuelle.

Il pose d'abord une catégorie $\mathcal{M}^{+}(k)$ de \emph{motifs effectifs}
sur $k$, abélienne, semi-simple, munie d'un produit tensoriel, et des
foncteurs de réalisation, tous multiplicatifs et fidèles : $T_{DR}$ et
$T_{\mathrm{Hdg}}$ si $k$ est de caractéristique nulle, $T_{B}$ si
$k = \mathbb{C}$, $T^{\ell}$ pour $\ell \neq \operatorname{car} k$. La
catégorie est graduée, compatiblement au produit tensoriel, et chaque variété
$X$ donne ses composantes $h^{i}(X)$. En degré $1$, les motifs effectifs sont
les variétés abéliennes à isogénie près : $X \mapsto
\operatorname{Pic}^{0}_{X}$ d'un côté, $A \mapsto h^{1}(\hat{A})$ de l'autre,
où $\hat{A}$ est la variété duale.\footnote{La page écrit
$\mathcal{M}^{+1}(k)^{\circ} \xrightarrow{\sim} \mathrm{Isab}(k)$ et lit
l'exposant de $\operatorname{Pic}$ « $00$ », sous réserve. Avec les deux
formules qu'elle donne, le foncteur $h^{1}(X) \mapsto \operatorname{Pic}^{0}_{X}$
est \emph{covariant} --- deux contravariances se composent --- et son
quasi-inverse est bien $A \mapsto h^{1}(\hat{A})$, puisque $h^{1}(X) \simeq
h^{1}(\operatorname{Alb} X)$ et que $\operatorname{Alb} X$ est la duale de
$\operatorname{Pic}^{0}_{X}$. Si le « $\circ$ » désigne la catégorie opposée,
c'est l'autre convention, $A \mapsto h^{1}(A)$ ; la transcription ne permet pas
de trancher ce que désigne le signe.} Il note que les motifs de poids
$\geqslant 1$ généralisent les variétés abéliennes, et qu'on garde la
« connexion discrète » : pour une extension $k \to k'$ de corps
algébriquement clos, le foncteur de changement de base est pleinement fidèle.

Vient alors le motif $\mathbb{Q}(-1)$, qu'il appelle « anti-Tate », et la
catégorie $\mathcal{M}(k)$ des motifs, obtenue en inversant $\mathbb{Q}(-1)$ :
elle a, en plus du produit tensoriel et de la graduation, des $\operatorname{Hom}$
internes. Et elle se caractérise par une propriété universelle. Soit
$\mathcal{C}$ une catégorie pseudo-abélienne munie d'un produit tensoriel. Les
foncteurs tensoriels $\mathcal{M}(k) \to \mathcal{C}$ correspondent aux
foncteurs $T \colon \mathcal{V}(k)^{\circ} \to \mathcal{C}$ sur la catégorie
opposée des variétés projectives et lisses qui satisfont la formule de Künneth
et pour lesquels l'image $T(\mathbb{Q}(-1)) = \operatorname{Im} T(\pi)$ du
projecteur $\pi$ de $h(\mathbb{P}^{1}_{k})$ d'image $h^{2}(\mathbb{P}^{1}_{k})$
est inversible ; ces $T$ commutent alors automatiquement aux $\operatorname{Hom}$.
C'est, au mot près, la propriété universelle qu'on attache aujourd'hui aux
motifs purs comme source des cohomologies de Weil.\footnote{La page ne dit pas
à quelle relation d'équivalence sur les cycles la catégorie est construite ; la
propriété universelle telle qu'elle l'énonce n'a de sens qu'une fois la
catégorie construite, ce que les pages 143 à 151 feront.}

\pagerange{2}{3}
\subsection*{Le niveau d'un motif (pages 2 et 3)}

Un motif $M$ est \emph{de niveau} $\leqslant i$ s'il est de la forme $M'(n)$ avec
$M'$ effectif à composantes homogènes de degrés $\leqslant i$ ; il est
\emph{purement de niveau} $i$ s'il est de niveau $\leqslant i$ et n'a aucun
sous-motif non nul de niveau $\leqslant i-1$. Cet $i$ est unique et positif,
et tout motif se décompose de façon unique en somme de motifs purs de niveau
$i$. Si $M$ est de poids $n$, les $i$ qui interviennent ont la parité de $n$,
et
\[
M \;=\; \bigoplus_{i} M'_{i}\Big(\frac{i-n}{2}\Big) ,
\qquad M'_{i} \text{ effectif, pur de poids et de niveau } i .
\]\footnote{La page écrit $M'_{i}\big(\frac{n-i}{2}\big)$. Avec sa propre
convention, où $\mathbb{Q}(-1)$ est de poids $2$ et $M(-1) = M \otimes
\mathbb{Q}(-1)$ (pages 15 et 27), il faut tordre $M'_{i}$, de poids $i$, par
$\mathbb{Q}(-\frac{n-i}{2})$ pour atteindre le poids $n$ ; le signe est
rétabli. Il écrit aussi la somme avec $\coprod$.}

Les motifs de niveau $\leqslant i$ sont ceux qui, à torsion de Tate près,
proviennent de variétés de dimension $\leqslant i$ : $h(X)$ est de niveau
$\leqslant \dim X$, et si $i \leqslant \dim X$, $h^{i}(X)$ s'injecte dans
$h^{i}(Y)$ pour $Y$ une section linéaire de dimension $i$ --- c'est le théorème
de Lefschetz faible. Si $i \geqslant n = \dim X$, $h^{i}(X)$ est de niveau
$\leqslant 2n-i$, par le théorème de Lefschetz fort.

Il en donne l'interprétation $\ell$-adique. Soit la « filtration géométrique »
\[
\operatorname{Filt}^{j} H^{i}(\overline{X}, \mathbb{Q}_{\ell})
\;=\; \varinjlim_{U}\, \operatorname{Ker}\bigl(H^{i}(\overline{X}, \mathbb{Q}_{\ell})
\to H^{i}(\overline{U}, \mathbb{Q}_{\ell})\bigr) ,
\qquad \operatorname{codim}(X - U) \geqslant j ,
\]
qu'on appelle aujourd'hui la \emph{filtration par le coniveau}. Il attend que
$\operatorname{Filt}^{j} H^{i}$ soit la réalisation du plus grand sous-motif de
$h^{i}(X)$ de niveau $\leqslant i - 2j$. Moyennant la résolution des
singularités et les conjectures standard, les dimensions des
$\operatorname{Filt}^{j}$ seraient alors indépendantes de $\ell$, et les
quotients successifs seraient des motifs de niveau pur $i-2j$.\footnote{La page
dit ces quotients « purs de poids $i-2j$ resp.\ $2n-i-j$ ». Ils sont de poids
$i$, comme $H^{i}$ lui-même ; ce qui vaut $i-2j$ est leur niveau, et c'est
ainsi qu'on l'écrit. La variante « $2n-i-j$ », qui semble viser l'autre moitié
des degrés, n'a pas été reconstituée. L'énoncé qu'on attend ici --- le
coniveau comme niveau motivique --- est la forme $\ell$-adique de ce qu'on
appelle la \emph{conjecture de Hodge généralisée}, qu'il reformulera en 1969 ;
elle est ouverte.} Il ajoute : « parler des invariants birationnels », et, en
marge, que les conjectures de Tate et de Hodge pourraient venir ici. Les pages
41 à 51 le feront.

\pagerange{3}{5}
\subsection*{Les $\otimes$-catégories (pages 3 à 5)}

Une catégorie abélienne munie d'un produit tensoriel associatif, commutatif et
unitaire, avec des $\operatorname{Hom}$ internes, est dite \emph{rigide} si
$\check{M} \otimes N \xrightarrow{\sim} \underline{\operatorname{Hom}}(M,N)$ ;
elle est une $\otimes$-catégorie sur un corps $K$ si elle est $K$-linéaire à
$\operatorname{Hom}$ de dimension finie, et \emph{connexe} si
$\operatorname{End}(\mathbf{1}) = K$.\footnote{Il écrit tantôt $K$, tantôt $k$
pour le corps de coefficients ; on écrit $K$.}

L'exemple universel est la catégorie $\operatorname{Rep}_{K}(G)$ des
représentations de dimension finie d'un schéma en groupes affine $G$ sur $K$.
Étant donnée une $\otimes$-catégorie $\mathcal{C}$ munie d'un \emph{foncteur
fibre} $F$ à valeurs dans les $K$-espaces vectoriels, on lui associe
canoniquement $G = \underline{\operatorname{Aut}}^{\otimes}(F)$, défini par
$G(A) = \operatorname{Aut}^{\otimes}(F \otimes A)$, et une équivalence
\[
(\mathcal{C}, F) \;\simeq\; (\operatorname{Rep}_{K}(G), \text{oubli}) ;
\]
le foncteur qui envoie $G$ sur $(\operatorname{Rep}_{K}(G), \text{oubli})$ est
pleinement fidèle. $\mathcal{C}$ est de type fini si et seulement si $G$ l'est ;
en caractéristique nulle, $\mathcal{C}$ est semi-simple si et seulement si $G$
est pro-réductif.\footnote{C'est le théorème de reconstruction tannakien, sous
la forme qu'on lit aujourd'hui dans la thèse de Saavedra (1972) et dans
Deligne--Milne (1982). Les feuillets l'énoncent sans démonstration et sans le
mot « tannakien ».}

On en tire les autres opérations tensorielles --- puissances extérieures et
symétriques --- indépendamment du choix de $F$. Tout autre foncteur fibre $F'$
s'obtient en tordant $F$ par un torseur $P$ sous $G$, $F' = P \times^{G} F$, et
son groupe est la forme intérieure $G' = P \times^{G} G = \operatorname{ad}(P)$.
Deux foncteurs fibres sur une catégorie de type fini deviennent isomorphes sur
une extension finie. Il peut n'exister aucun foncteur fibre ; mais sur une
catégorie de type fini, il « doit pouvoir » s'en trouver un sur une extension
finie, et l'on décrit alors $\mathcal{C}$ intrinsèquement par une
\emph{gerbe}.\footnote{L'existence d'un foncteur fibre sur une extension, pour
les catégories tensorielles rigides en caractéristique nulle dont les dimensions
sont des entiers positifs, est un théorème de Deligne (1990) ; la description
par une gerbe est celle de Saavedra. La page dit « doit pouvoir montrer » :
c'est une attente.}

\pagerange{5}{6}
\subsection*{Motifs sur une base, et le groupe de Galois motivique (pages 5 et 6)}

Pour $S$ régulier et connexe, de point générique $\eta$, on définit
$\mathcal{M}^{+}(S)$ et $\mathcal{M}(S)$ comme sur un corps, à partir des
schémas projectifs et lisses sur $S$, mais en prenant pour correspondances
celles de la fibre générique, $C(X) := C(X_{\eta})$. Le foncteur $\mathcal{M}(S)
\to \mathcal{M}(\eta)$ est donc pleinement fidèle par construction, et un
morphisme $S' \to S$ de schémas réguliers doit donner un foncteur tensoriel
$f^{*}$ --- c'est là que sert la régularité, et il note « à expliciter ».

Tout point géométrique complexe $\xi$ de $S$ définit un foncteur fibre de Betti
$F^{B}_{\xi}$ à valeurs dans les $\mathbb{Q}$-espaces vectoriels, d'où un groupe
pro-algébrique pro-réductif $G$ sur $\mathbb{Q}$ qui décrit $\mathcal{M}(S)$ :
c'est le \emph{groupe de Galois motivique}. En pratique on ne considère qu'une
sous-catégorie de type fini, décrite par un groupe algébrique réductif. Il
demande comment ces groupes dépendent de $\xi$ --- ils se déduisent les uns des
autres par torsion --- et note que la question est déjà intéressante pour $S$
le spectre d'un corps de nombres.

En caractéristique $p > 0$, il n'existe en général pas de foncteur fibre à
valeurs dans $\mathbb{Q}$ : c'est le contre-exemple de Serre, lié aux courbes
elliptiques d'invariant de Hasse nul.\footnote{Les endomorphismes d'une courbe
elliptique supersingulière forment une algèbre de quaternions sur $\mathbb{Q}$,
ramifiée en $p$ et en l'infini, qui ne peut opérer sur un $\mathbb{Q}$-espace
vectoriel de dimension $2$ --- ni d'ailleurs sur un $\mathbb{Q}_{p}$- ou un
$\mathbb{R}$-espace de dimension $2$.} On dispose en revanche d'un foncteur fibre
$\ell$-adique en tout point géométrique de caractéristique $\neq \ell$, et donc,
sur une sous-catégorie de type fini, d'un foncteur fibre sur une extension
finie $K$ de $\mathbb{Q}$, d'où un groupe réductif $G_{K}$.

\pagerange{6}{9}
\subsection*{Le groupe fondamental (pages 6 à 9)}

Un foncteur fibre $F$ sur $\mathcal{M}(S)$ et un point géométrique $\xi$ sont
\emph{associés} si l'on s'est donné un isomorphisme entre la restriction de $F$
aux motifs effectifs de poids $0$ et le foncteur fibre en $\xi$ --- c'est une
donnée, non une condition. Ainsi le foncteur de Betti sur $\operatorname{Spec}
\mathbb{C}$, ou le foncteur $\ell$-adique ; mais pas le foncteur de De Rham sur un
corps de caractéristique nulle.\footnote{Son exemple 4 --- sur un corps
algébriquement clos, tout $F$ serait associé à $\xi$, « du moins si $k$ est
dénombrable » --- est donné par la page avec une précision de marge (« de degré
de transcendance dénombrable ») et reste allusif ; on ne le développe pas.}

Les représentations continues de dimension finie de $\pi_{1}(S, \xi)$ sur
$\mathbb{Q}$ --- qui se factorisent par un quotient fini --- forment une
sous-catégorie pleine des motifs, celle qu'on appelle aujourd'hui les
\emph{motifs d'Artin}. D'où un épimorphisme $G = \varprojlim G_{\alpha} \to
\pi_{1}(S, \xi)$, qui se factorise par les composantes connexes :
\[
(*) \qquad \varprojlim_{\alpha} G_{\alpha}/G_{\alpha}^{0}
\;\twoheadrightarrow\; \pi_{1}(S, \xi) .
\]
C'est un isomorphisme si les motifs dont le groupe de Galois motivique est
fini sont exactement les motifs de niveau et de poids zéro, et c'est ce que
donnent les conjectures de Tate, ou de Hodge si $S$ est de caractéristique
générique nulle.\footnote{La page encadre $(*)$ comme un isomorphisme dès qu'elle
l'écrit, puis le justifie par les conjectures quatre lignes plus bas. Ce qui
est acquis sans elles est l'épimorphisme ; on l'écrit ainsi. Un passage
antérieur de la page 7, barré de hachures, disait la même chose et est repris
en tête de la page 8.} Enfin, si $S' \to S$ est un revêtement étale connexe et
si $F$ provient de $S'$, le groupe $G'$ s'identifie à l'image inverse dans $G$
du sous-groupe ouvert de $\pi_{1}$ correspondant ; d'où il suit que la
bijectivité de $(*)$ pour tout $S$ équivaut à la connexité des groupes $G$ sur
les corps algébriquement clos.

\pagerange{9}{13}
\subsection*{Le groupe fondamental profini, et le Frobenius (pages 9 à 13)}

Supposons $K \subset \mathbb{Q}_{\ell}$ et $F$ \emph{compatible} au foncteur de
Tate, c'est-à-dire muni d'un isomorphisme $F \otimes_{K} \mathbb{Q}_{\ell}
\simeq T^{\ell}_{\xi}$. L'action de $\pi_{1}$ sur $T^{\ell}_{\xi}$ donne
\[
\rho_{\ell} \colon \pi_{1}(S, \xi) \longrightarrow G(\mathbb{Q}_{\ell}) ,
\]
et, $F$ étant a fortiori associé à $\xi$, la flèche $G \to \pi_{1}$ donne
$\eta_{\ell} \colon G(\mathbb{Q}_{\ell}) \to \pi_{1}(S, \xi)$, avec
$\eta_{\ell}\rho_{\ell} = \mathrm{id}$. Pour une sous-catégorie de groupe
$G_{\alpha}$, le quotient $\pi_{\alpha}$ de $\pi_{1}$ défini par
$G_{\alpha}/G_{\alpha}^{0}$ s'insère dans un diagramme commutatif :

\begin{tikzcd}
 & G_{\alpha}/G_{\alpha}^{0} \arrow[dr, dashed] & \\
\pi_{1}(S, \xi) \arrow[r, "\rho_{\ell}"] \arrow[rr, bend right=40, "\text{can.}"'] & G_{\alpha}(\mathbb{Q}_{\ell}) \arrow[r] \arrow[u, dashed] & \pi_{\alpha}
\end{tikzcd}

La \emph{conjecture de Tate}, avec la semi-simplicité de l'action de Galois,
équivaut à dire que, pour $S$ localisé d'un schéma de type fini sur
$\mathbb{Z}$, $G_{\alpha}$ est l'enveloppe algébrique --- l'adhérence de
Zariski --- de l'image de $\rho_{\ell}$.\footnote{C'est aujourd'hui la forme
courante de la conjecture de Tate en termes de groupe de Galois motivique.
Tate l'a démontrée pour les endomorphismes des variétés abéliennes sur un corps
fini (1966), Faltings sur un corps de nombres (1983) ; elle est ouverte en
général. La marge ajoute que la conjecture sous la forme des classes de cycles
invariantes implique, avec les conjectures standard, la semi-simplicité.}

Sur un corps fini $k$, le Frobenius est un élément canonique de
$\operatorname{Aut}(\mathrm{id}_{\mathcal{M}(k)})$, donc du centre : il donne
$\mathrm{Frob}_{k,F} \in Z_{G}(\mathbb{Q})$, défini sur $\mathbb{Q}$, dont
l'image dans $G(\mathbb{Q}_{\ell})$ est celle du Frobenius de $\pi_{1}(k) \simeq
\hat{\mathbb{Z}}$ par $\rho_{\ell}$. La conjecture de Tate pour $k$ dit alors
que $G$ est l'enveloppe algébrique du sous-groupe engendré par
$\mathrm{Frob}_{k,F}$ : $G$ est commutatif, extension de $\hat{\mathbb{Z}}$ par
un tore --- un protore pour la catégorie entière. Pour $S$ quelconque, chaque
point fermé $s$ à corps résiduel fini donne un élément $f_{s}$ du schéma des
invariants $I_{G} = G /\!\!/ \operatorname{Ad} G$, qui pour $G$ connexe est $T/W$ ;
c'est la famille des Frobenius, compatible aux flèches $G \to \pi_{1}$ et
$\pi_{1} \to G(\mathbb{Q}_{\ell})$.

\pagerange{14}{16}
\subsection*{Structures supplémentaires sur le groupe de Galois motivique (pages 14 à 16)}

Le paragraphe qu'il numérote 12 --- il n'y a pas de 11 --- traduit en
homomorphismes de groupes les structures dont la catégorie est munie.

\emph{Graduations.} Une graduation de type $M$ du foncteur fibre, compatible au
produit tensoriel, est la même chose qu'un homomorphisme $D_{K}(M) \to G$ du
groupe diagonalisable de caractères $M$ ; elle provient d'une graduation des
objets eux-mêmes si et seulement si cet homomorphisme est central.\footnote{Le
même critère, dans les mêmes termes, ouvre les « Notes Saavedra » du dossier 10,
pages 101 et suivantes.} Ainsi la graduation par le poids donne
un cocaractère central
\[
y \colon \mathbb{G}_{m,K} \longrightarrow G ;
\]
si $F = T_{\mathrm{Hdg}}$, la bigraduation de Hodge donne deux cocaractères qui
commutent, $i_{1}, i_{2} \colon \mathbb{G}_{m} \to G$, avec $i_{1}(\lambda)
i_{2}(\lambda) = y(\lambda)$ --- ce qui contient déjà qu'ils commutent. Sur
$\mathbb{C}$ et pour $F = T_{B}$, on peut prendre $K = \mathbb{Q}$ ; la
décomposition de Hodge donne un isomorphisme de foncteurs fibres $T_{B} \otimes
\mathbb{C} \simeq T_{\mathrm{Hdg}}$, donc $G^{B}_{\mathbb{C}} \simeq
G^{\mathrm{Hdg}}$, et les deux cocaractères vivent dans $G^{B}_{\mathbb{C}}$ ;
ils sont conjugués l'un de l'autre :
\[
\overline{i_{1}(\lambda)} = i_{2}(\bar{\lambda}) , \qquad\text{c'est-à-dire}\qquad
i_{2} = \bar{i}_{1} .
\]\footnote{La page écrit $\overline{i_{1}(\lambda)} = i_{2}(\lambda)$. Pour la
conjugaison complexe de $G^{B}(\mathbb{C})$, on vérifie sur $V^{p,q}$ que
$\overline{i_{1}(\lambda)}$ agit par $\bar{\lambda}^{q}$, c'est-à-dire comme
$i_{2}(\bar{\lambda})$ ; l'énoncé juste porte sur les cocaractères, $i_{2}$
étant le conjugué de $i_{1}$ comme morphisme. C'est ce que la page veut dire
quand elle parle de graduations conjuguées.}

\emph{L'objet de Tate.} Si $\mathbb{Q}(-1)$ est dans la sous-catégorie, il
définit un caractère $\varepsilon \colon G \to \mathbb{G}_{m,K}$ et, puisque
$\mathbb{Q}(-1)$ est de poids $2$,
\[
\varepsilon(y(\lambda)) = \lambda^{2} .
\]

\emph{Filtrations.} Soit sur $F$ une filtration fonctorielle compatible au
produit tensoriel. Le gradué $F' = \operatorname{Gr}(F)$ est un autre foncteur
fibre, défini par un torseur $P$ sous $G$ ; inversement $F$ s'obtient de $F'$
par un torseur $P'$ sous $G' = \underline{\operatorname{Aut}}^{\otimes}(F')$.
\emph{Supposons} que, sur une extension $K'$ de $K$, il existe un isomorphisme
$F \simeq F'$ compatible aux filtrations et induisant l'identité sur les
gradués --- c'est le cas pour $T_{DR}$ muni de sa filtration de Hodge sur un
corps de caractéristique nulle. Alors $P$ contient un sous-torseur canonique $Q$
sous le groupe $H$ des automorphismes qui respectent la filtration et induisent
l'identité sur le gradué, et de même $Q' \subset P'$ sous $H'$ ; $H$ et $H'$
sont unipotents. Donc $F$ se déduit canoniquement du foncteur gradué $F'$ ---
graduation définie par $\alpha \colon \mathbb{G}_{m,K} \to G'$ --- et d'un torseur
$Q'$ sous le sous-groupe unipotent $H(\alpha)$ de $G'$ défini par
$\alpha$.\footnote{L'hypothèse d'existence d'un isomorphisme filtré est la sienne,
et elle est nécessaire : sans elle $Q'$ n'est qu'un pseudo-torseur. Le dossier 10
(pages 88 à 98) reprend ce dictionnaire en détail et bute sur le même point. La
construction de $H(\alpha)$ est laissée, en marge, « à expliciter » : c'est le
radical unipotent du parabolique défini par $\alpha$, $\mathfrak{h} =
\bigoplus_{n \geqslant 1} \mathfrak{g}_{n}$.}

\pagerange{16}{18}
\subsection*{Le groupe de Mumford--Tate, et la conjecture de Hodge (pages 16 à 18)}

Les structures de Hodge polarisables forment une $\otimes$-catégorie
semi-simple, dont toute sous-catégorie de type fini est définie par un groupe
réductif sur $\mathbb{Q}$. Soit $(V, V_{0}, (V^{p,q}))$ une structure de Hodge,
$\Gamma = \underline{\operatorname{Aut}}_{\mathbb{Q}}(V_{0})$, et $i_{1}, i_{2}
\colon \mathbb{G}_{m,\mathbb{C}} \to \Gamma_{\mathbb{C}}$ les cocaractères de la
bigraduation, $i(\lambda, \mu)$ agissant par $\lambda^{p}\mu^{q}$ sur $V^{p,q}$.
Le plus petit sous-groupe algébrique $G \subset \Gamma$ défini sur $\mathbb{Q}$
tel que $G_{\mathbb{C}}$ contienne les images de $i_{1}$ et $i_{2}$ --- c'est
aussi le sous-groupe engendré par leurs conjugués sous
$\operatorname{Aut}(\mathbb{C})$ --- est réductif : c'est le \emph{groupe de
Mumford--Tate}. Les éléments des constructions tensorielles sur $V_{0}$
invariants par $G$ sont exactement ceux de bidegré $(0,0)$ ; $G$ est donc le
groupe de Galois de la structure de Hodge, au sens tannakien.\footnote{La
réductivité demande la polarisabilité, que la phrase d'ouverture suppose. Le
dossier 13 (pages 3 à 11) construit le même groupe par un morphisme du tore
$\mathbb{S}$ dans $G_{\mathbb{R}}$ ; ici il est défini par ses deux cocaractères
complexes.}

Si $V$ provient d'un motif $M$, engendrant une $\otimes$-catégorie de motifs
$\mathcal{C}_{0}$, le foncteur de réalisation $\mathcal{C}_{0} \to \mathcal{C}$
vers les structures de Hodge donne une inclusion du groupe de Mumford--Tate
$G$ dans le groupe de Galois motivique $G_{0} \subset \Gamma$. La conjecture de
Hodge s'interprète alors comme l'\emph{égalité}
\[
G = G_{0} .
\]\footnote{L'inclusion $G \subset G_{0}$ est acquise ; l'égalité est la
conjecture de Hodge --- les classes de Hodge sur les puissances de la variété
sont algébriques --- et elle est ouverte. Cette reformulation est aujourd'hui
standard.}

Le paragraphe 14, « Structure de positivité », n'a que son titre, et la page est
blanche dessous. La positivité reviendra deux fois : pages 85 à 88, et pages 157
à 166.

\pagerange{20}{20}
\section*{II. Le plan « Motifs » (page 20)}

Sur un feuillet séparé, sous le titre encadré « Motifs », dix-neuf points
numérotés par lui : la catégorie $\mathcal{M}^{+}(X)$ ; les variances ; le cas
$X = \varprojlim X_{i}$ ; les foncteurs $T_{\ell}$ ; les $\mathbb{Q}_{\ell}(n)$ ; la
catégorie $\mathcal{M}(X)$ ; les $\underline{\operatorname{Hom}}$ et
$R\underline{\operatorname{Hom}}$ ; motifs constants, tordus ; filtrations de
$\mathcal{M}(X)$ et $\mathcal{M}^{+}(X)$ ; anneaux $\mathbf{M}^{+}(X)$ et
$\mathbf{M}(X)$ ; interprétation topologique des niveaux ; l'homomorphisme
$L(K) \to \mathbf{M}^{+}(K)$ et les invariants birationnels fondamentaux ;
caractérisation géométrique des filtrations (conjecture de Tate) ; invariants de
Galois et théorèmes de commutation ; cohomologie absolue ; relation avec les
variétés abéliennes sur un schéma de type fini ; formes positives ; motifs et
fonctions $L$ ; relation avec la cohomologie de Hodge.

C'est la table de ce qui suit. Les points 1 à 10 sont traités un à un aux pages
22 à 40, les points 11 à 16 sont ses paragraphes 11 à 16 des pages 41 à 79, les
points 17 et 18 ses paragraphes 17 et 18 des pages 85 à 94. Le point 19 n'est
écrit nulle part.

\pagerange{22}{40}
\section*{III. Les motifs sur une base : les points 1 à 10 du plan (pages 22 à 40)}

Une rédaction rapide et continue, non paginée, sous le titre « Motifs », en
paragraphes numérotés 1) à 10) qui suivent un à un le plan de la page 20. Tout y
est au mode du souhait --- « on veut que », « on désire que » : c'est un cahier
des charges pour une théorie des motifs \emph{mixtes} sur une base quelconque,
écrit à une époque où même les motifs purs sur un corps n'étaient pas
construits.\footnote{Le formalisme qui répond aujourd'hui en partie à ce cahier
des charges est celui des catégories \emph{triangulées} de motifs mixtes
(Voevodsky, vers 2000 ; Ayoub, 2007 ; Cisinski--Déglise, 2019), avec les six
opérations. Une catégorie \emph{abélienne} de motifs mixtes ayant les propriétés
demandées ici reste conjecturale.}

\pagerange{22}{24}
\subsection*{La catégorie, les opérations, les limites (pages 22 à 24)}

À tout préschéma noethérien $X$ est associée une catégorie abélienne
$\mathbb{Q}$-linéaire $\mathcal{M}^{+}(X)$ des motifs effectifs, munie d'un
produit tensoriel commutatif et unitaire, exact à droite, d'unité
$\mathbf{1}_{X} = \mathbb{Q}_{X}(0)$, et sa catégorie dérivée bornée
$D^{b}(\mathcal{M}^{+}(X))$ avec $\overset{L}{\otimes}$. Un morphisme
$f \colon X \to Y$ donne un foncteur exact $f^{*}$, compatible au produit
tensoriel ; si $f$ est de type fini, et propre ou $Y$ excellent, un foncteur
$Rf_{*}$, avec transitivité et formule de projection
\[
Rf_{*}(M \otimes Lf^{*}N) \simeq Rf_{*}(M) \otimes N .
\]
Il note en marge qu'il faudra aussi tout le formulaire de dualité, avec $f^{!}$
et $Rf_{!}$. Pour $X = \varprojlim X_{i}$, limite filtrante à morphismes de
transition affines, $\mathcal{M}^{+}(X) \simeq \varinjlim \mathcal{M}^{+}(X_{i})$,
et $Rf_{*}$ commute au passage à la limite : on se ramène aux préschémas de type
fini sur $\mathbb{Z}$.

\pagerange{24}{28}
\subsection*{La réalisation $\ell$-adique, l'objet de Tate, les motifs non effectifs (pages 24 à 28)}

Pour $\ell$ inversible sur $X$, un foncteur
\[
T_{\ell} \colon \mathcal{M}^{+}(X) \longrightarrow \mathcal{M}_{\ell}(X)
\]
vers les $\mathbb{Q}_{\ell}$-faisceaux constructibles, compatible au produit
tensoriel et aux unités, exact, fidèle mais non pleinement fidèle, compatible à
$f^{*}$, à $Rf_{*}$, à la formule de Künneth, et qui s'étend aux catégories
dérivées.\footnote{Un paragraphe sur un foncteur $T_{\infty}$ vers les Modules
cohérents, pour $X$ de type fini sur $\mathbb{Q}$, est encadré puis barré sur la
page 25. La marge de la même page observe que l'existence de $T_{\ell}$ exclut
la solution triviale $\mathcal{M}^{+}(X) = 0$.}

L'objet de Tate $\mathbb{Q}_{X}(-1)$, compatible au changement de base, se
réalise en $\mathbb{Q}_{\ell}(-1)$, dual du module de Tate de $\mathbb{G}_{m}$ ;
on peut le définir comme $R^{2}f_{*}(\mathbf{1})$ pour $f \colon
\mathbb{P}^{1}_{X} \to X$. Posant $\mathbb{Q}(-n) = \mathbb{Q}(-1)^{\otimes n}$,
on obtient, pour $f \colon X \to S$ lisse et projectif à fibres géométriques
connexes non vides de dimension relative $d$,
\[
R^{2d}f_{*}(\mathbf{1}_{X}) \simeq \mathbb{Q}_{S}(-d) ,
\]
et $R^{2d}f_{!}(\mathbf{1}_{X}) \simeq \mathbb{Q}_{S}(-d)$ si $f$ est seulement
quasi-projectif.\footnote{La page écrit $\mathbb{Z}_{S}(-d)$ ; la catégorie
étant $\mathbb{Q}$-linéaire, on lit $\mathbb{Q}_{S}(-d)$. Il écrit indifféremment
$\mathbb{Z}(-1)$ et $\mathbb{Q}(-1)$ dans la suite.} Pour une immersion fermée
$i \colon Y \to X$ de codimension $d$ de schémas lisses sur $S$, il demande
l'isomorphisme de pureté
\[
Ri^{!}(\mathbf{1}_{X}) \simeq \mathbb{Q}_{Y}(-d)[-2d] ,
\]
compatible aux isomorphismes $\ell$-adiques connus.\footnote{La page écrit
$Ri^{!}(\mathbf{1}_{X}) \simeq \mathbb{Q}_{Y}(-d)$, sans le décalage ; la
cohomologie à supports de $Y$ est concentrée en degré $2d$, et c'est
$R^{2d}i^{!}$ qui vaut $\mathbb{Q}_{Y}(-d)$.}

Le foncteur $M \mapsto M(-1) = M \otimes \mathbb{Q}(-1)$ est pleinement fidèle
sur $\mathcal{M}^{+}(X)$ sans en être une équivalence ;\footnote{C'est ce qu'on
appelle aujourd'hui la \emph{simplification} (\emph{cancellation}), démontrée par
Voevodsky pour ses motifs triangulés sur un corps parfait (2010).} on obtient
$\mathcal{M}(X)$ comme pseudo-limite inductive de la suite $\mathcal{M}^{+}
\xrightarrow{\otimes\mathbb{Q}(-1)} \mathcal{M}^{+} \to \cdots$, où
$\mathbb{Q}(-1)$ devient inversible d'inverse $\mathbb{Q}(1)$. Tout objet de
$\mathcal{M}(X)$ s'écrit $M_{n}(n)$ avec $M_{n}$ effectif pour $n$ assez grand,
et $M_{n}(n) \simeq M_{n+1}(n+1)$ équivaut à $M_{n+1} \simeq M_{n}(-1)$ ;
$\mathcal{M}^{+}(X)$ est une sous-catégorie pleine épaisse de $\mathcal{M}(X)$, et
$T_{\ell}$ s'étend de façon unique. Dans $\mathcal{M}(X)$ il y a des
$\underline{\operatorname{Hom}}$ et des $R\underline{\operatorname{Hom}}$, adjoints
au produit tensoriel et compatibles à $T_{\ell}$.

\pagerange{29}{30}
\subsection*{Indépendance de $\ell$ (pages 29 et 30)}

Pour $\ell$, $\ell'$ premiers aux caractéristiques résiduelles, on veut que
$T_{\ell}(M)$ soit un faisceau lisse --- « constant tordu » --- si et seulement
si $T_{\ell'}(M)$ l'est, avec le même rang. Il suffit de le voir sur un corps, et
plus généralement, pour $u \in \operatorname{End}(M)$,
\[
\operatorname{Tr} T_{\ell}(u) = \operatorname{Tr} T_{\ell'}(u) \in \mathbb{Q} ,
\qquad
\det(1 - t\,T_{\ell}(u)) = \det(1 - t\,T_{\ell'}(u)) \in \mathbb{Q}[t] ,
\]
la seconde égalité se déduisant de la première appliquée aux $\wedge^{n}u$.
L'argument est celui qu'on donne aujourd'hui : $u$ est un élément de
$\operatorname{Hom}(\mathbf{1}_{X}, M^{\vee} \otimes M)$, la contraction
$M^{\vee} \otimes M \to \mathbf{1}_{X}$ en fait un $c(u) \in
\operatorname{End}(\mathbf{1}_{X})$, et $\operatorname{Tr} T_{\ell}(u) =
T_{\ell}(c(u))$ ; il suffit donc que $\operatorname{End}(\mathbf{1}_{X}) =
\mathbb{Q}$ pour $X$ connexe.\footnote{L'argument demande que $M$ ait un dual
$M^{\vee}$ avec contraction et coévaluation, c'est-à-dire que la catégorie soit
rigide ; les motifs effectifs ne le sont pas, mais $\mathcal{M}(X)$ doit l'être,
et c'est là qu'on se place. La page le note seulement en marge : « à inclure
dans le formalisme tensoriel ».}

\pagerange{30}{33}
\subsection*{La filtration par le poids (pages 30 à 33)}

Les sous-catégories $\mathcal{M}^{+}_{i}(X)$ des motifs effectifs à poids
$\leqslant i$ forment une filtration exhaustive,\footnote{Il écrit d'abord
les indices en haut, $\mathcal{M}^{+i}(X)$, puis note quatre fois en marge
« indices en bas » et passe aux indices inférieurs à la page 33 ; on les met en
bas partout.} dont l'appartenance se vérifie fibre à fibre, en un point fermé
suffisant quand $X$ est de type fini --- et, sur un corps fini, par les valeurs
absolues du Frobenius. Elle est locale, compatible à $f^{*}$, et
\[
R^{j}f_{!} \colon \mathcal{M}^{+}_{i}(X) \longrightarrow \mathcal{M}^{+}_{i+j}(Y) ,
\]
ce qui est, transcrit, le théorème principal de Weil II.\footnote{Pour les
$\mathbb{Q}_{\ell}$-faisceaux mixtes, que $R^{j}f_{!}$ augmente les poids d'au
plus $j$ est le théorème de Deligne (« La conjecture de Weil II », 1980). La
page le pose comme exigence.} Sur un corps algébriquement clos,
$\mathcal{M}^{+}_{i}(k)$ est engendrée, par extensions, par des sous-quotients de
motifs $R^{j}f_{!}(\mathbf{1}_{X})$.\footnote{La page dit les
$R^{i}f_{!}(\mathbb{Z}\cdot\mathbf{1}_{X})$ « pour $X$ si on veut projectif lisse
de dim $\leqslant j$ ». Pour $X$ projectif et lisse, $R^{i}f_{!}$ est pur de poids
$i$ ; les poids $< i$ demandent d'autres degrés ou des $X$ non propres. Les
indices de la phrase ne se lisent pas sûrement, et on n'en garde que le
principe.} La filtration est multiplicative, $\mathcal{M}^{+}_{i} \otimes
\mathcal{M}^{+}_{k} \subset \mathcal{M}^{+}_{i+k}$, et $\mathbb{Q}(-1)$ est de
poids $2$ ; mieux, $M \in \mathcal{M}^{+}_{i}$ si et seulement si $M(-j) \in
\mathcal{M}^{+}_{i+2j}$. Elle se prolonge donc à $\mathcal{M}(X)$ en une filtration
$(\mathcal{M}_{i}(X))_{i \in \mathbb{Z}}$, avec $\mathcal{M}^{+}_{i} =
\mathcal{M}_{i} \cap \mathcal{M}^{+}$ (nul pour $i < 0$), et le produit par
$\mathbb{Q}(-1)$ est une équivalence $\mathcal{M}_{i} \simeq \mathcal{M}_{i+2}$ ;
en revanche $\mathbb{Q}(-j) \otimes \mathcal{M}^{+}_{i} \subset
\mathcal{M}^{+}_{i+2j}$ est stricte en général. Chaque cran est épais dans le
suivant, et les quotients
\[
\mathcal{G}^{+}_{i}(X) = \mathcal{M}^{+}_{i}(X) / \mathcal{M}^{+}_{i-1}(X) ,
\qquad
\mathcal{G}_{i}(X) = \mathcal{M}_{i}(X) / \mathcal{M}_{i-1}(X)
\]
sont « fort intéressants ». Le produit par $\mathbb{Q}(-j)$ donne
$\mathcal{G}_{i} \simeq \mathcal{G}_{i+2j}$, donc chaque $\mathcal{G}_{i}$ est
équivalent à $\mathcal{G}_{0}$ ou à $\mathcal{G}_{1}$ selon la parité de $i$, et
$\mathcal{G}^{+}_{i}$ s'identifie à une sous-catégorie pleine épaisse de
$\mathcal{G}_{i}$.\footnote{La page écrit $\mathcal{G}_{i}(X) \simeq
\mathcal{G}_{0}(X)$ pour tout $i$. La torsion par $\mathbb{Q}(-j)$ déplace les
poids de $2j$ : elle ne relie que des indices de même parité.}

\pagerange{34}{35}
\subsection*{La filtration par le niveau (pages 34 et 35)}

La filtration par le « type dimensionnel » $D_{i}$ est la plus petite filtration
épaisse qui majore celle par le poids et qui soit stable par $\otimes\,
\mathbb{Q}(-1)$ : $D_{i}(\mathcal{M}^{+}(X))$ est formée des objets qui, localement
pour la topologie étale, sont extensions d'objets $M(-j)$ avec $M \in
\mathcal{M}^{+}_{i}(X)$ et $j \geqslant 0$. Sur les motifs semi-simples sur un
corps, c'est exactement le niveau de la page 2 ; la filtration $D_{i}$ l'étend aux
motifs qui ne sont plus semi-simples. Elle est stable par image inverse, et
\[
R^{j}f_{!}\bigl(D_{i}\mathcal{M}^{+}(X)\bigr) \subset D_{i+j}\mathcal{M}^{+}(Y) ,
\]
l'analogue pour $R^{j}f_{*}$ restant marqué d'un point d'interrogation, avec
cette seule certitude que $R^{j}f_{*}$ envoie les poids $\leqslant i$ dans les
poids $\leqslant i+2j$. Elle est stable par torsion --- $M \in D_{i}$ entraîne
$M(-j) \in D_{i}$, « pas $D_{i+j}$ !! » ---, se prolonge à $\mathcal{M}(X)$, est
multiplicative, $D_{i} \otimes D_{j} \subset D_{i+j}$, avec $\mathbb{Q}(-j) \in D_{0}$
pour tout $j$, et
\[
\underline{\underline{\operatorname{Ext}}}^{i}(D_{j}, D_{k}) \subset D_{j+k} .
\]\footnote{La marge droite de la page 34, écrite de bas en haut et en partie
biffée, ne se lit pas ; une autre note dit que tout doit se passer « en termes
de la dimension de $X$ ».}

\pagerange{35}{40}
\subsection*{Motifs constants tordus, et leurs anneaux (pages 35 à 40)}

Soit $\mathcal{M}^{+}(X)_{\mathrm{sp}}$ la sous-catégorie des motifs effectifs
\emph{constants tordus} --- ceux dont la réalisation $\ell$-adique est un faisceau
lisse, condition locale qui se vérifie sur $T_{\ell}(M)$ par le point 8 --- et de
même $\mathcal{M}(X)_{\mathrm{sp}}$ ; $M$ est constant tordu si et seulement si
$M(j)$ l'est. Ce sont des sous-catégories abéliennes stables par extensions et
par produit tensoriel, et, pour la seconde, par $\underline{\operatorname{Hom}}$ ;
tout objet y est de longueur finie.\footnote{La page ajoute, en une ligne en
partie illisible, que les $\operatorname{Ext}^{i}$, $i > 0$, y seraient nuls
(lecture incertaine). Ce ne peut être vrai tel quel : sur $X = \operatorname{Spec}
k$, une unité $a \in k^{*}$ non de torsion définit une extension de Kummer non
triviale de $\mathbb{Q}(0)$ par $\mathbb{Q}(1)$, dont la réalisation est lisse.
Rien de ce qui suit ne s'en sert : les anneaux ne dépendent que de la longueur
finie.} Leur anneau de Grothendieck, muni de ses opérations $\lambda$, est noté
$\mathbf{M}^{+}(X)$, resp.\ $\mathbf{M}(X)$.

Soit $\xi = \operatorname{cl}\mathbb{Q}_{X}(-1)$. Alors
\[
\mathbf{M}(X) = \mathbf{M}^{+}(X)\bigl[\xi^{-1}\bigr] ,
\]
et $\xi$ n'est pas diviseur de zéro : $\mathbf{M}^{+}(X)$ est le
$\mathbb{Z}$-module libre de base les classes d'objets simples, et la
multiplication par $\xi$ y est induite par une injection de l'ensemble
$\Sigma^{+}(X)$ des classes de simples dans lui-même, puisque $M$ est simple si
et seulement si $M(-1)$ l'est. D'où $\mathbf{M}^{+}(X) =
\mathbb{Z}^{(\Sigma^{+}(X))} \subset \mathbf{M}(X) = \mathbb{Z}^{(\Sigma(X))}$,
où $\Sigma(X)$ est la limite inductive de $\Sigma^{+}(X)$ par $e \mapsto \xi e$.

Les objets simples étant purs,\footnote{C'est une attente --- un motif simple
est pur --- que la page emploie sans la dire.} la partition $\Sigma^{+}(X) =
\bigsqcup_{i} \Sigma^{+}_{i}(X)$ suivant le poids fait de $\mathbf{M}^{+}(X)$ un
anneau gradué, $\Sigma^{+}_{i}\cdot\Sigma^{+}_{j} \subset
\mathbb{Z}^{(\Sigma^{+}_{i+j})}$, où $\xi$ est de degré $2$, et de $\mathbf{M}(X)$
un anneau gradué en tous degrés où $\xi$ est inversible. La partie de degré $0$
est $\mathbf{M}_{0}(X) = \mathbb{Z}^{(\Sigma_{0}(X))}$, avec $\Sigma_{0}(X) =
\varinjlim_{j} \xi^{-j}\Sigma^{+}_{2j}(X)$, et
\[
\mathbf{M}^{\mathrm{pair}}(X) = \mathbf{M}_{0}(X)[\xi, \xi^{-1}] ,
\qquad
\mathbf{M}^{\mathrm{impair}}(X) = \mathbf{M}_{1}(X) \otimes_{\mathbf{M}_{0}(X)}
\mathbf{M}_{0}(X)[\xi, \xi^{-1}] ,
\]
de sorte que la structure d'anneau gradué de $\mathbf{M}(X)$ se lit sur l'anneau
$\mathbf{M}_{0}(X)$, le module $\mathbf{M}_{1}(X)$ et la multiplication
$\mathbf{M}_{1} \otimes \mathbf{M}_{1} \to \mathbf{M}_{2} \simeq
\mathbf{M}_{0}\,\xi$.\footnote{La page 38 dit $\xi$ « de degré $1$ » et écrit
$\mathbf{M}_{0}(X)[\xi]$, et $\Sigma_{0}$ comme limite des $\Sigma^{+}_{i}$. Mais
$\mathbb{Q}(-1)$ est de poids $2$ (pages 31 et 32), et la page 39 elle-même
décompose $\Sigma^{+}_{i}$ en $\sigma_{i} \cup \xi\sigma_{i-2} \cup \cdots$, où
$\xi$ déplace le degré de $2$. On a rétabli le degré $2$, les indices pairs dans
$\Sigma_{0}$, et l'inversion de $\xi$ dans $\mathbf{M}(X)$. L'énoncé sur
$\mathbf{M}_{1}$ est surchargé de petits ajouts illisibles sur la page.}

Les \emph{éléments primitifs} de poids $i$ sont ceux qui ne sont pas des
multiples de $\xi$ :\footnote{La page écrit la réunion pour
$\alpha \geqslant 0$, ce qui retirerait tout $\Sigma^{+}_{i}$ ; il faut
$\alpha \geqslant 1$.}
\[
\sigma_{i}(X) \;=\; \Sigma^{+}_{i}(X) \smallsetminus \bigcup_{\alpha \geqslant 1}
\xi^{\alpha}\,\Sigma^{+}_{i-2\alpha}(X) ,
\qquad
\Sigma^{+}_{i}(X) \;=\; \sigma_{i}(X) \sqcup \xi\,\sigma_{i-2}(X) \sqcup
\xi^{2}\sigma_{i-4}(X) \sqcup \cdots
\]
Un simple $\xi^{a}s$ avec $s \in \sigma_{k}$ est de poids $k + 2a$ et de niveau
$k$. La filtration par le niveau se lit donc ainsi :
\[
D_{i}\bigl(\Sigma^{+}(X)\bigr) \cap \Sigma^{+}_{j}(X) \;=\;
\bigsqcup_{\substack{k \leqslant \inf(i,j) \\ k \equiv j \ (2)}}
\xi^{(j-k)/2}\,\sigma_{k}(X) ,
\qquad
D_{i} \smallsetminus D_{i-1} \;=\; \bigsqcup_{a \geqslant 0} \xi^{a}\,\sigma_{i}(X) ,
\]
et $\mathbf{M}^{+}(X)$, $\mathbf{M}(X)$ avec leurs deux filtrations et la
multiplication par $\xi$ sont entièrement déterminés, comme groupes, par la famille
$(\sigma_{i}(X))_{i \geqslant 0}$.\footnote{La page 39 écrit, pour $j > i$,
$\xi^{j-i}\sigma_{i} \cup \xi^{j-i+1}\sigma_{i-1} \cup \cdots \cup \xi^{j}\sigma_{0}$,
formule juste si $\xi$ est de degré $1$, et donne pour le quotient
$D_{i}/D_{i-1}$ l'ensemble $\sigma_{i} + \xi\sigma_{i-2} + \xi^{2}\sigma_{i-4} +
\cdots$, qui est $\Sigma^{+}_{i}$, l'ensemble des simples de \emph{poids} $i$.
Avec $\xi$ de degré $2$ et $D_{i}$ la filtration par le niveau ($D_{i}$ formé des
$\xi^{j}s$, $s$ de poids $\leqslant i$, page 38), les deux formules sont celles
qu'on donne ici.} Pour la structure multiplicative il faut de plus les produits
$\sigma_{i} \times \sigma_{j} \to \mathbb{Z}^{(\cdots)}$. Il promet d'expliciter
plus bas $\sigma_{0}(X)$ et $\sigma_{1}(X)$ ; il ne le fait pas.

Restent les opérations. Pour la dualité, $\mathbb{Q}(-i) \otimes
\bigl[\mathcal{M}^{+}_{i}\bigr]^{\vee} \subset \mathcal{M}^{+}_{i}$, et $x \mapsto
\check{x}\,\xi^{i}$ est une involution de $\mathbf{M}_{i}$ ; la théorie devrait
montrer qu'elle est l'identité, c'est-à-dire que $\check{x} = \xi^{-i}x$ sur la
partie de poids $i$.\footnote{C'est l'existence, pour tout motif pur $M$ de poids
$i$, d'un isomorphisme $M^{\vee} \simeq M(i)$, c'est-à-dire d'une forme non
dégénérée $M \otimes M \to \mathbb{Q}(-i)$ : une polarisation. Sur les $h^{i}(X)$
elle vient de la dualité de Poincaré et du théorème de Lefschetz fort ; c'est
le sujet des pages 85 à 88 et 157 à 166.} Les opérations $\lambda^{i}$ devraient
se déterminer en termes de groupes finis d'opérateurs, tels les groupes
symétriques. La page 40 s'arrête là, sa moitié basse blanche.

\pagerange{41}{59}
\section*{IV. Coniveau, invariants birationnels, caractérisations galoisiennes (pages 41 à 59)}

La même rédaction continue, avec ses paragraphes 11 à 15, dont les titres
soulignés reprennent les points 11 à 15 du plan.

\pagerange{41}{42}
\subsection*{11) Le coniveau (pages 41, 42 et 57)}

Soit $X$ lisse de dimension $n$ sur un corps algébriquement clos. La filtration
de $X$ par la codimension donne la \emph{suite spectrale du coniveau}
\[
E_{1}^{pq} = \bigoplus_{x \in X^{(p)}} H^{q-p}(x, \mathbb{Q}_{\ell})(-p)
\;\Longrightarrow\; H^{p+q}(X, \mathbb{Q}_{\ell}) ,
\]
où $X^{(p)}$ est l'ensemble des points de codimension $p$ et $H^{*}(x) =
\varinjlim_{U} H^{*}(U)$ sur les ouverts denses de $\overline{\{x\}}$ ; les termes
sont nuls hors du secteur $0 \leqslant p \leqslant q$, et la marge souligne qu'on
utilise le terme $E_{1}$, non $E_{2}$.\footnote{Cette suite spectrale est celle
qu'étudieront Bloch et Ogus (1974). La page écrit $\mathbb{Z}_{\ell}$ dans l'une de
ses deux lignes de définition, $\mathbb{Q}_{\ell}$ partout ailleurs.} Elle
provient d'une suite spectrale de motifs, au moins à partir de $E_{2}$ --- sinon
il faudrait des ind-motifs, ou des filtrations finies convenables par des fermés.
Si $X$ est projective et lisse --- « hypothèse essentielle ! » ---, le cran
$\geqslant p$ de la filtration aboutissement de $H^{i}$ est de niveau
$\leqslant i-2p$ : c'est l'image, par des morphismes de Gysin, de la cohomologie
de degré $i-2p$ de variétés de dimension $n-p$, tordue par $\mathbb{Q}(-p)$. Le
feuillet 57, barré et hors de sa place, reprend ce point pour $i = n$ : les
$\operatorname{gr}^{p}H^{n}$ sont de niveau pur $n-2p$ et de poids
$n$.\footnote{Le feuillet 57 est barré de quatre traits ; son verso, écrit
dans l'autre sens, porte des suites de localisation $H^{i}_{Y}(X) \to H^{i}(X)
\to H^{i}(U) \to$ et $H^{i}_{c}$ avec des bornes en $\xi$ dont les indices
$m$ et $n$ se distinguent mal ; on n'en tire rien de plus.}

Un encadré de la page 42 ajoute l'observation qui ouvre le paragraphe suivant :
si $U$ est un ouvert dense de $X$, $H^{i}(X) \to H^{i}(U)$ est un isomorphisme
modulo $D_{i-1}$ ; les deux ont donc la même composante dans $D_{i}/D_{i-1}$, et
c'est un \emph{invariant birationnel}.\footnote{Le haut de la page 42, qui
parlait de remplacer le coefficient constant par un motif constant tordu de
niveau $\delta$ et rapprochait l'énoncé de la conjecture de Tate, est barré.}

\pagerange{42}{51}
\subsection*{12) L'homomorphisme $L(K) \to \mathbf{M}^{+}(K)$ et les invariants birationnels (pages 42 à 51)}

Soit $K$ un corps et $L(K)$ le groupe engendré par les classes des $K$-schémas
de type fini, soumis aux relations de découpage $[X] = [Z] + [X \smallsetminus
Z]$ ; c'est un anneau pour le produit, l'\emph{anneau de Grothendieck des
variétés}. L'homomorphisme d'anneaux
\[
\sigma \colon L(K) \longrightarrow \mathbf{M}^{+}(K) ,
\qquad
[X] \longmapsto \sum_{i} (-1)^{i}\operatorname{cl} H^{i}_{c}(X, \mathbb{Q}(0)) ,
\]
où $H^{i}_{c}(X) = R^{i}f_{!}(\mathbb{Q}(0))$ pour $f \colon X \to
\operatorname{Spec} K$, est ce qu'on appelle aujourd'hui la
\emph{caractéristique d'Euler motivique}. Ses composantes de poids sont des motifs
virtuels, et leurs rangs les « nombres de Betti virtuels » « demandés par
Serre ».\footnote{L'existence inconditionnelle de $\sigma$, à valeurs dans
l'anneau de Grothendieck des motifs de Chow, est due à Gillet--Soulé et à
Guillén--Navarro Aznar (1996) en caractéristique nulle. Les nombres de Betti
virtuels, eux, se définissent sans motifs dès qu'on a les poids : par la théorie
de Hodge mixte sur $\mathbb{C}$ (Deligne, 1971-1974) ou par Weil II sur un corps
fini.} La droite affine a pour image $\xi = \operatorname{cl}\mathbb{Q}(-1)$ :
c'est le motif de Lefschetz, qui élève la dimension d'une unité.

Filtrons $L(K)$ par la dimension. Une variété de dimension $\leqslant i$ a ses
$H^{i+k}_{c}$ dans $\xi^{k}\mathcal{M}^{+}_{(i-k)}$, d'où
\[
L_{(i)}(K) \longrightarrow D'_{i} = \mathbf{M}^{+}_{(i)} + \xi\,\mathbf{M}^{+}_{(i-1)}
+ \xi^{2}\mathbf{M}^{+}_{(i-2)} + \cdots
\]
et un homomorphisme naturel
\[
\operatorname{gr}_{n} L(K) \longrightarrow
\mathbb{Z}^{(\sigma_{n} \sqcup \xi\sigma_{n-1} \sqcup \xi^{2}\sigma_{n-2} \sqcup \cdots)} ,
\]
où $\xi^{a}\sigma_{n-a}$ est de poids $n+a$ et où les $\sigma_{k} = \sigma_{k}(K)$
sont les ensembles d'éléments primitifs des pages 39 et 40.\footnote{Les pages 43
à 57 écrivent $\zeta$ pour ce qu'on note ici $\xi$ ; c'est le même élément, celui
de la page 36. La page 43 donne une première version de ce paragraphe, barrée.}
À $X$ de dimension $\leqslant n$ est ainsi attaché un invariant
$\gamma_{n}(X) + \xi\gamma_{n-1}(X) + \xi^{2}\gamma_{n-2}(X) + \cdots$, avec
$\gamma_{j}(X) \in \mathbb{Z}^{(\sigma_{j}(K))}$. On l'obtient à partir des
composantes de poids pur $\beta_{k}(X)$ de Serre pour $k \geqslant n$ : un tel
$\beta_{k}$ s'écrit $\xi^{k-n}\beta'_{2n-k}$ avec $\beta'_{2n-k}$ de poids $2n-k$,
et $\gamma_{2n-k}$ est, au signe $(-1)^{k}$ près, la composante primitive de
$\beta'_{2n-k}$. Par construction ce sont des \emph{invariants $n$-birationnels}.
Si $X$ est projective et lisse, $\beta'_{j} = \operatorname{cl}H^{j}(X)$ pour
$0 \leqslant j \leqslant n$, par le théorème de Lefschetz fort, et $\gamma_{j}(X)$
est la composante primitive --- au sens du niveau --- de $H^{j}(X)$.\footnote{La
page dit « le dual (lecture incertaine) des $H^{j}$ ». Par Lefschetz,
$H^{2n-j}(X) \simeq H^{j}(X)(j-n)$ ; c'est donc $H^{j}$ lui-même qu'on trouve, et
comme $H^{j}$ est autodual à torsion près, les deux lectures donnent la même
classe.}

La même chose se dit mieux avec des objets. Pour $X$ lisse de dimension $n$ et
$i \leqslant n$, et pour tout ouvert dense $U$, la flèche $H^{i}(X) \to H^{i}(U)$
est injective modulo $D_{i-2}$ et bijective modulo $D_{i-1}$ : on obtient un
invariant birationnel dans $D_{i}/D_{i-1}$, qui provient de
$\mathcal{M}^{+}_{(i)}$, d'où un élément de $\mathbb{Z}^{(\sigma_{i}(K))}$ :
c'est $\gamma_{i}(X)$. Mieux encore, $K$ étant parfait, soit
\[
\Gamma^{i}(X) \subset \varinjlim_{U} H^{i}(U, \mathbb{Q}(0))
\]
le plus grand sous-motif \emph{bipur} de type $i$ ; c'est un motif semi-simple,
fonctoriel en $X$ --- ou en son corps de fonctions ---, nul pour les morphismes
non dominants et \emph{injectif} pour les morphismes dominants, parce que
$H^{i}(X) \to H^{i}(X')$ est injectif pour $X' \to X$ propre et surjectif entre
variétés lisses.\footnote{La page écrit ici $\mathcal{M}^{+}(X)$ et $\sigma_{i}(X)$
là où les pages voisines ont $\mathcal{M}^{+}(K)$, $\sigma_{i}(K)$ ; on a corrigé.
Le « TMB » qui introduit le cas non dominant est lu tel quel.}

Suivent des questions de « fidélité » de ce foncteur. Si deux corps de fonctions
$R$, $R'$ ont des systèmes $(\Gamma_{i})$ isomorphes, chacun
s'immerge-t-il dans une extension transcendante pure de l'autre ? Deux exemples
précisent la question :
\begin{enumerate}
  \item si deux courbes $C$, $C'$ lisses, complètes, connexes sur $K$
  algébriquement clos ont des jacobiennes reliées par un homomorphisme à noyau
  fini, existe-t-il une application rationnelle dominante $C' \times
  \mathbb{A}^{N} \dashrightarrow C$ ?
  \item si les $\Gamma_{i}$ sont nuls pour $i > 0$ --- par exemple si la
  cohomologie d'une variété projective lisse est entièrement algébrique ---, la
  variété est-elle unirationnelle ? Et pour une surface --- « Cast.\ ok ! ».
\end{enumerate}
Et une remarque : si l'on néglige les variétés réglées, le seul invariant
cohomologique qui reste est $\Gamma_{n}$ ; a-t-on un isomorphisme entre
$\operatorname{gr}_{n}L(K)$, modulo le sous-groupe « réglé » $\xi\operatorname{gr}_{n-1}
L(K)$, et $\mathbb{Z}^{(\sigma_{n}(K))}$ ? En particulier, si $\Gamma_{n}(X) = 0$,
$X$ est-elle réglée ? Pour une surface dont le $H^{2}$ est algébrique, la marge
répond : « Non en car.\ $p > 0$ », avec $E \times E$ pour $E$ supersingulière.

Les trois questions ont une réponse négative, et en caractéristique nulle
aussi.\footnote{Ceci est de nous. (1) Si $C$ est de genre $\geqslant 1$, une
application rationnelle $C' \times \mathbb{A}^{N} \dashrightarrow C$ est
constante sur les droites, donc se factorise par $C'$ ; si $C$ et $C'$ sont de
même genre $g \geqslant 2$, une application dominante $C' \to C$ est de degré $1$
par Riemann--Hurwitz, donc un isomorphisme. Or il existe des courbes de genre $2$
non isomorphes à jacobiennes isogènes (par exemple isogènes à un même produit
$E \times E$). (2) et la remarque : une surface d'Enriques complexe a $H^{1} = 0$
et $H^{2}$ entièrement algébrique, donc $\Gamma_{1} = \Gamma_{2} = 0$ ; elle n'est
pas réglée, étant de dimension de Kodaira nulle, ni unirationnelle, son groupe
fondamental étant d'ordre $2$ (une variété unirationnelle complexe est
simplement connexe, Serre 1959). Le critère de Castelnuovo demande $q = P_{2} =
0$, et $H^{2}$ algébrique ne donne que $q = p_{g} = 0$ : le « Cast.\ ok ! » de la
page va trop vite.} C'est ce que la théorie des invariants birationnels
stables a précisé depuis.\footnote{Larsen et Lunts (2003) ont montré, en
caractéristique nulle, que le quotient de $L(K)$ par l'idéal engendré par la
classe de la droite affine est l'anneau libre sur les classes de rationalité
stable des variétés projectives et lisses. C'est la forme actuelle de la question
de la page 49.}

La page 50 dresse, pour $X$ lisse purement de dimension $n$, les bornes que ces
calculs utilisent : $H^{i}_{c}(X)$ est de poids $\leqslant i$ et de niveau
$\leqslant n$, et mieux $H^{n+k}_{c}(X) \in \xi^{k}\mathcal{M}^{+}_{(n-k)}$ ;
$H^{i}(X)$ est dans $D_{i}$ et dans $D_{n}$, et mieux dans $\mathcal{M}^{+}_{(i)} +
\xi\mathcal{M}^{+}_{(i-1)} + \cdots$ si $i \leqslant n$, dans
$\xi^{i-n}[\mathcal{M}^{+}_{(2n-i)} + \xi\mathcal{M}^{+}_{(2n-i-1)} + \cdots]$ si
$i \geqslant n$ ; enfin les poids de $H^{i}(X, M)$, pour $M$ pur de poids $\rho$,
sont compris entre $\rho + i$ et $\rho + \inf(2i, 2n)$.\footnote{La page écrit
$\sum \operatorname{Gr}_{i}$ où l'indice de sommation est $\alpha$ ; on lit
$\operatorname{Gr}_{\alpha}$. Ce sont les bornes que Deligne établira en théorie de
Hodge mixte (1971) et, $\ell$-adiquement, par Weil II.} Le feuillet 51, barré,
demande si $L_{i}(K)/L_{i-1}(K) \to \mathbb{Z}^{(\sigma_{i}(K))}$ est injectif ou
surjectif, au moins modulo torsion, et commence par observer que la classe de
$X$ dans $L(K)$ détermine l'Albanese à isogénie près d'une surface.

\pagerange{52}{54}
\subsection*{13) Caractérisation galoisienne des filtrations (pages 52 à 54)}

Les classes de faisceaux $\ell$-adiques constants tordus qui proviennent de
motifs forment un anneau $\mathbf{M}^{+}_{\ell}(X)$, quotient de
$\mathbf{M}^{+}(X)$. Pour $X$ de type fini sur $\mathbb{Z}$, la surjection
$\mathbf{M}^{+}(X) \to \mathbf{M}^{+}_{\ell}(X)$ est bijective si deux motifs
constants tordus simples non isomorphes ont des réalisations sans facteur
simple commun --- ce que donnent les conjectures de Tate. Les filtrations se
lisent alors sur le Frobenius. Pour $E$ constant tordu sur $X$ de type fini
sur $\mathbb{Z}$ :
\begin{itemize}
  \item $\operatorname{cl}(E) \in \mathbf{M}_{i}(X)$ si et seulement si, en tout
  point fermé $x$, les valeurs propres de Frobenius sur $T_{\ell}(E)_{x}$ sont de
  valeur absolue $N(x)^{i/2}$ --- ce sont, dit la marge, les conjectures de
  Weil ;
  \item pour $E$ simple, $\operatorname{cl}(E) \in \mathbf{M}^{+}(X)$ si et
  seulement si, en tout point fermé, le polynôme caractéristique du Frobenius
  géométrique $\mathrm{Frob}_{x}^{-1}$ est à coefficients entiers, c'est-à-dire
  si ses valeurs propres sont des entiers algébriques --- « Conjectures de
  Tate ! ».\footnote{Le sens direct est un théorème : les valeurs propres du
  Frobenius géométrique sur la cohomologie à supports propres sont des entiers
  algébriques (Deligne, SGA 7, exposé XXI). La réciproque, qu'un motif simple à
  valeurs propres entières soit effectif, est conjecturale, et c'est elle que la
  marge rapporte à Tate.}
\end{itemize}
Pour le niveau, prenons $E$ homogène de degré $\alpha$. Pour $i = 0$ : $E \in
D_{0}$ si et seulement si $\alpha = 2\beta$ est pair et $E(\beta)$ effectif,
c'est-à-dire si les $\rho(x)/N(x)^{\beta}$ sont des entiers algébriques, $\rho(x)$
parcourant les valeurs propres de $\mathrm{Frob}_{x}^{-1}$ ; ce sont alors des
racines de l'unité, et $\pi_{1}(X)$ opère sur $T_{\ell}(E)(\beta)$ à travers un
quotient fini --- énoncé qui, note-t-il, ne se généralise guère par une
filtration de $\pi_{1}$, comme on le voit sur le spectre d'un corps
fini.\footnote{Le passage des racines de l'unité au quotient fini demande la
semi-simplicité de l'action, que la page suppose pour la question suivante.}
Pour $i \geqslant 1$ : soit $\xi^{\beta}$ la plus grande puissance telle que
$E(\beta)$ soit de degré $i-1$ ou $i$ selon la parité de $\alpha$ ; on demande
$E(\beta)$ effectif. Enfin, pour $E$ semi-simple, la composante $E_{i}$ de niveau
$i$ se caractérise « de façon purement combinatoire en termes de Frobenius » ; le
bas de la page 54, qui le montrait sur un corps fini, est barré.

\pagerange{55}{58}
\subsection*{14) Invariants de Galois, théorèmes de commutation, et le motif de Lie (pages 55, 56 et 58)}

Soit $X$ propre et lisse sur un corps $K$ de type fini, $\Pi =
\operatorname{Gal}(\overline{K}/K)$. Le groupe $H^{2i}(X_{\overline{K}},
\mathbb{Q}_{\ell}(i))$ est de poids $0$ et peut avoir des invariants ; la
\emph{conjecture de Tate} dit
\[
\mathcal{A}^{i}(X) \otimes_{\mathbb{Q}} \mathbb{Q}_{\ell} \;\xrightarrow{\ \sim\ }\;
H^{2i}(X_{\overline{K}}, \mathbb{Q}_{\ell}(i))^{\Pi} ,
\]
$\mathcal{A}^{i}(X)$ désignant les cycles de codimension $i$ définis sur $K$,
modulo l'équivalence homologique. Il la précise ainsi. L'anneau
$\mathcal{U}^{i}(X)$ des correspondances algébriques qui opèrent sur $H^{i}$ est
semi-simple --- « via Hodge », c'est-à-dire par la conjecture standard de type
Hodge --- et l'on veut
\begin{enumerate}
  \item que $\mathcal{U}^{i}(X) \otimes \mathbb{Q}_{\ell}$ soit le commutant de
  $\Pi$ dans $\operatorname{End}_{\mathbb{Q}_{\ell}}(H^{i}(X_{\overline{K}},
  \mathbb{Q}_{\ell}))$ ;
  \item que l'algèbre enveloppante de $\Pi$ soit le commutant de
  $\mathcal{U}^{i}(X)$ --- de sorte que l'action de $\Pi$ est semi-simple.
\end{enumerate}
Le second énoncé, joint à l'injectivité de $\mathcal{U}^{i}(X) \otimes
\mathbb{Q}_{\ell} \to \operatorname{End}(H^{i})$, entraîne tout : c'est le
théorème du double commutant, l'algèbre semi-simple $\mathcal{U}^{i} \otimes
\mathbb{Q}_{\ell}$ étant égale à son bicommutant.\footnote{L'argument est
complet dès que $\mathcal{U}^{i}(X)$ est semi-simple, ce qui est l'hypothèse
« via Hodge ». La page ne nomme pas le théorème du double commutant.} Ce qu'on
caractérise ainsi, note-t-il, ce ne sont ni les classes algébriques elles-mêmes
ni l'image de $\Pi$, mais seulement les $\mathbb{Q}_{\ell}$-espaces qu'elles
engendrent et l'algèbre enveloppante. Passant à la limite sur les extensions
finies de $K$, on remplace $\Pi$ par l'image $\mathfrak{G}_{\ell}$ de Galois et
les cycles de $X$ par ceux de $X_{\overline{K}}$. Et une question, restée ouverte
sous cette forme : comment caractériser l'image de $\mathfrak{G}_{\ell}$ dans
$\operatorname{End}(H^{i})$ ? Ce doit être une algèbre de Lie
réductive.\footnote{La réductivité est une conséquence de la semi-simplicité
conjecturée. La question de l'image elle-même est aujourd'hui celle de la
conjecture de Mumford--Tate et de ses analogues, ouverte en général.}

La page 58 tire de là un objet. Pour un motif $M$ sur $K$, semi-simple, l'image
de $\mathfrak{G}_{\ell}$ dans $\operatorname{End}(M_{\ell}) = \check{M}_{\ell}
\otimes M_{\ell}$ doit provenir d'un motif $\mathfrak{G}_{M} \subset \check{M}
\otimes M$, indépendant de $\ell$, qui est un motif en algèbres de Lie ; quand
$M$ grandit on obtient un pro-motif en algèbres de Lie canonique, qu'il appelle
le \emph{pro-motif fondamental}, et qui doit pouvoir se construire à partir de la
seule catégorie des motifs avec ses structures habituelles.\footnote{C'est
l'algèbre de Lie du groupe de Galois motivique vue comme objet de la catégorie :
le groupe de Galois motivique opère sur sa propre algèbre de Lie par l'action
adjointe, qui est donc une représentation, c'est-à-dire un (pro-)motif ; sa
réalisation est $\operatorname{Lie} G_{M} \subset \operatorname{End}(\omega(M))$.
Dans le vocabulaire de Deligne (1989-1990), c'est l'algèbre de Lie du groupe
fondamental de la catégorie tannakienne. Que l'image de Galois l'atteigne est une
forme de la conjecture de Tate.} Quelle en est la « partie algébrique » ? La fin
de la page, sur le cas d'un corps fini, est barrée et ne se lit qu'en fragments.

\pagerange{59}{59}
\subsection*{15) La cohomologie absolue : les exigences (page 59)}

En plus des foncteurs relatifs $R^{i}f_{*}$, il faut des foncteurs « absolus »
$H^{i}(X, M)$, provenant d'un foncteur de catégories triangulées
\[
R\Gamma_{X} \colon D^{b}(\mathcal{M}^{+}(X)) \longrightarrow D^{b}(\mathrm{Ab}) ,
\]
qui commutent aux limites projectives filtrantes de préschémas à morphismes de
transition affines --- ce qui ramène au cas de type fini sur $\mathbb{Z}$ --- et
satisfont la transitivité : pour $f \colon X \to Y$ de type fini, $R\Gamma_{X} =
R\Gamma_{Y} \circ Rf_{*}$, d'où une suite spectrale de Leray
$H^{p}(Y, R^{q}f_{*}M) \Rightarrow H^{p+q}(X, M)$.\footnote{C'est ce qu'on appelle
aujourd'hui la \emph{cohomologie motivique}, que Beilinson placera au centre de
ses conjectures (1985) ; mais on verra aux pages 63 à 68 que les propriétés
demandées ici ne sont pas toutes celles de la cohomologie motivique de
Voevodsky.}

\pagerange{61}{79}
\section*{V. Cohomologie absolue, dualité, et points rationnels des variétés abéliennes (pages 61 à 79)}

La suite du paragraphe 15, puis le paragraphe 16. Les pages impaires 65 à 71 et
la page 80 sont des épreuves d'EGA IV sur lesquelles il a écrit au verso, et
n'entrent pas ici.

\pagerange{61}{64}
\subsection*{Les propriétés demandées (pages 61 à 64)}

On veut $H^{i}(X, M) = 0$ pour $i < 0$ et $H^{0}(X, M) = \operatorname{Hom}(
\mathbb{Q}_{X}(0), M)$. Rien n'empêche d'introduire des hyper-Ext,
\[
\operatorname{Ext}^{i}(X ; M, N) = \mathbb{H}^{i}\bigl(X,
R\underline{\operatorname{Hom}}(M, N)\bigr) ,
\]
avec $\operatorname{Ext}^{0} = \operatorname{Hom}$, et l'on espère que
$\operatorname{Ext}^{1}$ classe les extensions de $M$ par $N$ ; on aimerait que
les $\operatorname{Ext}^{*}$ soient ceux de la catégorie des motifs, mais il a des
doutes pour $i > 0$, peut-être faut-il passer aux Ind-motifs. Il
note les suites spectrales $H^{p}(X, H^{q}(M^{\bullet})) \Rightarrow
\mathbb{H}^{*}(X, M^{\bullet})$ et $H^{p}(X, \underline{\operatorname{Ext}}^{q}(M,N))
\Rightarrow \operatorname{Ext}^{*}(X ; M, N)$.\footnote{Que la cohomologie
motivique soit un groupe d'extensions dans une catégorie abélienne de motifs
mixtes est l'une des conjectures de Beilinson ; elle reste ouverte. La page pose
la question et ne la tranche pas.}

Puis la comparaison $\ell$-adique : les flèches
\[
H^{i}(X, M) \otimes_{\mathbb{Q}} \mathbb{Q}_{\ell} \longrightarrow
H^{i}(X, T_{\ell}(M)) ,
\qquad
\operatorname{Ext}^{i}(X ; M, N) \otimes \mathbb{Q}_{\ell} \longrightarrow
\operatorname{Ext}^{i}(X ; T_{\ell}M, T_{\ell}N)
\]
doivent être des isomorphismes pour $X$ de type fini sur $\mathbb{Z}$, et
$H^{*}(X, M)$ un $\mathbb{Q}$-espace de dimension finie. En particulier, sur un
corps algébriquement clos, $H^{i}(k, M) = 0$ pour $i > 0$ et tout
motif.\footnote{Ces exigences ne sont pas celles de la cohomologie motivique de
Voevodsky : sur $\mathbb{C}$, $H^{1}_{\mathcal{M}}(\mathbb{C}, \mathbb{Q}(1)) =
\mathbb{C}^{*} \otimes \mathbb{Q}$ n'est pas nul, et sur $\mathbb{F}_{p}$,
$H^{1}_{\mathcal{M}}(\mathbb{F}_{p}, \mathbb{Q}(0)) = 0$ alors que la page 68
demande $H^{1}(\mathbb{F}_{p}, \mathbb{Q}(0)) \simeq \mathbb{Q}$. La cohomologie
demandée ici, qui sur un corps fini a la taille de la cohomologie galoisienne de
$\hat{\mathbb{Z}}$ mais des coefficients rationnels, est plus proche de ce que
Lichtenbaum appellera la cohomologie Weil-étale (2005). Le rapprochement est de
nous.} Le théorème de Tate, qu'il appelle ainsi, se reformule alors : pour $X$
projective et lisse sur $k$ algébriquement clos,
\[
a^{i}(X) \otimes_{\mathbb{Z}} \mathbb{Q} \;\simeq\;
\operatorname{Im}\bigl[ H^{2i}(X, \mathbb{Q}(i)) \longrightarrow
H^{2i}(X, \mathbb{Q}_{\ell}(i)) \bigr] ,
\]
où $a^{i}(X)$ est le groupe des classes de cycles de codimension
$i$.\footnote{La page 63 donne d'abord $a^{i}(X) \otimes \mathbb{Q} \simeq
\Gamma_{k}(R^{i}f_{*}\mathbb{Q}_{X}(i))$, avec l'exposant $i$ là où la formule
encadrée de la page 64 porte $2i$ ; on lit $2i$.} Pour $j \neq 2i$, l'image de
$H^{j}(X, \mathbb{Q}(i))$ dans la cohomologie $\ell$-adique est nulle, pour une
raison de poids.

Et, pour $k = \overline{\mathbb{F}}_{p}$, une table :
\[
H^{i}(X, \mathbb{Q}(j)) =
\begin{cases}
0 & i \neq 2j, 2j+1 , \\
a^{j}(X) \otimes \mathbb{Q} & i = 2j \quad (\text{canoniquement}) , \\
a^{j}(X) \otimes \mathbb{Q} & i = 2j+1 \quad (\text{non canoniquement}) .
\end{cases}
\]
La page 70 la justifie comme limite sur les extensions finies de $\mathbb{F}_{p}$
(voir plus bas). Elle n'est pas compatible, telle quelle, avec l'annulation
demandée à la page 63 : sur un corps algébriquement clos, $H^{i}(k, -) = 0$ pour
$i > 0$ et la suite spectrale de Leray donneraient $H^{2j+1}(X, \mathbb{Q}(j)) =
\operatorname{Hom}(\mathbf{1}, R^{2j+1}f_{*}\mathbb{Q}(j)) = 0$, puisque
$R^{2j+1}f_{*}\mathbb{Q}(j)$ est de poids $1$.\footnote{La tension est entre deux
définitions de la cohomologie absolue sur $\overline{\mathbb{F}}_{p}$ : directement,
ou comme limite inductive des cohomologies sur les corps finis. Le dossier ne la
remarque pas et ne la tranche pas ; on la signale sans la résoudre.}

\pagerange{66}{68}
\subsection*{La dualité absolue (pages 66 et 68)}

Soit $X$ régulier de dimension $n$, de type fini sur $\mathbb{Z}$. Le motif
dualisant est $\mathbb{Q}(n)[2n+1]$, et pour deux motifs constructibles $M$,
$M^{*} = R\underline{\operatorname{Hom}}(M, \mathbb{Q}(n))$, les complexes
$R\Gamma_{X}(M)$ et $R\Gamma_{c}(X, M^{*})$ doivent être duaux ; les coefficients
étant rationnels, sans Tor, cela revient à
\[
H^{i}(X, M) \ \text{dual de}\ \operatorname{Ext}_{c}^{2n+1-i}(X ; M, \mathbb{Q}(n)) ,
\qquad
H^{i}_{c}(X, M) \ \text{dual de}\ \operatorname{Ext}^{2n+1-i}(X ; M, \mathbb{Q}(n)) ,
\]
et les deux coïncident si $X$ est propre sur $\operatorname{Spec}\mathbb{Z}$.
C'est la forme motivique de la dualité d'Artin--Verdier : pour $n = 1$, le
dualisant $\mathbb{Q}(1)[3]$ est $\mathbb{G}_{m}[2]$ à coefficients rationnels.
Par la dualité relative, il la ramène au cas $X = \operatorname{Spec}\mathbb{Z}$,
puis au cas $X = \operatorname{Spec}\mathbb{F}_{p}$, où il s'agit de la dualité
entre $H^{i}(\mathbb{F}_{p}, M)$ et $H^{1-i}(\mathbb{F}_{p}, M^{*})$ par le
cup-produit et $H^{1}(\mathbb{F}_{p}, \mathbb{Q}(0)) \simeq \mathbb{Q}$ ; il suffit
de la vérifier $\ell$-adiquement, où c'est la trivialité que $H^{i}(\hat{\mathbb{Z}},
V)$ et $H^{1-i}(\hat{\mathbb{Z}}, V^{*})$ sont duaux --- invariants et coinvariants
d'un seul endomorphisme.\footnote{Le premier pas est la dualité relative ; le
second ne peut pas être formel : la dualité d'Artin--Verdier pour
$\operatorname{Spec}\mathbb{Z}$ contient de la théorie du corps de classes global
et ne se ramène pas aux seuls points fermés. La page écrit « et lui-même le cas
géométrique », sans argument. La marge ajoute « à développer au sujet des
cup-produits à coefficients abstraits ».}

\pagerange{70}{72}
\subsection*{Sur un corps fini : la suite de Hochschild--Serre et les poids (pages 70 et 72)}

Soit $X$ projective, lisse, géométriquement connexe sur $k$ fini à $q$ éléments,
$\pi = \operatorname{Gal}(\overline{k}/k)$. La suite de Hochschild--Serre donne
\[
0 \to H^{i-1}(\overline{X}, \mathbb{Q}_{\ell}(j))_{\pi} \to H^{i}(X, \mathbb{Q}_{\ell}(j))
\to H^{i}(\overline{X}, \mathbb{Q}_{\ell}(j))^{\pi} \to 0 ,
\]
et comme $H^{m}(\overline{X}, \mathbb{Q}_{\ell}(j))$ est pur de poids $m - 2j$, les
invariants et coinvariants ne peuvent être non nuls qu'en poids $0$ : $H^{i}(X,
\mathbb{Q}_{\ell}(j)) = 0$ sauf pour $i = 2j$ et $i = 2j+1$. Pour $i = 2j$ on trouve
$H^{2j}(\overline{X}, \mathbb{Q}_{\ell}(j))^{\pi}$, qui est $a^{j}(X) \otimes
\mathbb{Q}_{\ell}$ par la conjecture de Tate ; pour $i = 2j+1$, les coinvariants
$H^{2j}(\overline{X}, \mathbb{Q}_{\ell}(j))_{\pi}$, qui s'identifient aux invariants
si l'action est semi-simple, donc encore à $a^{j}(X) \otimes
\mathbb{Q}_{\ell}$.\footnote{La conjecture de Tate sur les corps finis est ouverte
en général ; elle est démontrée pour les diviseurs sur les variétés abéliennes
(Tate, 1966), cas dont la page 77 va se servir.} Passant à la limite sur les
extensions $k_{n}$, on obtient $0$ hors des degrés $2j$, $2j+1$, et $a^{j}(\overline{X})
\otimes \mathbb{Q}_{\ell}$ dans les deux, non canoniquement pour le second : c'est
la table de la page 64.

La page 72 fait le compte des poids pour $S$ lisse de dimension $n$, de type
fini sur $\mathbb{F}_{p}$, et $M$ pur de poids $\rho$. Dans la suite de
Hochschild--Serre, $H^{i-1}(\overline{S}, \overline{M}_{\ell})_{\pi}$ a ses poids
entre $\rho + i - 1$ et $\rho + \inf(2i-2, 2n)$, et $H^{i}(\overline{S},
\overline{M}_{\ell})^{\pi}$ entre $\rho + i$ et $\rho + \inf(2i, 2n)$ ; il faut le
poids $0$ dans l'un des deux intervalles. D'où la condition nécessaire, pour que
$H^{i}(S, M_{\ell}) \neq 0$ :
\[
-2i \leqslant \rho \leqslant 1 - i \quad (i \geqslant 1) ,
\qquad
\rho = 0 \quad (i = 0) .
\]\footnote{La page donne $-2i \leqslant \rho \leqslant -i+1$ pour $i \geqslant 2$,
mais $-2 \leqslant \rho \leqslant 1$ pour $i = 1$. Pour $i = 1$ le premier intervalle
est réduit à $\rho = 0$ ($H^{0}(\overline{S})$ est de poids $\rho$) et le second est
$[-2, -1]$ : la réunion est $[-2, 0]$, que la formule générale donne aussi. On a
unifié. Les bornes de poids utilisées sont celles de Weil II et, pour la borne
supérieure en $2n$, des arguments de résolution --- la page les suppose.}

\pagerange{73}{79}
\subsection*{16) Points rationnels et cohomologie des variétés abéliennes (pages 73 à 79)}

Soit $A_{K}$ une variété abélienne sur un corps $K$ de type fini sur son corps
premier. $K$ est le corps des fonctions d'un schéma $S$ intègre et régulier, de
type fini sur $\mathbb{Z}$, sur lequel $A_{K}$ s'étend en un schéma abélien $A$,
et $A_{K}(K) = A(S)$ est de type fini par le théorème de
Mordell--Weil--Néron--Lang. Le motif
\[
T_{\mathrm{mot}}(A) = R^{1}f_{*}\mathbb{Q}_{A}(1) ,
\]
pour $f \colon A \to S$, est pur de poids $-1$ ; il redonne $A_{K}$ à isogénie
près, par les conjectures de Tate.\footnote{La page le dit « pur de poids $1$ ».
$R^{1}f_{*}$ est de poids $1$ et la torsion par $\mathbb{Q}(1)$ retire $2$ : le
poids est $-1$, celui d'un module de Tate, et c'est ce poids que le raisonnement
suivant utilise, puisqu'il annule $H^{0}$. Sa réalisation $\ell$-adique est
$V_{\ell}$ de la variété duale, isogène à $A$.} On attend donc
\[
A(S) \otimes_{\mathbb{Z}} \mathbb{Q} \;\simeq\; H^{1}(S, T_{\mathrm{mot}}(A)) ,
\]
avec $H^{0}(S, T_{\mathrm{mot}}(A)) = 0$ pour des raisons de poids.\footnote{C'est,
dans le langage d'aujourd'hui, l'attente que le groupe de Mordell--Weil
rationalisé soit un $\operatorname{Ext}^{1}$ de motifs mixtes, $\operatorname{Ext}^{1}(
\mathbb{Q}(0), h_{1}(A))$ --- la forme que Beilinson et Bloch--Kato donneront à
cette heuristique. Elle est conjecturale.} La suite de Kummer $0 \to
{}_{\ell^{\nu}}A \to A \to A \to 0$ donne, par passage à la limite,
\[
0 \longrightarrow A(S) \otimes \mathbb{Q}_{\ell} \longrightarrow
H^{1}(S, V_{\ell}A) \longrightarrow V_{\ell}\bigl(H^{1}(S, A)\bigr) \longrightarrow 0 ,
\]
et le résultat escompté équivaut vraisemblablement à
\[
H^{1}(S, A)(\ell) \ \text{fini pour } \ell \text{ premier aux caractéristiques
résiduelles de } S .
\]
C'est vrai par Lang pour $S = \operatorname{Spec}\mathbb{F}_{q}$ ; en général,
c'est la finitude de la composante $\ell$-primaire du groupe de
Tate--Šafarevič.\footnote{Ouverte en général. Pour les courbes elliptiques sur
$\mathbb{Q}$ de rang analytique $\leqslant 1$, elle résulte des travaux de
Gross--Zagier et Kolyvagin (vers 1988-1990).}

Pour $S$ de type fini sur $\mathbb{F}_{p}$, les suites de Leray pour $S \to
\operatorname{Spec}\mathbb{F}_{p}$ et de Kummer, jointes à la suite $0 \to
A^{\mathrm{tr}} \to R^{0}f_{*}A \to M \to 0$ où $A^{\mathrm{tr}}$ est la trace
de $A$ sur $\mathbb{F}_{p}$ et $M = H^{0}(\overline{S}, \overline{A})/
H^{0}(\overline{\mathbb{F}}_{p}, \overline{A}^{\mathrm{tr}})$ un $\pi$-module de
type fini sur $\mathbb{Z}$, donnent le diagramme qu'il nomme (D) :

\begin{tikzcd}[column sep=small, row sep=small, nodes={font=\scriptsize}]
0 \arrow[r] & H^{0}(S, A) \otimes \mathbb{Q}_{\ell}/\mathbb{Z}_{\ell} \arrow[r] \arrow[d] & H^{1}(S, {}_{\ell^{\infty}}A) \arrow[r] \arrow[d] & H^{1}(S, A)(\ell) \arrow[r] \arrow[d] & 0 \\
0 \arrow[r] & (H^{0}(\overline{S}, \overline{A}) \otimes \mathbb{Q}_{\ell}/\mathbb{Z}_{\ell})^{\pi} \arrow[r] & H^{1}(\overline{S}, {}_{\ell^{\infty}}\overline{A})^{\pi} \arrow[r] & H^{1}(\overline{S}, \overline{A})(\ell)^{\pi} \arrow[r] & 0
\end{tikzcd}

La première flèche verticale est un isomorphisme modulo des groupes finis,
$H^{0}(S, A) \simeq M^{\pi}$ et $H^{0}(\overline{S}, \overline{A}) \otimes
\mathbb{Q}_{\ell}/\mathbb{Z}_{\ell} \simeq M \otimes \mathbb{Q}_{\ell}/\mathbb{Z}_{\ell}$
étant standard. Et la suite de Leray, avec $H^{1}(\mathbb{F}_{p},
A^{\mathrm{tr}}) = 0$ par Lang et $H^{1}(\mathbb{F}_{p}, M)$ fini --- il s'agit de
la cohomologie étale, c'est-à-dire galoisienne de $\hat{\mathbb{Z}}$, insiste-t-il
---, montre que la dernière flèche verticale est bijective à noyau fini près :
\[
H^{1}(S, A)(\ell) \longrightarrow H^{1}(\overline{S}, \overline{A})(\ell)^{\pi} .
\]\footnote{La phrase de la page 76 qui établit ce point s'arrête sur « et le »,
et reprend en tête de la page 78 ; la page 77 est un ajout appelé depuis la marge
de la page 78.} La conjecture signifie donc aussi que $H^{1}(\overline{S},
\overline{A})(\ell)^{\pi}$ est fini --- et, à peu près, par les jacobiennes, que
pour $X \to S$ propre et lisse de dimension relative $1$ et $S$ régulier de type
fini sur $\mathbb{Z}$, l'application $H^{2}(S, \mathbb{G}_{m}) \to H^{2}(X,
\mathbb{G}_{m})$ est surjective modulo des groupes finis --- question que la page
laisse avec deux points d'interrogation.\footnote{C'est la relation entre le
groupe de Brauer d'une surface fibrée et le groupe de Tate--Šafarevič de la
jacobienne de sa fibre générique, qu'établit, pour une surface fibrée sur une
courbe, « Le groupe de Brauer III » (Dix exposés, 1968). La page ne cite rien. Une
phrase biffée juste avant dit : « Je vois une façon de la prouver ».}

Il note encore que $H^{i}(S, A)(\ell)$ devrait être fini pour $i \geqslant 3$, et
que $T_{\ell}(H^{i}(S, A)(\ell)) \simeq H^{i}(S, T_{\ell}(A))$ pour $i \geqslant 2$,
de sorte qu'il ne reste à examiner que $H^{0}$ et $H^{2}$ pour $A$, $H^{1}$ et
$H^{2}$ pour $T_{\ell}(A)$. Et l'ajout de la page 77 montre que $H^{2}$ n'est pas
nul en général : pour $S$ une courbe propre, lisse, géométriquement connexe sur
$\mathbb{F}_{p}$ et $A = J \times S$, où $J$ est la jacobienne de $S$,
\[
H^{2}(S, T_{\ell}(A)) \simeq H^{1}(\overline{S}, T_{\ell}(\overline{A}))_{\pi}
= \bigl[T_{\ell}(\overline{J}) \otimes T_{\ell}(\overline{J})\bigr]_{\pi}
\simeq \operatorname{End}(J) \otimes \mathbb{Q}_{\ell} ,
\]
les classes de correspondances divisorielles sur $S \times S$, qui ne sont pas
nulles.\footnote{Ici rien n'est conjectural : l'isomorphisme
$\operatorname{End}(J) \otimes \mathbb{Q}_{\ell} \simeq
\operatorname{End}_{\pi}(V_{\ell}J)$ et la semi-simplicité qui identifie invariants
et coinvariants sont les théorèmes de Tate (1966).}

Enfin, $H^{1}(S, T_{\ell}(A)) \simeq H^{1}(\overline{S}, T_{\ell}(\overline{A}))^{\pi}$,
de sorte que la conjecture signifie essentiellement
\[
\operatorname{rang}_{\mathbb{Q}_{\ell}} H^{1}(\overline{S}, V_{\ell}(\overline{A}))^{\pi}
= \operatorname{rang}_{\mathbb{Z}} A(K) ,
\]
et, sous cette forme, elle est essentiellement équivalente à la conjecture de
\emph{Birch et Swinnerton-Dyer}.\footnote{Sur un corps de fonctions d'une
variable sur un corps fini, l'équivalence entre la finitude de la composante
$\ell$-primaire du groupe de Tate--Šafarevič et la conjecture de Birch et
Swinnerton-Dyer est due à Tate (Séminaire Bourbaki, 1966) et Milne (1975). Elle
n'est démontrée que dans des cas particuliers. La page écrit l'égalité des rangs
avec un « $\simeq$ » souligné.}

\pagerange{81}{83}
\section*{VI. Deux feuillets isolés (pages 81 et 83)}

\pagerange{81}{81}
\subsection*{Une courbe sur un corps fini (page 81)}

Soit $C$ une courbe propre et lisse sur $\mathbb{F}_{q}$, $A$ un schéma abélien
sur $C$, $x$ un point fermé, $U = C \smallsetminus \{x\}$, $K_{x}$ le corps
local en $x$. La suite de localisation
\[
0 \longrightarrow H^{1}(C, A) \longrightarrow H^{1}(U, A) \longrightarrow
H^{2}_{x}(C, A) \longrightarrow H^{2}(C, A)
\]
--- l'injectivité à gauche vient de ce que $A(\mathcal{O}^{h}_{x}) =
A(K^{h}_{x})$ --- et l'excision $H^{2}_{x}(C, A) \simeq H^{1}(K_{x}, A)$, dual de
$A^{*}(K_{x})$ par la dualité locale de Tate, donnent : $H^{1}(U, A)(\ell)$ est fini
pour $\ell \neq p$, dès que $H^{1}(C, A)(\ell)$ l'est. Il se demande si l'on peut
généraliser au cas d'une variété régulière dont on enlève un diviseur
régulier.\footnote{La page va plus vite, et trop : elle affirme $H^{i}(C, A) = 0$
pour $i \geqslant 2$ « des conjectures de Weil », et « sans doute, utilisant Tate »,
$H^{1}(C, A) = 0$, d'où $H^{1}(U, A) \simeq H^{2}_{x}(C, A)$, « groupe fini !
OK ». Trois corrections. $H^{1}(C, A)$ est le groupe de Tate--Šafarevič, fini
conjecturalement mais non nul en général. $H^{2}(C, A)$ n'est pas nul : par la
dualité de Milne, il est dual de $A^{*}(K)$ --- et la page 77 de ce dossier le
montre elle-même, avec $H^{2}(S, T_{\ell}(A)) \neq 0$ pour $A = J \times S$. Enfin
$H^{1}(K_{x}, A)$, dual du groupe profini $A^{*}(K_{x})$, n'est fini que dans sa
partie première à $p$. L'argument ci-dessus n'a besoin d'aucun des trois
énoncés fautifs.}

\pagerange{83}{83}
\subsection*{Calculs sur $\mathbb{Q}_{\ell}(1)$ et $\mathbb{G}_{m}$ (page 83)}

Pour $X$ de type fini sur $\mathbb{Z}$, $\ell$ inversible, la suite de Kummer
donne : $H^{0}(X, \mathbb{Q}_{\ell}(1)) = 0$ ; $H^{1}(X, \mathbb{Q}_{\ell}(1)) =
\mathcal{O}(X)^{*} \otimes \mathbb{Q}_{\ell}$ ; $H^{2}(X, \mathbb{Q}_{\ell}(1))$
extension de $T_{\ell}(H^{2}(X, \mathbb{G}_{m}))$ par $\operatorname{Pic}(X) \otimes
\mathbb{Q}_{\ell}$ ; $H^{3}(X, \mathbb{Q}_{\ell}(1)) \simeq T_{\ell}(H^{3}(X,
\mathbb{G}_{m}))$, qui peut valoir $\mathbb{Z}_{\ell}$, par exemple pour un modèle
d'un corps de fonctions d'une variable sur $\mathbb{F}_{p}$ ; et il demande si
$H^{i}(X, \mathbb{Q}_{\ell}(1)) = T_{\ell}(H^{i}(X, \mathbb{G}_{m}))$ est nul pour
$i \geqslant 4$.\footnote{L'annulation de $H^{0}$ vient de ce que $X$, de type fini
sur $\mathbb{Z}$, n'a qu'un nombre fini de racines de l'unité ; celle de
$T_{\ell}\operatorname{Pic}$ de la finitude de $\operatorname{Pic}(X)[\ell^{\infty}]$.
La marge appelle « coBrauer » le terme $T_{\ell}H^{2}(X, \mathbb{G}_{m})$.} Il
rappelle la conjecture d'Artin, que le groupe de Brauer est fini pour $X$ propre
sur $\mathbb{Z}$ --- vrai pour une courbe projective et lisse sur $\mathbb{F}_{p}$,
faux dès qu'on lui enlève des points, le corang du $H^{2}$ augmentant
d'autant.\footnote{La conjecture d'Artin est ouverte en général ; pour une courbe
sur un corps fini, $\operatorname{Br}(C) = 0$ par la théorie du corps de classes.}

\pagerange{85}{88}
\section*{VII. 17) Formes positives (pages 85 à 88)}

C'est le point 17 du plan, et la première mouture de la positivité ; la seconde
est aux pages 157 à 166. Le théorème étant « géométrique », on se place sur un
corps $k$ algébriquement clos, et l'on ne considère que des motifs semi-simples
purs de poids $p$. Pour un tel $M$, les formes $M \times M \to \mathbb{Q}(-p)$
sont les éléments de
\[
\mathcal{L}(M) = \Gamma_{k}\bigl(\underline{\operatorname{Hom}}(M \otimes M,
\mathbb{Q}(-p))\bigr) ,
\]
$\Gamma_{k} = \operatorname{Hom}(\mathbf{1}, -)$, le $\underline{\operatorname{Hom}}$
étant bien de poids $0$ ; c'est un $\mathbb{Q}$-espace de dimension finie, dans
lequel on considère le sous-espace $\mathcal{L}^{s}(M)$ des formes symétriques si
$p$ est pair, alternées si $p$ est impair. Le produit tensoriel des formes donne
des injections $\mathcal{L}^{s}(M) \otimes \mathcal{L}^{s}(N) \to \mathcal{L}^{s}(M
\otimes N)$, non bijectives en général --- le produit de deux formes de même
symétrie est symétrique, de symétries opposées antisymétrique ---, et
\[
\mathcal{L}^{s}(M' \oplus M'') \simeq \mathcal{L}^{s}(M') \times \mathcal{L}^{s}(M'')
\times \mathcal{B}(M', M'') ,
\qquad
\mathcal{B}(M', M'') = \operatorname{Hom}(M' \otimes M'', \mathbb{Q}(-p)) .
\]\footnote{La page écrit $\mathcal{L}^{s}(M' \times M'')$ pour la somme directe et
$\mathbb{Q}(p)$ dans $\mathcal{B}$, là où la page 85 a $\mathbb{Q}(-p)$ ; les
formes sont à valeurs dans $\mathbb{Q}(-p)$.}

On se donne dans $\mathcal{L}^{s}(M)$ un cône convexe $\mathcal{L}^{+}(M)$ ---
ses éléments sont les formes \emph{positives}, ceux de son intérieur les formes
\emph{amples} --- soumis à huit axiomes :
\begin{enumerate}
  \item[(i)] $\mathcal{L}^{+}(M)$ est d'intérieur non vide ;
  \item[(ii)] un élément de $\mathcal{L}^{+}(M)$ est intérieur si et seulement s'il
  est non dégénéré ;
  \item[(iii)] tout morphisme $M \to N$ envoie $\mathcal{L}^{+}(N)$ dans
  $\mathcal{L}^{+}(M)$ ;
  \item[(iv)] si $M \to N$ est injectif, il envoie formes amples sur formes amples ;
  en particulier la restriction d'une forme ample à $M$ est non dégénérée, et
  $N = M \oplus M^{\perp}$ ;
  \item[(v)] la somme directe de deux formes amples est ample ;
  \item[(vi)] le produit tensoriel de deux formes amples est ample ;
  \item[(vii)] si $p = 0$, de sorte que $M$ est essentiellement un
  $\mathbb{Q}$-espace vectoriel de dimension finie, les formes amples sont les
  formes définies positives ;
  \item[(viii)] $\mathcal{L}^{s}$ d'un motif de Tate est $\mathbb{Q}$, et les formes
  amples y sont les nombres strictement positifs.
\end{enumerate}
De (vi) et (viii) il résulte que l'isomorphisme $\mathcal{L}^{s}(M) \simeq
\mathcal{L}^{s}(M(n))$ respecte les cônes ; de (iv) et (vii), que si $p = 0$ une
forme ample sur $M$ induit une forme définie positive sur $\Gamma_{k}M$, par
restriction au sous-motif $\Gamma_{k}M \otimes \mathbf{1} \subset M$.\footnote{La
page renvoie à (iv) et (vi) --- lecture incertaine ; l'argument demande (iv) et
(vii).} Enfin, décomposant $M$ en composantes isotypiques $M_{\alpha}$, on a
$\mathcal{L}^{s}(M) \simeq \bigoplus_{\alpha}\mathcal{L}^{s}(M_{\alpha})$, et il suffit
d'étudier les formes positives sur les morceaux isotypiques $I \otimes_{A} V$, où
$I$ est simple, $A = \operatorname{End}(I)$ un corps gauche, et $V$ un $A$-espace
vectoriel.\footnote{La décomposition suppose qu'il n'y ait pas de forme non nulle
entre deux composantes isotypiques distinctes, c'est-à-dire que chaque
$M_{\alpha}$ soit isomorphe à son dual tordu $M_{\alpha}^{\vee}(-p)$ --- ce que
donne une polarisation. L'hypothèse n'est pas écrite.}

Ces axiomes sont ceux que satisfont les formes de polarisation des structures de
Hodge, et que la conjecture standard de type Hodge doit fournir pour les
motifs.\footnote{La notion de catégorie tannakienne polarisée sera axiomatisée
par Saavedra (1972) et reprise par Deligne--Milne (1982). Le dossier 10, pages 3
à 16, contient une autre axiomatique des polarisations, par les formes de Weil ;
la parenté est de nous.}

\pagerange{89}{94}
\section*{VIII. 18) Fonctions $L$ et cohomologie à action galoisienne (pages 89 à 94)}

\pagerange{89}{91}
\subsection*{Le dictionnaire algébrique (pages 89 à 91)}

Soit $k$ un anneau et $\mathcal{C}$ la catégorie des $k[t]$-modules projectifs de
type fini sur $k$, c'est-à-dire des couples $(V, f)$ d'un tel module et d'un
endomorphisme --- des $k[\mathbb{N}]$-modules. Elle donne un $\lambda$-anneau
$R(\mathbb{N}; k)$, qui pour $k$ un corps est le $\mathbb{Z}$-module libre sur les
classes des $k[t]/(P)$, $P$ unitaire irréductible, le polynôme caractéristique de
$(k[t]/(P), t)$ étant $P$. Il préfère $Q(t) = t^{d}P(1/t) = \det(1 - tf)$, et
obtient
\[
R(\mathbb{N}; k) \longrightarrow 1 + t\,k[[t]] ,
\qquad
(V, f) \longmapsto L(V, f) = \frac{1}{\det(1 - tf)} ,
\]
homomorphisme de $\lambda$-anneaux vers le groupe multiplicatif $1 + t\,k[[t]]$, dont
l'image est formée de fonctions rationnelles valant $1$ en $0$. Il le dit injectif
pour $k$ un corps. Il l'est sur le sous-anneau engendré par les $(V, f)$ avec $f$
inversible, et c'est le seul cas qui servira --- le Frobenius opère
inversiblement sur la cohomologie ; mais les classes où $f$ est nilpotent ont
$L = 1$, et forment le noyau.\footnote{La page affirme l'injectivité sans
restriction, par factorialité de $k[t]$ ; l'argument oublie le facteur $P = t$,
pour lequel $Q = 1$. Que le quotient de l'anneau des endomorphismes par celui des
nilpotents s'identifie aux vecteurs de Witt rationnels est un théorème
d'Almkvist (1974). L'anneau $1 + t\,k[[t]]$ est l'anneau des grands vecteurs de
Witt, que la page ne nomme pas.} La dérivée logarithmique
\[
\Lambda(V, f) = \frac{L'(V, f)}{L(V, f)} = \sum_{n \geqslant 1}
(\operatorname{Tr} f^{n})\, t^{n-1}
\]
permet, en caractéristique nulle, de retrouver $L = \exp \int_{0}^{t} \Lambda$ ;
$\Lambda$ est alors injectif sur les classes à $f$ inversible, ce qu'il cesse d'être
en caractéristique $p > 0$.

\pagerange{91}{94}
\subsection*{Ce que la fonction $L$ sait de la cohomologie (pages 91 à 94)}

Par la formule des traces de Lefschetz--Weil--Grothendieck, la connaissance de la
fonction $L$ de $(X, F)$ sur $\mathbb{F}_{p}$ équivaut à celle de l'élément
\[
\chi(F) = \sum_{i} (-1)^{i} \operatorname{cl} H^{i}_{c}(\overline{X}, \overline{F})
\in R(\mathbb{N}; \mathbb{Q}_{\ell}) ,
\]
$f$ étant le Frobenius. Si l'on admet les conjectures de Weil et que $F$ vient
d'un motif, $\chi(F)$ se décompose suivant les poids en $\sum (-1)^{i}h^{i}(F)$, et
$L(X, F)$ équivaut aux $h^{i}(F)$ --- les zéros et pôles se groupant par valeur
absolue ; si $X$ est complète et lisse et $F$ lisse et pur, $h^{i}(F) =
H^{i}(\overline{X}, \overline{F})$ avec son Frobenius.

Pour $U \subset X$ à complémentaire de dimension $< n$, les $L(U, F|U)$ ne
diffèrent de $L(X, F)$, pour $F$ pur de poids $p$, que par des fonctions de
poids $\leqslant 2n-2+p$ : $L(X, F)$ modulo ces fonctions est un \emph{invariant
birationnel}, et équivaut à la connaissance de $\operatorname{cl}H^{2n}_{c} -
\operatorname{cl}H^{2n-1}_{c}$ modulo les poids bas. Pour $X$ lisse, connexe, de
dimension $n$, $H^{2n}_{c}(X, F)$ est dual de $H^{0}(X, \check{F})(n)$ et
$H^{2n-1}_{c}(X, F)$ de $H^{1}(X, \check{F}(n))$, qui a en général des poids
$2n-1$ et $2n-2$. Pour $F = \mathbb{Q}_{\ell}$, la connaissance birationnelle
grossière de la fonction zêta équivaut, modulo les conjectures de Tate, à celle
de $\pi_{0}(\overline{X})$ comme ensemble galoisien et de l'Albanese de $X$. Une
analyse plus fine, par le niveau et non seulement par le poids, doit récupérer
$\zeta(X)$ modulo les $\zeta(Y)$, $\dim Y < n$, c'est-à-dire les « invariants
birationnels fondamentaux » relativement à $\mathbb{F}_{p}$, « envisagés dans les
notes sur les \emph{Motifs} » : ce sont ceux du paragraphe 12, pages 42 à 51, que
la marge de la page 92 appelle son « N° 12 ».\footnote{Un passage encadré et
annulé de la page 92 commençait le même calcul pour une fibration $f \colon X \to
Y$, avec les ouverts denses de $Y$ ; les marges obliques des pages 92 et 93 sont
en grande partie perdues.}

Il propose enfin un programme : pour $f \colon X \to Y$ lisse et propre et $F$
localement constant, la connaissance de $L(X, F)$ modulo les $L(f^{-1}(Y'), F)$,
$Y'$ parcourant les fermés stricts de $Y$, permet-elle de récupérer les invariants
birationnels fondamentaux de la fibre générique sur le corps des fonctions de
$Y$ ? Et dans quelle mesure faut-il supposer $f$ lisse et propre et $F$
localement constant ? Le programme n'est pas repris.

\pagerange{96}{100}
\section*{IX. Poids et niveaux d'espaces de cohomologie $\ell$-adiques (pages 95 à 100)}

Une chemise de sa main, page 95, porte ces mots, suivis de « base var.\ quelconque »
. Dessous, un tapuscrit de deux pages, son plan manuscrit, un
feuillet anglais, et un calcul.

\pagerange{96}{97}
\subsection*{L'introduction d'un article (pages 96 et 97)}

« Filtration et poids des espaces de cohomologie des variétés algébriques, par
A.\ Grothendieck » : l'introduction, dactylographiée et corrigée de sa main. On y
lit le programme de toute une partie du dossier, et on le résume. Depuis
l'invention de la cohomologie $\ell$-adique, on sait que l'action du groupe de
Galois sur la cohomologie d'une variété définie sur un corps $k$ est une
structure supplémentaire remarquable, qui donne lieu aux conjectures de Tate ;
celles-ci semblent plus loin de leur solution que les conjectures de Weil. Le
propos est de tirer des conjectures de Weil et, au besoin, de la résolution des
singularités des conséquences en apparence \emph{géométriques} --- relatives à un
corps de base algébriquement clos --- en ramenant l'énoncé géométrique au cas
d'un corps fini. L'idée de ces raisonnements par les poids « remonte à 1964 » ;
des exemples analogues sont développés dans SGA 7, et le séminaire oral contenait
l'exemple du n° 2.

La portée des « résultats » est « sujette à caution », puisqu'ils dépendent des
conjectures de Weil et de la résolution des singularités « sous la forme forte de
Hironaka » en toute caractéristique. Mais, écrit-il, la compréhension ainsi
obtenue, « quoiqu'elle soit pour l'instant entièrement heuristique », peut
conduire à des propriétés importantes et inattendues, et elle « mérite pour cette
raison de ne pas être dédaignée ». En caractéristique nulle, le travail récent de
Deligne sur la théorie de Hodge des variétés quelconques justifie déjà, par voie
transcendante, le « yoga » des poids et des niveaux ; en caractéristique $p$ il
faudra attendre la preuve des conjectures de Weil, « ou mieux, des conjectures
standart ». Les filtrations par poids et niveaux, pour $\ell$ variable, doivent
donner des espaces de rang indépendant de $\ell$ : ce que la théorie conjecturale
des motifs exprime en disant qu'elles vivent déjà sur la cohomologie motivique ---
en accord avec la théorie de Deligne, où la filtration par le poids de la
cohomologie de De Rham provient d'une filtration à coefficients
rationnels.\footnote{Le contenu situe la frappe : après SGA 7 (séminaire de
1967-1969) et après l'annonce de la théorie de Hodge mixte de Deligne (Hodge II,
1971), avant la démonstration des conjectures de Weil (Deligne, Weil I, 1974),
qu'il dit « non prouvées à l'heure actuelle ». Les conjectures de Weil sont
démontrées depuis ; les conjectures standard ne le sont pas ; la résolution des
singularités en caractéristique $p$ reste ouverte, mais les altérations de de Jong
(1996) suffisent aux arguments de poids.}

\pagerange{98}{98}
\subsection*{Le plan (page 98)}

Le plan manuscrit du même article compte dix numéros : introduction ; notions de
poids et de niveau pour un module galoisien sur un corps ; un schéma privé d'un
diviseur à croisements normaux ; poids et niveaux de $H^{i}(X)$ pour $X$ lisse, de
$H^{i}_{c}(X)$ pour $X$ séparé, de $H^{i}(X)$ pour $X$ quelconque ; le cas propre,
avec les plus grands quotients de poids $i$ et la résolution ; le cas propre et
lisse, avec la filtration par le niveau ; les nombres de Betti virtuels de Serre ;
des contre-exemples ; les coefficients tordus. Suit une liste de références aux
crochets vides : SGA 4, Deligne (Hodge), Tate (conjectures), Kleiman (conjectures
standard), Weil, SGA 7.\footnote{Le numéro 8 figure deux fois, les numéros 5, 7
et 8 étant surchargés. Des dix numéros, seule l'introduction est rédigée ; les
notes des pages 103 à 121 en sont le matériau, et le « Contre-exemples divers »
est peut-être le calcul de la page 100.}

\pagerange{99}{99}
\subsection*{Un feuillet anglais sur les classes de cycles (page 99)}

Un feuillet dactylographié en anglais, « 2.\ Formalism of Algebraic Cohomology
Classes », dont l'auteur n'est pas nommé et dont les corrections manuscrites ne
sont pas sûrement de sa main. Il définit, pour $X$ projective et lisse de
dimension $n$, l'application classe de cycle $\rho_{X} \colon C^{i}(X) \to
H^{2i}_{\ell}(X)$ sur l'anneau de Chow rationnel : pour $Z$ irréductible de
codimension $i$, $\rho_{X}(Z)$ est l'unique classe telle que
\[
\kappa_{X}\bigl(\alpha \cdot \rho_{X}(Z)\bigr) = \kappa_{Z}\bigl(\operatorname{res}_{Z}\alpha\bigr)
\qquad \text{pour tout } \alpha \in H^{2n-2i}_{\ell}(X) ,
\]
$\kappa$ désignant les traces, ce qui a un sens par la dualité de Poincaré ; elle
transforme produit d'intersection en cup-produit, et son noyau définit
l'équivalence cohomologique.\footnote{Les torsions à la Tate sont omises, comme il
est d'usage dans les exposés de l'époque. La trace $\kappa_{Z}$ existe pour $Z$
irréductible propre quelconque, $H^{2d}(Z)$ étant de rang $1$. Le seul lien du
feuillet avec le dossier est le renvoi « Kleiman / conj.\ standart » du plan de la
page 98 ; on n'en tire aucune attribution.}

\pagerange{100}{100}
\subsection*{Un contre-exemple (page 100)}

Dans $\mathbb{P}^{2}$, soit $T$ la réunion de deux droites qui se coupent en un
point $c$, $a$ et $b$ deux points de $T$ distincts de $c$, un sur chaque droite,
$T' = \{a, b\}$, et $X = \mathbb{P}^{2} \smallsetminus T'$, $Y = T \smallsetminus
T'$, $U = \mathbb{P}^{2} \smallsetminus T$ ; $\widetilde{T}$ est la normalisée de
$T$, deux droites disjointes, et $\widetilde{Y}$ l'image inverse de $Y$. Il écrit
les suites exactes de la cohomologie à supports propres. Refait, le calcul donne :
$H^{1}_{c}(U) = 0$ parce que $T$ est connexe ; $H^{1}_{c}(X) \simeq
H^{0}(T')/H^{0}(\mathbb{P}^{2}) \simeq \mathbb{Q}$, de poids $0$ ; la flèche
$H^{1}_{c}(X) \to H^{1}_{c}(Y) \simeq \mathbb{Q}$ est injective ; mais
$H^{1}_{c}(\widetilde{Y}) = 0$, $\widetilde{Y}$ étant réunion disjointe de deux
droites affines. La classe de poids $0$ disparaît donc sur la normalisation, et
la réponse à la question « inj ? » qu'il inscrit sur la flèche vers la
normalisation est négative.\footnote{Ce calcul est de nous. La transcription note
que certaines flèches de la page sont d'un sens incertain, que la position de
$a$ et $b$ est lue sur le croquis, et qu'un « $\neq 0$ » sous
$H^{1}_{c}(\widetilde{Y})$ est douteux ; avec les données du croquis, ce groupe
est nul.}

\pagerange{101}{119}
\section*{X. Les notes de travail : poids et niveaux, faisceaux admissibles (pages 101 à 119)}

\pagerange{101}{103}
\subsection*{Deux feuillets d'approche (pages 101 et 103)}

La page 101 est un feuillet anglais dactylographié, isolé, qui continue le
formalisme de la page 99 sans en être la suite matérielle : fonctorialité de la
classe de cycle pour $f^{*}$ et pour $f_{*}$, $f_{*} \colon C^{i}(X) \to
C^{i-r}(Z)$ et $H^{2i}_{\ell}(X) \to H^{2i-2r}_{\ell}(Z)$ avec $r = \dim X - \dim Z$ ;
puis, par Künneth et Poincaré, $H^{*}_{\ell}(X \times Z) \simeq
\operatorname{Hom}(H^{*}_{\ell}(X), H^{*}_{\ell}(Z))$, un homomorphisme étant dit
\emph{algébrique} si la classe associée l'est. Il s'arrête sur une composition
annoncée et non écrite.

La page 103, feuillet isolé à l'encre bleue, dresse les tables qui vont être
démontrées. Pour $H^{1}$ : dans le cas propre (ou pour $H^{1}_{c}$), poids $0$ et
$1$, de niveaux $0$ et $1$ ; dans le cas lisse, poids $1$ et $2$, de niveaux $1$ et
$0$. Pour $H^{2}$ : dans le cas propre, poids $0$, $1$, $2$, de niveaux $0$, $1$,
$2$ ; dans le cas lisse, poids $2$, $3$, $4$, de niveaux $2$, $1$, $0$. Les morceaux
purs seraient semi-simples, sauf, pour $H^{2}$, celui de poids $2$, extension par
un motif de Tate $\mathbb{Q}(-1)^{r}$ ; et pour $H^{2}_{c}$, il semble qu'il puisse
y avoir « mélange de poids dans le mauvais sens ».

\pagerange{104}{108}
\subsection*{Poids de la cohomologie $\ell$-adique d'une variété singulière ou non propre (pages 104 à 108)}

Sous ce titre encadré --- avec, en haut à droite, « sur un corps de type fini »
---, la démonstration, conditionnelle aux conjectures de Weil et à la résolution
des singularités, de ce que Deligne établira en théorie de Hodge mixte et en
cohomologie $\ell$-adique.\footnote{Les énoncés qui suivent sont aujourd'hui des
théorèmes : en caractéristique nulle par la théorie de Hodge mixte (Deligne,
Hodge II et III, 1971-1974), et $\ell$-adiquement par Weil II (1980), les
altérations de de Jong (1996) remplaçant la résolution là où elle manque. Les
poids se définissent sur un corps fini par les valeurs absolues du Frobenius, et
sur un corps de type fini par spécialisation.}

\emph{A) Coefficients constants.} (1) Soit $X$ lisse et quasi-projective de
dimension $n$, écrite par résolution forte $X = X' \smallsetminus Y$ avec $X'$
projective et lisse et $Y = \bigcup Y_{i}$ un diviseur à croisements normaux
stricts. La suite exacte
\[
H^{i-1}(X') \to H^{i-1}(Y) \to H^{i}_{c}(X) \to H^{i}(X') \to H^{i}(Y) \to \cdots
\]
et la suite spectrale de Mayer--Vietoris de $Y$, dont le terme $E_{2}^{pq}$ est
calculé à partir des $H^{q}(Y_{i_{0}} \cap \cdots \cap Y_{i_{p}})$, intersections
projectives et lisses de dimension $n-p-1$, donc pures de poids $q$, donnent :

\textbf{Proposition 1 (page 104).} \emph{$H^{i}_{c}(X)$ a ses poids dans
$[0, i]$. Sa partie de poids $i$ est l'image de $H^{i}_{c}(X) \to H^{i}(X')$, sa
partie de poids $\leqslant i-1$ l'image de $H^{i-1}(Y)$. Si $i \geqslant n$, son
niveau est $\leqslant 2n-i$, ses poids sont $\geqslant 2(i-n)$, et le morceau de
poids $2(i-n) + k$ est de niveau $\leqslant k$.}

La filtration obtenue est la filtration par les poids, donc indépendante de la
compactification ; la suite spectrale dégénère pour des raisons de
poids.\footnote{La borne $2(i-n)$ vient de ce que $E_{2}^{pq}$, avec $p + q =
i-1$, n'est non nul que si $q \leqslant 2(n-p-1)$. La page note en marge que
« tous les niveaux de la filtration sont effectifs » et qu'« en tous cas, niveau
$\leqslant \inf(i, n)$ ».}

\textbf{Corollaire 1 (page 105).} \emph{Par dualité, $H^{i}(X)$ a ses poids dans
$[i, 2n]$, et sa partie de poids $i$ est exactement l'image de $H^{i}(X')$. Si
$i \geqslant n$, les morceaux de poids $i, i+1, \ldots, 2n$ sont de niveaux
respectivement $\leqslant 2n-i, 2n-i-1, \ldots, 0$ ; si $i \leqslant n$, il n'y a
que les poids $i, \ldots, 2i$, de niveaux $\leqslant i, i-1, \ldots, 0$.}

(2) Si $X$ est lisse, séparée mais non quasi-projective, une compactification de
Nagata et la résolution donnent les mêmes résultats. Si $X$ n'est pas séparée,
$H^{i}_{c}$ n'est plus défini, mais la suite spectrale d'un recouvrement affine
montre que $H^{i}(X)$ est effectif, à poids dans $[0, 2i]$ si $i \leqslant n$ et
dans $[0, 2n]$ si $i \geqslant n$, de niveau $\leqslant i$ ; la marge demande si le
niveau est $\leqslant 2n-i$.

(3) Pour $X$ séparée non lisse, $H^{i}_{c}(X)$ est encore effectif, de poids
$\leqslant i$, de niveau $\leqslant \inf(i, 2n-i)$ : par récurrence sur la dimension,
avec un ouvert lisse dense $U$ et la suite $H^{i}_{c}(U) \to H^{i}_{c}(X) \to
H^{i}_{c}(X \smallsetminus U)$.

(4) Pour $X$ quelconque, $H^{i}(X)$ est effectif, de poids $\leqslant 2\inf(i, n)$
et de niveau $\leqslant \inf(i, n)$ : on se ramène à $X$ affine, on raisonne par
récurrence sur $i$ avec une résolution $X' \to X$ et la suite spectrale de
descente des produits fibrés itérés $(X'/X)^{s}$ ; les termes $E_{2}^{pq}$ avec
$q < i$ relèvent de la récurrence, et $E_{2}^{0,i} \subset H^{i}(X')$ est borné
par le cas lisse ; si $i > n$, $H^{i}(X) = 0$, $X$ étant affine.\footnote{La
descente est licite parce qu'un morphisme propre et surjectif est de descente
cohomologique (SGA 4, exposé V bis) ; les produits fibrés itérés ne sont pas
lisses, mais la récurrence porte sur tous les schémas. La page note l'ajout
« $\leqslant i$ et $\leqslant n$ » sans qu'on voie s'il porte sur les poids ou sur
le niveau ; la lecture donnée est celle que confirme le point C) 2).}

\emph{B) Coefficients tordus.} Si les coefficients proviennent d'un motif lisse
sur $X$ lisse, on se ramène par dévissage à $F = R^{j}f_{*}\mathbb{Q}_{\ell}$ pour
$f \colon Y \to X$ projectif et lisse, qu'on peut supposer de dimension relative
$j$. D'après Deligne, la suite spectrale de Leray de $f$ dégénère, et $H^{i}(X, F)$
est un morceau de $H^{i+j}(Y)$ : effectif, à poids entre $i+j$ et $2n+2j$, de
niveaux $\leqslant 2n-i+j, \ldots, 0$ ; si $i \leqslant n$, à poids $i+j, \ldots,
2(i+j)$ et niveaux $i+j, \ldots, 0$. Dualement, $H^{i}_{c}(X, F)$ est effectif de
poids $\leqslant i+j$. Si $X$ est propre, $H^{i}(X, F)$ est pur de poids $i+j$, « et
semi-simple ici ».\footnote{La dégénérescence est le théorème de Deligne,
« Théorème de Lefschetz et critères de dégénérescence de suites spectrales »
(1968), démontré modulo le théorème de Lefschetz fort --- ce que dit le rapport
de la page 122. Trois questions de marge restent sur la page : peut-on remplacer
$2n+2j$ par $2n+j$, $2i+2j$ par $2i+j$, et les poids sont-ils $\geqslant j$ ? À la
dernière la réponse est oui, puisqu'ils sont $\geqslant i+j$.} Si $F$ n'a qu'une
représentation locale sur $X$, on passe par la cohomologie à supports propres,
qui permet de découper $X$, puis par dualité.

\emph{C) Coefficients effectifs de poids $j$ et de niveau $j$, $X$ quelconque.}
$H^{i}_{c}$ se traite comme en B). Et $H^{i}(X, F)$ est effectif, de poids
$\leqslant 2\inf(i, n) + 2j$ et de niveau $\leqslant \inf(i, n) + j$ : ce sont
les bornes que donne l'argument de B), $H^{i}(X, F)$ étant un morceau de
$H^{i+j}(Y)$ avec $\dim Y = n + j$.\footnote{La page écrit « poids $\leqslant
\inf(2i, 2n)$, niveau $\leqslant \inf(i, n)$ », sans terme en $j$, ce qui ne peut
valoir pour $j > 0$ (déjà $H^{0}(X, F)$ est de poids $j$), et s'arrête sur un
tiret et un point d'interrogation. La page 119 proposera, pour $F$ spécial de
poids $j$, « poids $\leqslant 2\inf(i,n) + j$ ? » : c'est la question de marge de
B), qu'elle laisse ouverte, et qu'on laisse ouverte.}

\pagerange{110}{114}
\subsection*{Faisceaux $\ell$-adiques admissibles et motifs virtuels (pages 110 à 114)}

La même matière, reprise en numéros --- « Le bon point de vue commence à se
voir ! », dit la marge --- pour des $\mathbb{Q}_{\ell}$-faisceaux sur des schémas
où $\ell$ est inversible, le cas intéressant étant celui des schémas localement
de type fini sur $\mathbb{Z}$. L'idée est de remplacer les motifs, qu'on n'a pas,
par la classe des faisceaux $\ell$-adiques qui « ont l'air de venir de la
géométrie ».

\emph{2.1.} Un faisceau $F$ sur $S$ est \emph{spécial de poids $i$} s'il est
isomorphe à un sous-quotient tordu d'un $R^{i}f_{*}\mathbb{Q}_{\ell}$, avec $f
\colon X \to S$ projectif et lisse ; pour $i < 0$ cela veut dire $F = 0$.

\textbf{Lemme 2.2.} \emph{a) La propriété d'être spécial de poids $i$ est stable
par sommes directes. b) Si $F$, $F'$ sont spéciaux de poids $i$, $j$, $F \otimes F'$
est spécial de poids $i+j$. c) Elle est stable par image inverse. d) Si $F$ est
spécial de poids $i$, $\check{F}(-i)$ l'est aussi. e) $\mathbb{Q}_{\ell}(-i)$ est
spécial de poids $2i$ pour $i \geqslant 0$.}

Les démonstrations sont immédiates : réunion disjointe, produit fibré et Künneth,
changement de base propre, théorème de Lefschetz fort relatif et dualité de
Poincaré, $\mathbb{Q}_{\ell}(-i) = R^{2i}f_{*}\mathbb{Q}_{\ell}$ pour $\mathbb{P}^{i}_{S}$.
D'où, en marge, le \textbf{Corollaire 2.3} : \emph{si $F$ est spécial de poids $i$
et $n \geqslant i$, $\check{F}(-n)$ est spécial de poids $2n-i$.}\footnote{Le point
d) demande le théorème de Lefschetz fort relatif en cohomologie $\ell$-adique, que
Deligne démontrera dans Weil II (1980) ; en caractéristique $p$ il était alors
conjectural. Une marge qualifie la définition de « canulée » --- lecture
incertaine --- faute de pouvoir se ramener à une situation définie sur
$\mathbb{Z}$.}

\emph{3.1-3.2.} $F$ est \emph{effectivement admissible à poids dans} $P \subset
\mathbb{Z}$ si, pour $S$ quasi-compact, il existe une partition finie de $S$ en
sous-schémas localement fermés $S_{\alpha}$ sur chacun desquels $F|S_{\alpha}$
admet une filtration dont les quotients successifs sont spéciaux de poids dans
$P$.\footnote{La définition 3.1 est barrée et récrite dans l'interligne, illisible
en partie ; on donne la forme 3.2, dont une accolade de la page dit qu'elle doit
précéder la Proposition 3.3. Une marge ajoute qu'il faudrait permettre une
extension finie radicielle locale.}

\textbf{Proposition 3.3.} \emph{a) La propriété « effectivement admissible à poids
dans $P$ » est stable par sous-faisceaux et quotients constructibles, sommes finies
et extensions. b) Si $F$, $F'$ le sont à poids dans $P$, $P'$, alors $F \otimes F'$
l'est à poids dans $P + P'$. c) Elle est stable par changement de base. d) Si $F$
est constant tordu, effectivement admissible à poids dans $P$, et si $P \leqslant
n$, alors $\check{F}(-n)$ est effectivement admissible constant tordu, à poids dans
$2n - P = \{2n - i \mid i \in P\}$.}\footnote{La page écrit $2n - P = \{n - i \mid i
\in P\}$ et un twist de $\check{F}$ surchargé ; le Corollaire 2.3 donne $2n - i$ et
la torsion $(-n)$. La formation de $\check{F}$ commute au changement de base pour
$F$ constant tordu, d'où d) par 2.2.}

\emph{4.1.} Les \emph{niveaux} s'ajoutent en exigeant, dans 3.1, que les quotients
soient de la forme $G'(-k)$, $k \geqslant 0$, avec $G'$ spécial de poids $\leqslant
\nu$ ; et « à poids dans $P$ et de niveau $\leqslant \nu$ » si de plus le poids
$j + 2k$ de $G'(-k)$ est dans $P$ --- ce qui, prévient-il, n'est pas a priori la
conjonction des deux conditions.

\textbf{Proposition 4.2.} \emph{a) « Effectivement admissible à poids dans $P$ et
de niveau $\leqslant \nu$ » est stable par sous-faisceaux et quotients
constructibles, sommes finies, extensions. b) Pour $F_{1}$, $F_{2}$ à poids dans
$P_{1}$, $P_{2}$ et niveaux $\leqslant \nu_{1}$, $\nu_{2}$, $F_{1} \otimes F_{2}$ est
à poids dans $P_{1} + P_{2}$ et de niveau $\leqslant \nu_{1} + \nu_{2}$. c) Stabilité
par changement de base. d) Si $F$ est effectivement admissible constant tordu, de
poids $\nu + 2k$ et de niveau $\leqslant \nu$, alors pour $n = \nu + k + h$, $h
\geqslant 0$, $\check{F}(-n)$ l'est aussi, de poids $2n - \nu - 2k = \nu + 2h$ et de
niveau $\leqslant \nu$.}\footnote{La page écrit $\check{F}(n)$, lecture probable
d'un twist surchargé ; avec sa convention, où $F(n)$ est la torsion par
$\mathbb{Q}_{\ell}(1)^{\otimes n}$ (page 110), le poids $2n - \nu - 2k$ qu'elle
annonce est celui de $\check{F}(-n)$.}

\textbf{Théorème 5.1.} \emph{Soit $f \colon X \to S$ séparé de présentation finie,
$n$ le maximum des dimensions des fibres, et $F$ strictement spécial de poids $j$
sur $X$, c'est-à-dire facteur direct d'un $R^{j}g_{*}\mathbb{Q}_{\ell}$ avec $g
\colon Y \to X$ projectif et lisse. Alors $R^{i}f_{!}F$ est effectivement
admissible, à poids $\leqslant i+j$ et à niveaux $\leqslant \inf(i+j, 2n-i+j)$. Si
$i \geqslant n$, ses poids sont compris entre $2(i-n)$ et $i+j$, le morceau de poids
$2(i-n) + k$ étant de niveau $\leqslant k$, jusqu'au niveau $2n-i+j$ en poids $i+j$.}

La borne inférieure $2(i-n)$ ne dépend pas de $j$, et c'est juste : $R^{i}f_{!}F$
est un morceau de $R^{i+j}(fg)_{!}\mathbb{Q}_{\ell}$, dont les fibres sont de
dimension $n + j$, et $2((i+j) - (n+j)) = 2(i-n)$. La démonstration traite d'abord
$F = R^{j}g_{*}\mathbb{Q}_{\ell}$ : par changement de base propre et dévissage, on se
ramène à $S$ affine intègre, et, après un changement de base fini radiciel, à $X$
affine lisse sur $S$ compactifié en un $\widehat{X}$ projectif et lisse dont le bord
est un diviseur relatif à croisements normaux ; puis à $F = \mathbb{Q}_{\ell}$, qui
est le point A) (1) des pages 104 et 105 --- sans avoir besoin, note-t-il, de la
dégénérescence qui y intervient. Pour $F$ quelconque, on se ramène à $j$ égal à la
dimension relative de $g$, par Lefschetz ou en remplaçant $Y$ par $Y \times
\mathbb{P}^{j-d}$ ; la suite spectrale de Leray de $fg$ dégénère par Deligne, et
$R^{i}f_{!}F$ est un facteur direct de $R^{i+j}(fg)_{!}\mathbb{Q}_{\ell}$, auquel on
applique le premier cas.\footnote{La page dit « quotient d'un sous-faisceau » ;
le théorème de décomposition de Deligne (1968) donne mieux, $Rg_{*}\mathbb{Q}_{\ell}
\simeq \bigoplus_{q} R^{q}g_{*}\mathbb{Q}_{\ell}[-q]$ pour $g$ projectif et lisse,
d'où un facteur direct. Les marges ajoutent les hypothèses dont la démonstration a
besoin : Lefschetz fort « à la Deligne », résolution des singularités sur les corps
résiduels de $S$, et peut-être « semi-simplicité arithmétique ».}

\textbf{Corollaire 5.2.} \emph{Si la partie « semi-simplicité » des conjectures de
Tate est vraie, la conclusion de 5.1 vaut pour $F$ seulement effectivement
admissible de poids $\leqslant j$ : $R^{i}f_{!}F$ est effectivement admissible, de
poids $\leqslant i+j$, de niveau $\leqslant \inf(i+j, 2n-i+j)$, et, si $i \geqslant n$,
de poids entre $2(i-n)$ et $i+j$ avec les niveaux $0, 1, \ldots, 2n-i+j$.} Par
semi-simplicité, « effectivement admissible » devient « strictement spécial ».

\textbf{Corollaire 5.3.} \emph{Sans conjecture de Tate : si $F$ est effectivement
admissible, $R^{i}f_{!}F$ l'est.} On se ramène, par dévissage et changement de base
propre, au cas où $F$ est spécial et $S$ de type fini sur $\mathbb{Z}$.

\textbf{Corollaire 5.4.} \emph{Sur un corps $k$, si $X$ est lisse, séparé, de type
fini, et $F$ strictement spécial de poids $j$ --- ou effectivement admissible de
poids $j$ si l'on admet la semi-simplicité ---, alors $R^{i}f_{*}F$ est
effectivement admissible, à poids entre $i+j$ et $2n+2j$ et à niveaux $\leqslant
2n-i+j, \ldots, 0$ ; si $i \leqslant n$, à poids $i+j, \ldots, 2(i+j)$ seulement, de
niveaux $i+j, \ldots, 0$.} On peut le démontrer sans dualité, en reprenant la
démonstration de 5.1 par Leray et en se ramenant à $F = \mathbb{Q}_{\ell}$ et à une
compactification lisse à bord à croisements normaux --- « à titre de vérification
de cohérence du yoga ! ».

Le numéro 6 ouvre le cas de $R^{i}f_{*}F$ pour $f$ de type fini et $F$
effectivement admissible de poids $j$, en commençant par $F = \mathbb{Q}_{\ell}$ et
le cas d'une immersion ; la page 114 s'arrête sur ce mot.

\pagerange{115}{119}
\subsection*{Le cas d'une immersion : les feuilles de calcul (pages 115, 117 et 119)}

Trois feuilles de calcul, sans phrases suivies, qui continuent le numéro 6 sous
la forme d'un point « (1) », que le « (2°) » de la page 119 et le « 2) » de la page
121 poursuivent. Soit $X$ régulier, $F = \mathbb{Q}_{\ell}$, $f \colon X \to Y$ une
immersion, $n = \dim \overline{f(X)}$. On factorise $f = h \circ g$, avec $g \colon X
\to Y'$ immersion ouverte dans $Y'$ régulier, de complémentaire un diviseur à
croisements normaux $Z = \sum Z_{i}$, et $h \colon Y' \to Y$ propre. La suite
spectrale $R^{p}h_{*}(R^{q}g_{*}\mathbb{Q}_{\ell}) \Rightarrow R^{p+q}f_{*}
\mathbb{Q}_{\ell}$, avec
\[
R^{q}g_{*}\mathbb{Q}_{\ell} = \bigoplus_{i_{1} < \cdots < i_{q}}
\mathbb{Q}_{\ell, Z_{i_{1}} \cap \cdots \cap Z_{i_{q}}}(-q) ,
\]
pur de poids $2q$ par la pureté,\footnote{Sur une base régulière quelconque,
c'est la conjecture de pureté absolue de Grothendieck, démontrée par Gabber
(publiée par Fujiwara, 2002) ; la page l'emploie sans la nommer.} donne des
termes $E_{2}^{pq}$ effectivement admissibles de poids $\leqslant p + 2q$. Les
tables des pages 117 et 119 --- poids et niveaux de chaque $E_{2}^{p,i-p}$, selon
que $i \leqslant n-1$ ou $i \geqslant n$ --- aboutissent à l'encadré
\[
\text{poids} \leqslant 2\inf(i, n) , \qquad \text{niveau} \leqslant \inf(i, n-1)
\]
pour $R^{i}f_{*}\mathbb{Q}_{\ell}$, les niveaux croissant de $0$ à leur maximum puis
décroissant jusqu'à $0$ en poids $2n$, lequel n'intervient que par
$E_{2}^{2(i-n), 2n-i}$, faisceau porté par des intersections de $2n-i$ composantes,
de dimension $i-n$.\footnote{Les niveaux des dernières colonnes des tables sont
écrits vite et en partie surchargés ; la transcription les donne sans les
vérifier, et on ne les reproduit pas colonne par colonne. L'encadré, lui, est net.}
Puis, avec un « ? » : pour $F$ strictement spécial de poids $j$, poids $\leqslant
2\inf(i,n) + j$ et niveau $\leqslant \inf(i, n-1) + j$ ; et, dans le cas relatif
« (2°) », au voisinage d'un point $y$ de $Y$, avec $n = d_{y}(f)$ défini par une
borne sur les degrés de transcendance, poids $\leqslant 2\inf(i, n) + j$ et niveau
$\leqslant \inf(i, n) + j$, encore suivis d'un « ? ». La marge demande si l'on peut
remplacer $n$ par la dimension des fibres.

\pagerange{120}{122}
\section*{XI. Le cas relatif, et un rapport sur les travaux de Deligne (pages 120 à 122)}

\pagerange{121}{121}
\subsection*{Le point 2) : coefficients strictement spéciaux (page 121)}

La page 121, de sa main, est le point « 2) » de la suite des pages 115 à 119 :
$F$ strictement spécial, $X$ régulier, $F = R^{d}\varphi_{*}\mathbb{Q}_{\ell}$ pour un
morphisme $\varphi \colon X' \to X$ de dimension relative $d$, inscrit dans un carré
$X' \to Y'$ au-dessus de $f \colon X \to Y$. La suite spectrale de Leray de
$f\varphi$ dégénère par Deligne, de sorte que $R^{i}f_{*}F$ est un morceau de
$R^{i+d}(f\varphi)_{*}\mathbb{Q}_{\ell}$ ; l'autre suite spectrale, par $Y'$ de
dimension $\leqslant n + d$, borne les poids et les niveaux, la fin de la page ne se
lisant que par fragments.\footnote{La page écrit tantôt $R^{d}\varphi_{*}$, tantôt
$R^{j}\varphi_{*}$ pour le même faisceau ; on écrit $d$. Le mot qui qualifie
$R^{i}f_{*}F$ se termine peut-être en « -quotient ».}

\pagerange{122}{122}
\subsection*{Deux feuillets d'un rapport, lus dans l'ordre 122 puis 120}

Deux feuillets dactylographiés en français, corrigés à la main, présentent des
travaux de Deligne en points numérotés ; le feuillet 122 porte le point 4, le
feuillet 120 le point 5. Ils forment un seul texte : la page 122 s'achève, après
les mots « théorie des motifs », sur « qui se », et la page 120 commence par « une sorte de
synthèse géométrico-arithmétique des nombreuses théories cohomologiques dont on
dispose à présent pour les variétés algébriques, et des représentations
linéaires de groupes de Galois arithmétiques auxquelles ces théories donnent
naissance ».\footnote{Le raccord est de nous ; la transcription donne chaque
feuillet comme isolé. Le dernier mot de la page 122 est surchargé à la main, de
sorte que la lecture « qui se[ra] une sorte de synthèse » reste une conjecture
de lecture, mais le sens et la syntaxe s'enchaînent sans reste.}

\emph{Le point 4.} Un morphisme projectif et lisse de schémas est
\emph{cohomologiquement trivial} : la suite spectrale de Leray en cohomologie
$\ell$-adique --- ou réelle sur $\mathbb{C}$ --- dégénère. Le théorème est prouvé par
Deligne modulo le théorème de Lefschetz pour la cohomologie des fibres, et la
démonstration vaut aussi pour les fibrations analytiques à fibres kählériennes
compactes. Le cas où le système local de la cohomologie des fibres est trivial
était le résultat principal de la thèse de Blanchard, hypothèse qui n'est réalisée
dans presque aucun cas intéressant ; la démonstration de Deligne est différente,
et il ne connaissait pas celle de Blanchard. L'intérêt, dans l'optique du
rapporteur, est surtout en géométrie algébrique abstraite et en théorie des
nombres ; ce sont d'ailleurs des considérations arithmétiques qui l'avaient amené
à communiquer à Deligne l'énoncé comme une conjecture. Il en tire que les
hypothèses de Riemann--Weil impliquent leurs analogues à coefficients tordus, et
que le théorème sera l'un des ingrédients indispensables de la « théorie des
motifs ».\footnote{C'est l'article de Deligne de 1968 cité plus haut ; les pages
107, 113, 121 s'en servent, sous le nom de Deligne. Le feuillet remplace une fois
« rédacteur » par « rapporteur », pas l'autre.}

\pagerange{120}{120}
\subsection*{Le point 5 (page 120)}

\emph{Le point 5} est la solution, par Deligne, du problème de l'introduction d'une
$\lambda$-structure --- les puissances extérieures --- dans l'anneau $K^{\bullet}(X)$
construit sur les complexes parfaits, pour un espace localement annelé ; il était
signalé comme le plus important problème resté en suspens dans le séminaire de
l'IHES de 1966-1967 sur le théorème de Riemann--Roch (SGA 6). La nécessité de
travailler avec des complexes parfaits plutôt qu'avec des fibrés vectoriels est
imposée « par des impératifs techniques irrécusables ». La solution passe par des
foncteurs puissances extérieures dans la catégorie dérivée, construits par la
théorie de Dold--Puppe sur les complexes de chaînes ; la difficulté sérieuse était
de prouver les identités universelles des $\lambda$-anneaux, et Deligne y parvient
en déterminant d'abord le $\lambda$-anneau de toutes les opérations tensorielles
universelles sur les fibrés vectoriels. Il n'a écrit sa théorie qu'au-dessus d'un
corps de base ; la page s'arrête là.\footnote{Le contenu place le rapport après
SGA 6 et après le théorème de dégénérescence ; la transcription enregistre que des
feuillets voisins portent au verso des dates de mars 1968, que cette lecture ne
reprend pas comme datation.}

\pagerange{124}{136}
\section*{XII. Divers types de $\mathbb{Q}_{\ell}$-faisceaux (pages 124 à 136)}

Un brouillon rapide sur les rectos d'un exposé Bourbaki, dont les formules se
lisent et la prose souvent pas. Il reprend la notion de faisceau admissible des
pages 110 à 114 en ne gardant qu'une donnée, la dimension des fibres.\footnote{Le
premier mot du titre, « Divers », est une lecture douteuse.}

\emph{Définition.} Un $\mathbb{Q}_{\ell}$-faisceau constructible $F$ sur $X$ est
\emph{admissible de type} $\leqslant d$ si, localement --- et, comme le dit le
crochet qui suit, en vérifiant fibre par fibre, c'est-à-dire sur les strates d'une
stratification ---, il admet une filtration finie dont les constituants sont des
sous-quotients de sommes directes de faisceaux $R^{n}g_{*}\mathbb{Q}_{\ell}$, avec $g
\colon V \to U$ projectif et lisse à fibres connexes de dimension $\leqslant
d$.\footnote{La page écrit « $\operatorname{gr}^{i}$ isomorphes à des sommes
directes de faisceaux de la forme $R^{n}g_{*}(\mathbb{Q}_{\ell})$ », avec en marge
« en.\ constituants ». Les deux précisions retenues --- sous-quotients, strates ---
sont nécessaires : sans la première le Lemme 1 est faux, sans la seconde un
faisceau porté par un fermé ne serait jamais admissible.}

\textbf{Lemme 1.} \emph{Les faisceaux admissibles de type $\leqslant d$ forment une
sous-catégorie abélienne épaisse --- stable par sous-objets, quotients et
extensions --- des $\mathbb{Q}_{\ell}$-faisceaux constructibles.}

\textbf{Lemme 2.} \emph{Ils sont stables par image inverse} --- « trivial ».

\textbf{Lemme 3.} \emph{Soit $f \colon X \to Y$ un morphisme compactifiable de type
fini, de dimension relative $\leqslant d'$, et $F$ admissible de type $\leqslant d$
sur $X$. Alors les $R^{i}f_{!}F$ sont admissibles de type $\leqslant d + d'$.}

La démonstration, aux pages 126 à 132, se fait par récurrence sur $d + d'$, $Y$
étant un corps. Un ouvert $U$ lisse et dense où $F$ est filtré, et son
complémentaire, ramènent par la suite exacte des supports propres et les lemmes
1 et 2 au cas où $X$ est lisse et $F$ un constituant d'un $R^{m}g_{*}\mathbb{Q}_{\ell}$,
puis, par deux suites exactes courtes, au cas $F = R^{m}g_{*}\mathbb{Q}_{\ell}$ et aux
$R^{p}f_{!}(R^{q}g_{*}\mathbb{Q}_{\ell})$. Une \emph{factorisation d'Artin} de $U$,
licite parce qu'on peut prendre $U$ petit, ramène au cas où $U$ est de dimension
$1$ ;\footnote{Ce sont les « bons voisinages » d'Artin, SGA 4, exposé XI : une
fibration en courbes lisses élémentaires.} alors $R^{0}f_{!}(R^{q}g_{*}) = 0$ et la
suite spectrale se réduit à des suites exactes courtes
\[
0 \to R^{2}f_{!}(R^{n-2}g_{*}\mathbb{Q}_{\ell}) \to R^{n}(fg)_{!}\mathbb{Q}_{\ell} \to
R^{1}f_{!}(R^{n-1}g_{*}\mathbb{Q}_{\ell}) \to 0 ,
\]
de sorte qu'il suffit de voir que $R^{n}h_{!}\mathbb{Q}_{\ell}$, $h = fg$, est
admissible ; la résolution des singularités $V = \widehat{V} \smallsetminus Z$ et la
suite $R^{n-1}h_{Z!}\mathbb{Q}_{\ell} \to R^{n}h_{!}\mathbb{Q}_{\ell} \to
R^{n}\hat{h}_{*}\mathbb{Q}_{\ell}$ concluent, par définition pour le troisième terme et
par récurrence pour le premier.\footnote{Le haut de la page 132, qui achève ce
raisonnement, ne se lit que par mots isolés (« O.K. par définition », « par hyp.\ de récurrence
. O.K. ») ; un cas « b) $m = 0$ » par récurrence descendante est
barré page 128.}

\textbf{Proposition 1 (page 132).} \emph{Soit $i \colon U \to X$ une immersion
ouverte, $X$ excellent de dimension $\leqslant d$, et $F$ admissible de type
$\leqslant d'$ sur $U$. Alors les $R^{n}i_{*}F$ sont admissibles de type $\leqslant
d + d'$.}\footnote{La page écrit « $(d+d'-1)$-admissible », le dernier chiffre biffé
et surchargé, lecture douteuse. Avec $-1$, l'énoncé est faux déjà pour $F =
\mathbb{Q}_{\ell}$ ($d' = 0$) : si le bord de $U$ est un diviseur à croisements
normaux dont $d$ composantes se coupent en un point, $R^{d}i_{*}\mathbb{Q}_{\ell}$
contient en ce point $\mathbb{Q}_{\ell}(-d)$, qui est de poids $2d$ et ne peut être
un constituant d'un $R^{m}g_{*}$ à fibres de dimension $< d$. La borne $d + d'$ est
celle que la démonstration peut donner ; elle n'est pas établie sur la page.}

La démonstration commence par récurrence sur $d + d'$ et par récurrence
noethérienne : un ouvert régulier dense $V \subset U$ et le triangle $F \to
Rj_{*}F \to Q$, avec $Q$ porté par un fermé strict, ramènent au cas où $U$ est
régulier ; la résolution factorise $i$ en $i'$ suivi d'un $f$ projectif, avec $X'$
régulier et $X' \smallsetminus U$ à croisements normaux, et le Lemme 3 ramène à ce
cas. Pour $F = \mathbb{Q}_{\ell}$, on conclut par la pureté. Pour le cas général, on
compactifie $V$ par un morphisme projectif $g$, on écrit $F$ comme facteur direct
d'un $R^{n}t_{*}\mathbb{Q}_{\ell}$, on se ramène à $d' = 1$, et la suite spectrale de
la page 136 a ses termes $E_{2}^{pq}$ nuls hors de $0 \leqslant q \leqslant 2$ : les
$E_{2}^{p0}$ sont $(d-1)$-admissibles par récurrence, $E_{2}^{p2} \simeq
\mathbb{Q}_{\ell}(-1) \otimes E_{2}^{p0}$ aussi, et $E_{2}^{p1}$ est une extension de
trois faisceaux dont le dernier, $E_{\infty}^{p1}$, ramène au cas de
$R^{n}j_{*}\mathbb{Q}_{\ell}$, « qui reste à traiter…!! ». La Proposition 1 n'est donc
pas démontrée dans le dossier : il y manque exactement ce cas.

\pagerange{139}{141}
\section*{XIII. Motifs lisses et isomotifs lisses (pages 138 à 141)}

Une chemise encadrée, page 138, porte ce titre ; suivent deux pages au net et une
page de questions.

Soit $S$ normal et intègre, de point générique $\eta$, $K = k(\eta)$. Un
\emph{isomotif lisse} sur $S$ est un isomotif semi-simple $M$ sur $K$ tel que, pour
tout $\ell \neq \operatorname{car} K$, $T_{\ell}(M)$ soit non ramifié sur $S_{\ell} =
S[1/\ell]$ ; conjecturalement il suffit de le demander pour un seul $\ell$ premier à
toutes les caractéristiques résiduelles, ou pour deux $\ell \neq
\ell'$.\footnote{La page écrit « isomotif semi-simple sur $S$ » ; l'isomotif est pris
sur $K$, et c'est la condition de non-ramification qui le rattache à $S$. Pour
les variétés abéliennes sur une base de dimension $1$, l'indépendance de $\ell$
est le critère de Néron--Ogg--Šafarevič ; en général elle est conjecturale.}

Un \emph{motif lisse} sur $S$, pour $\operatorname{car} K = 0$, est la donnée
$(M, \{E_{\ell}\}, \{\alpha_{\ell}\})$ d'un isomotif lisse $M$, pour chaque $\ell$ d'un
$\mathbb{Z}_{\ell}$-faisceau lisse $E_{\ell}$ sur $S_{\ell}$ --- ou, ce qui revient au
même, sur $K$ et non ramifié sur $S_{\ell}$ --- et d'un isomorphisme
$\alpha_{\ell} \colon E_{\ell} \otimes \mathbb{Q}_{\ell} \xrightarrow{\sim} T_{\ell}(M)$,
avec l'axiome que pour presque tout $\ell$, $\alpha_{\ell}$ provient d'un isomorphisme
sur $\mathbb{Z}_{\ell}$ ; en particulier $E_{\ell}$ est sans torsion pour presque tout
$\ell$.\footnote{La page écrit $u_{\ell}$ dans l'axiome, là où les données ont
$\alpha_{\ell}$. L'axiome sous-entend une structure entière de référence sur $M$ ;
la page n'en dit pas plus.} Un morphisme est un couple $(u, \{u_{\ell}\})$ d'un
morphisme d'isomotifs et de morphismes de faisceaux compatibles aux
$\alpha_{\ell}$. C'est, pour les motifs, ce que sont les variétés abéliennes par
rapport à leur catégorie à isogénie près : une structure entière.

Le foncteur des motifs lisses sur $S$ vers les motifs sur $K$ est pleinement fidèle,
les $\operatorname{Hom}$ sont des $\mathbb{Z}$-modules de type fini, la catégorie est
abélienne. Le foncteur vers les isomotifs lisses est exact et essentiellement
surjectif ; les groupes commutatifs finis sur $K$ dont les composantes
$\ell$-primaires sont non ramifiées sur $S[1/\ell]$ forment une sous-catégorie épaisse,
et le quotient est la catégorie des isomotifs lisses. La catégorie a un produit
tensoriel et des $\underline{\operatorname{Hom}}$, compatibles aux réalisations. Un
schéma projectif et lisse sur $S$ définit des motifs $R^{i}f_{*}(\mathbb{Z}_{X})$.

\emph{Schémas en groupes.} On définit un foncteur contravariant des schémas en
groupes commutatifs propres et lisses $A$ sur $S$ vers les motifs lisses :
\[
\mathfrak{P}(A) = \bigl( R^{1}f_{A^{0}*}\mathbb{Q} ,\;
\underline{\operatorname{Hom}}({}_{\ell^{\infty}}A, \mathbb{Q}_{\ell}/\mathbb{Z}_{\ell}) ,\;
u \bigr) ,
\]
où $E_{\ell} = \underline{\operatorname{Hom}}({}_{\ell^{\infty}}A,
\mathbb{Q}_{\ell}/\mathbb{Z}_{\ell})$ est une extension de
$\underline{\operatorname{Hom}}(A/A^{0}, \mathbb{Q}_{\ell}/\mathbb{Z}_{\ell})$ par
$R^{1}f_{A^{0}*}\mathbb{Z}_{\ell}$ --- le dual du module de Tate de la partie
abélienne --- et $u_{\ell}$ la flèche induite. La page 141 affirme qu'un foncteur
noté d'une lettre qu'on lit $\mathcal{L}$, sans doute celui-ci, est pleinement
fidèle, en se ramenant au cas $S = \operatorname{Spec}K$. Son image
essentielle demande que le motif soit effectif de poids $1$, donc défini par une
variété abélienne $A_{K}$, et que $A_{K}$ ait bonne réduction sur $S$ : la question
devient de savoir si la non-ramification de tous les $T_{\ell}(A)$ sur les $S_{\ell}$
suffit à la bonne réduction. Oui si $S$ est de dimension $1$ --- critère de
Néron--Ogg--Šafarevič ---, et, dit-il, si $S$ est de caractéristique nulle ; peut-être
pas pour $S$ local régulier de dimension $2$ d'inégales caractéristiques. Pour un
contre-exemple, il suffirait d'un schéma abélien $A$ sur $S$ et d'un sous-groupe
fini de $A_{K}$ qui ne se prolonge pas en un sous-groupe fini et plat de
$A$.\footnote{Le critère de Néron--Ogg--Šafarevič est dû à Serre et Tate (1968). Au-delà
de la dimension $1$, c'est une question de pureté pour les schémas abéliens, que
la page laisse ouverte et qu'on ne tranche pas ici. Le nom lu entre parenthèses
sur la page, « Ogg-Šafarevič », est une lecture très douteuse.}

Reste un « Pb » : avec cette définition, qui fait d'un motif lisse un motif sur $K$
avec une propriété de non-ramification, on ne voit pas comment définir l'image
inverse par un morphisme $S' \to S$ --- par exemple la fibre en un point de $S$ ---,
difficulté qui vaut déjà pour les isomotifs effectifs de poids $1$. Rien ne la
résout dans le dossier.

\pagerange{143}{166}
\section*{XIV. Constructions abstraites de catégories de motifs (pages 142 à 166)}

Sous une chemise de sa main, page 142 --- « Enveloppes pseudo-abéliennes de
catégories multiplicatives (Constructions abstraites de catégories de motifs) » ---,
la troisième altitude du programme : on ne suppose plus la catégorie des motifs, on
la construit, et l'on reprend dans ce cadre la graduation, l'objet de Tate, le
niveau, la dualité et la positivité.\footnote{Dans cette section, l'objet de Tate
effectif de degré $p$ est noté $\Lambda(-1)$, comme $\mathbb{Q}(-1)$ ailleurs dans le
dossier ; voir la note de la page 152.}

\pagerange{143}{148}
\subsection*{L'enveloppe pseudo-abélienne (pages 143 à 148)}

On part d'une catégorie additive $\mathcal{V}$ et l'on cherche un foncteur additif
universel $\underline{h} \colon \mathcal{V}^{\circ} \to \mathcal{M}$ vers une
catégorie abélienne, tel que pour toute catégorie abélienne $\mathcal{C}$
\[
\operatorname{Fonc}_{\mathrm{ex}}(\mathcal{M}, \mathcal{C}) \longrightarrow
\operatorname{Fonc}_{\mathrm{add}}(\mathcal{V}^{\circ}, \mathcal{C}) ,
\qquad T \longmapsto T \circ \underline{h} ,
\]
soit une équivalence, et l'on voudrait des conditions pour que $\underline{h}$ soit
pleinement fidèle. On suppose $\mathcal{V}$ $k$-linéaire --- $\mathbb{Q}$-linéaire en
géométrie --- et chaque $\operatorname{Fl}(X,X)$ une $k$-algèbre semi-simple de
dimension finie. Soit $\mathcal{M}^{\#}$ la catégorie des couples $(X, \pi)$ d'un
objet et d'un projecteur $\pi = \pi^{2}$, les morphismes $(X, \pi) \to (X', \pi')$
étant les $f$ avec $\pi' f = f\pi$ modulo ceux avec $f\pi = \pi' f = 0$ --- autrement
dit $\pi' \operatorname{Fl}(X, X') \pi$. C'est l'\emph{enveloppe pseudo-abélienne}, ou
de Karoubi, de $\mathcal{V}$ ; $X \mapsto (X, \mathrm{id})$ est pleinement fidèle,
et pour toute catégorie additive $\mathcal{C}$ où les projecteurs ont une image,
$T \mapsto T \circ \underline{h}^{\#}$ est une équivalence
$\operatorname{Fonc}_{\mathrm{add}}(\mathcal{M}^{\#}, \mathcal{C}) \simeq
\operatorname{Fonc}_{\mathrm{add}}(\mathcal{V}, \mathcal{C})$, de quasi-inverse
$\psi(H)(X, \pi) = \operatorname{Im} H(\pi)$. Jusque-là rien n'utilise les hypothèses
sur $\mathcal{V}$.

Ce sont elles qui rendent $\mathcal{M}$ \emph{abélienne semi-simple} ; alors les
foncteurs additifs de $\mathcal{M}$ dans une catégorie abélienne sont automatiquement
exacts, et le problème universel est résolu. Cela revient à deux choses : (i) tout
objet vérifie les conditions de chaînes ascendante et descendante sur les
sous-objets, ce qui vient de la dimension finie des $\operatorname{Fl}(X,X)$ ; (ii)
tout morphisme $u \colon P \to Q$ est isomorphe à un morphisme de la forme
\[
A \oplus B \longrightarrow B \oplus C , \qquad (a, b) \longmapsto (b, 0) ,
\]
de sorte qu'un monomorphisme est l'injection d'un facteur direct.\footnote{La page
énonce (ii) pour les monomorphismes, avec la formule $A \times B \to B \times C$,
$\operatorname{inj}_{1} \circ \operatorname{pr}_{2}$, qui n'est pas un monomorphisme
si $A \neq 0$ : c'est la forme normale d'un morphisme quelconque, noyau $A$, image
$B$, conoyau $C$, et c'est ainsi qu'on l'énonce.} Pour (ii), il se ramène à $P = Q$
en considérant $v = \operatorname{inj}_{2}\,u\,\operatorname{pr}_{1}$ sur $P \oplus
Q$, et l'argument est celui des anneaux semi-simples : pour tout $u$ d'un anneau
semi-simple $U = \operatorname{Fl}(P,P)$ il existe $v$ tel que $vu = \pi$ et $uv =
\varpi$ soient des idempotents, avec $u\pi = u$, $\varpi u = u$, $\pi v = v$, $v\varpi
= v$, d'où $P \simeq \pi P \oplus (1-\pi)P$, $Q \simeq \varpi Q \oplus (1-\varpi)Q$
et $u$ un isomorphisme de $\pi P$ sur $\varpi Q$.\footnote{C'est ce qu'on appelle
aujourd'hui le lemme de Jannsen (1992) : une catégorie pseudo-abélienne
$\mathbb{Q}$-linéaire dont les algèbres d'endomorphismes sont semi-simples de
dimension finie est abélienne semi-simple. Jannsen en a tiré la semi-simplicité
inconditionnelle des motifs numériques. La page écrit $\operatorname{pr}_{2}$ dans
la réduction, là où il faut projeter sur $P$.}

Un foncteur biadditif $\mathcal{V} \times \mathcal{V}' \to \mathcal{V}''$ se prolonge
de façon essentiellement unique en un foncteur biadditif sur les enveloppes ; un
produit $\boxtimes$ symétrique, associatif et unitaire sur $\mathcal{V}$ donne ainsi
un produit tensoriel $\otimes$ sur $\mathcal{M}$, et, en caractéristique nulle,
toutes les opérations tensorielles, $\bigotimes^{n}$, $\bigwedge^{n}$,
$S^{n}$ --- images des idempotents de $\mathbb{Q}[\mathfrak{S}_{n}]$.

\pagerange{149}{151}
\subsection*{L'exemple des variétés (pages 149 à 151)}

Pour $H = T \circ \underline{h}$ avec $T$ exact, $T$ est fidèle si et seulement si
$T(M) = 0$ entraîne $M = 0$, si et seulement si $T$ est conservatif, si et seulement
si $T$ est injectif sur les $\operatorname{End}(M)$.\footnote{La phrase qui devait
dire à quelle condition sur $H$ il en est ainsi s'arrête inachevée, soulignée,
au-dessus d'un trait.}

\emph{Exemple.} Soit $k$ algébriquement clos, $\mathcal{V}$ la catégorie ayant pour
objets les schémas projectifs et lisses sur $k$, avec
\[
\operatorname{Fl}(X, Y) = \bigoplus_{i,j} C^{m_{j}}(X_{i} \times Y_{j}) ,
\]
$X_{i}$, $Y_{j}$ les composantes connexes, $m_{j} = \dim Y_{j}$, $C^{q}$ les classes
de cycles de codimension $q$ à coefficients rationnels modulo l'équivalence
homologique $\ell$-adique, et la composition des correspondances. C'est la catégorie
des correspondances de degré $0$ ; l'objet nul est $\varnothing$, la somme la
réunion disjointe.\footnote{Le dossier 11, pages 2 à 18, construit la même catégorie
des correspondances à partir d'une théorie bivariante quelconque.} Si l'on admet que
l'équivalence homologique coïncide avec l'équivalence numérique, les
$\operatorname{Fl}(X,Y)$ sont de dimension finie sur $\mathbb{Q}$ ; si l'on admet
l'inégalité de Hodge, les $\operatorname{Fl}(X,X)$ sont des algèbres à involution
positive « du type de Castelnuovo--Weil », donc semi-simples. Sous ces deux
conjectures standard --- $D$ et la conjecture de type Hodge --- la construction
précédente donne une catégorie abélienne semi-simple $\mathcal{M}$ de motifs, avec
le foncteur contravariant universel $\underline{h} \colon \mathcal{V}^{\circ} \to
\mathcal{M}$.\footnote{La semi-simplicité, on le sait depuis Jannsen (1992), vaut
sans conjecture pour l'équivalence numérique ; c'est l'égalité avec l'équivalence
homologique qui reste conjecturale, et avec elle l'existence du foncteur de Tate
ci-dessous sur les motifs numériques.} La cohomologie $\ell$-adique est fidèle sur
$\mathcal{V}$ par définition même de l'équivalence homologique et par la dualité ; elle
définit un foncteur exact, fidèle, conservatif
\[
T_{\ell} \colon \mathcal{M} \longrightarrow \operatorname{Mod}_{\mathrm{tf}}(\mathbb{Q}_{\ell}) ,
\]
qu'il appelle le \emph{foncteur de Tate}. S'y ajoutent les foncteurs de Hodge et de
De Rham, fidèles en caractéristique nulle, et de Betti si $k = \mathbb{C}$. Le produit
$X \times Y$ donne la structure multiplicative, et la formule de Künneth fait de
$T_{\ell}$ un foncteur tensoriel.\footnote{Avec une réserve que la page ne fait pas :
pour que $T_{\ell}$ respecte aussi la contrainte de commutativité, il faut modifier
celle-ci sur les motifs par le signe de Koszul des degrés, ce qui demande
l'algébricité des projecteurs de Künneth (conjecture standard $C$) ; voir
Deligne--Milne (1982) et, dans ce même groupe de dossiers, le dossier 10, pages 117
à 128.}

\pagerange{151}{153}
\subsection*{Graduation et objet de Tate (pages 151 à 153)}

Revenons au cas abstrait, et donnons-nous des projecteurs $\pi_{i} \in
\operatorname{End}(\mathrm{id}_{\mathcal{V}})$, orthogonaux deux à deux, de somme
$1$ --- « somme infinie ! », mais presque tous nuls sur chaque objet. Ils s'étendent à
$\mathcal{M}$, qui devient graduée : $\mathcal{M} = \prod_{i} \mathcal{M}^{i}$, les
objets de $\mathcal{M}^{i}$ étant les motifs effectifs de poids $i$. Si de plus
$\pi_{n}(X \boxtimes Y) = \sum_{i+j=n} \pi_{i}(X) \boxtimes \pi_{j}(Y)$, le produit
envoie $\mathcal{M}^{i} \times \mathcal{M}^{j}$ dans $\mathcal{M}^{i+j}$, et l'anneau
de Grothendieck de $\mathcal{M}$ est un $\lambda$-anneau gradué ; l'unité
$\Lambda(0)$ est de degré $0$.\footnote{Dans l'exemple géométrique, les $\pi_{i}$
sont les projecteurs de Künneth, et leur existence comme correspondances
algébriques est la conjecture standard $C$. Deux lignes barrées supposaient de plus
$\Lambda(0)$ seul objet simple de degré $0$, d'anneau $\mathbb{Q}$.}

On fixe ensuite un objet $\Lambda(-1)$ de degré $p \geqslant 1$ tel que $M \mapsto
M(-1) = M \otimes \Lambda(-1)$ soit pleinement fidèle, et l'on pose $\Lambda(-n) =
\Lambda(-1)^{\otimes n}$, avec $\operatorname{End}\Lambda(-n) =
\operatorname{End}\Lambda(0) = \mathbb{Q}$. La catégorie $\widetilde{\mathcal{M}}$ des
motifs non nécessairement effectifs est la pseudo-limite inductive des copies de
$\mathcal{M}$ par $M \mapsto M(-1)$ : ses objets sont les $M(i)$, $i \in \mathbb{Z}$,
et $\operatorname{Hom}(M(i), N(j)) = \operatorname{Hom}(M(i-r), N(j-r))$ pour $r$
assez grand. Elle est graduée, et périodique de période $p$ ; sa composante de
degré $i$ est $\varinjlim_{k}\mathcal{M}^{i+kp}$.\footnote{Les pages 152 à 154 notent
cet objet $\Lambda(1)$, de degré $p$, et $M(1) = M \otimes \Lambda(1)$, alors que les
pages 155 à 166 --- et tout le reste du dossier --- prennent l'objet effectif de
poids $p = 2$ sous le nom $\Lambda(-1)$ ou $\mathbb{Q}(-1)$, et écrivent les formes
à valeurs dans $\Lambda(-i)$. On a unifié sur la seconde convention ; les torsions
de ces trois pages sont donc de signe opposé à celui de la page. Une note de coin
ajoute : un sous-motif d'un motif effectif est effectif.}

\pagerange{154}{154}
\subsection*{Le niveau, troisième mouture (page 154)}

Un motif effectif $M$ de degré $i$ est de niveau $\leqslant i - kp$ s'il est
isomorphe à un $M'(-k)$ avec $M'$ effectif de degré $i - kp$. Ces motifs forment une
sous-catégorie épaisse de $\mathcal{M}^{i}$, et comme $\mathcal{M}$ est semi-simple,
tout objet se décompose canoniquement en la somme de son plus grand sous-objet de
niveau $\leqslant i - kp$ et d'un supplémentaire. D'où
\[
\mathcal{M}^{i} \;\simeq\; \mathcal{P}^{i} \times \mathcal{P}^{i-p}(-1) \times
\mathcal{P}^{i-2p}(-2) \times \cdots ,
\]
où $\mathcal{P}^{j}$ est la catégorie des motifs effectifs \emph{primitifs} de degré
$j$, n'ayant aucun sous-motif non nul de niveau $\leqslant j - p$. La structure
multiplicative se lit alors par les foncteurs $\mathcal{Q}^{i,j,k} \colon
\mathcal{P}^{i} \times \mathcal{P}^{j} \to \mathcal{P}^{i+j-kp}$, $(M, N) \mapsto
K_{i+j,k}(M \otimes N)(k)$, où $K_{i+j,k}$ est la projection de $\mathcal{M}^{i+j}$ sur
son facteur $\mathcal{P}^{i+j-kp}(-k)$.\footnote{La page écrit $K_{i,k}$ ; la
projection porte sur $M \otimes N$, de degré $i+j$. La lettre ronde des catégories
primitives et celle des foncteurs sont lues $\mathcal{P}$ et $\mathcal{Q}$ sans
certitude.} C'est, dans le cadre abstrait, la décomposition primitive par le niveau
de la page 2 et des pages 37 à 40. Le niveau d'un motif quelconque se définit en se
ramenant au cas effectif.

\pagerange{155}{157}
\subsection*{Duaux, formes bilinéaires, involutions (pages 155 à 157)}

On suppose désormais $p = 2$, et que pour tout motif effectif $M$ de degré $i$ le
$\underline{\operatorname{Hom}}(M, \Lambda(-i))$ existe dans les motifs effectifs, le
$\underline{\operatorname{Hom}}$ étant défini par l'adjonction
$\operatorname{Hom}(P, \underline{\operatorname{Hom}}(M, N)) \simeq \operatorname{Hom}(P
\otimes M, N)$. Si $M$, $N$ sont effectifs de degrés $i$, $j$ et que
$\underline{\operatorname{Hom}}(M, N)$ existe, il est de degré $j - i$ ; ainsi
\[
M^{*} = \underline{\operatorname{Hom}}(M, \Lambda(-i)) , \qquad \deg M^{*} = \deg M = i ,
\]
et $\underline{\operatorname{Hom}}(M, N(-i)) \simeq M^{*} \otimes N$ pour tout $N$.
Donc $\underline{\operatorname{Hom}}(M, N(-j))$ existe pour $j$ grand, puis
$\underline{\operatorname{Hom}}(M, N)$ dans $\widetilde{\mathcal{M}}$, avec les règles
évidentes ; on pose $\check{M} = \underline{\operatorname{Hom}}(M, \Lambda(0))$, et
$\underline{\operatorname{Hom}}(M, N) \simeq \check{M} \otimes N$, en prenant garde
que $\check{M}$ n'est pas effectif si $\deg M \neq 0$. La catégorie est
rigide.\footnote{La page note les $\underline{\operatorname{Hom}}$ par une majuscule
cursive grasse ; elle écrit ici $\widehat{\mathcal{M}}$, et $\overline{\mathcal{M}}$
page 158, pour ce que la page 153 note $\widetilde{\mathcal{M}}$ ; c'est la même
catégorie.}

Les \emph{formes bilinéaires} sur deux motifs $M$, $N$ de degré $i$ sont
\[
B(M, N) = \operatorname{Hom}(M \otimes N, \Lambda(-i)) \simeq \operatorname{Hom}(M, N^{*})
\simeq \operatorname{Hom}(N, M^{*}) ,
\]
avec, si $M = N$, les formes symétriques et alternées. Une forme est non dégénérée en
$M$ si $M \to N^{*}$ est un monomorphisme, ce qui équivaut à ce que $N \to M^{*}$ soit
un épimorphisme ; si $M = N$, les deux notions coïncident.\footnote{La page marque
cette dernière affirmation d'un point d'interrogation. Elle est juste : la dualité
est exacte et échange monomorphismes et épimorphismes, et $M$, $M^{*}$ ayant même
longueur, un monomorphisme $M \to M^{*}$ est un isomorphisme.}

\emph{Involutions.} Soient $\alpha \in B(M, M)$ et $\beta \in B(N, N)$ non dégénérées,
symétriques ou alternées, d'où $\tilde{\alpha} \colon M \xrightarrow{\sim} M^{*}$ et
$\tilde{\beta} \colon N \xrightarrow{\sim} N^{*}$. L'\emph{adjoint} de $u \colon M \to N$
est
\[
u' = \tilde{\alpha}^{-1} \circ {}^{t}u \circ \tilde{\beta} \colon N \longrightarrow M ,
\]
et l'on a $(u \pm v)' = u' \pm v'$, $(uv)' = v'u'$, $1' = 1$, et $(u')' = u$ si $\alpha$
et $\beta$ sont de même type. En particulier $\operatorname{End}(M)$ est une
$\mathbb{Q}$-algèbre involutive. Si l'on remplace $\alpha$, $\beta$ par $\alpha_{1}$,
$\beta_{1}$ avec $\tilde{\alpha}_{1} = \tilde{\alpha} A$ et $\tilde{\beta}_{1} =
\tilde{\beta} B$, le nouvel adjoint est
\[
u'' = A^{-1}\,u'\,B ,
\]
et pour $M = N$, $A = B$ autoadjoint, l'involution devient $u \mapsto A^{-1}u'A$ :
elle est déterminée par la forme modulo un automorphisme intérieur.\footnote{Le
diagramme de la page porte une seule étiquette sur une flèche à deux pointes, pour
$u$ et $u'$ ; on a écrit la formule. La même règle, $u^{\psi} = a^{-1}u^{\varphi}a$,
ouvre les « Notes anciennes » du dossier 10, pages 20 à 23.}

\pagerange{158}{166}
\subsection*{Formes définies positives (pages 160, 158, 162, 164 et 166)}

Le feuillet 160 pose les axiomes que le feuillet 158 emploie déjà : on les lit dans
cet ordre. On suppose donné, pour tout $X \in \mathcal{V}$, avec $\underline{h}(X) =
(\underline{h}^{i}(X))_{i}$, un ensemble $\Phi_{X}$ de \emph{formes fondamentales}
$\varphi = (\varphi_{i})$, $\varphi_{i} \in B(\underline{h}^{i}(X), \underline{h}^{i}(X))$,
tel que :
\begin{enumerate}
  \item[(i)] si $\varphi \in \Phi_{X}$ et $\psi \in \Phi_{Y}$, $\varphi \oplus \psi \in
  \Phi_{X \sqcup Y}$ ;
  \item[(ii)] si $\varphi \in \Phi_{X}$ et $\psi \in \Phi_{Y}$, $\varphi \otimes \psi \in
  \Phi_{X \boxtimes Y}$ ;
  \item[(iii$_{0}$)] sur $\underline{h}^{0}(X)$, les $\varphi \in \Phi_{X}$ sont définies
  positives ;
  \item[(iii)] si $\varphi \in \Phi_{X}$, $\psi \in \Phi_{Y}$ et $u \colon X \to Y$
  induit $u^{(i)} \colon \underline{h}^{i}(Y) \to \underline{h}^{i}(X)$, d'adjoint
  $s^{(i)}_{u}$ relativement à $\psi$ et $\varphi$, alors
  \[
  \operatorname{Tr}\bigl(s^{(i)}_{u}\,u^{(i)}\bigr) \geqslant 0 , \quad
  \text{avec égalité seulement si } u^{(i)} = 0 .
  \]\footnote{La page écrit $u^{(i)} \colon h^{i}(Y) \to h^{i}(Y)$ ; la
variance de $\underline{h}$ demande $h^{i}(Y) \to h^{i}(X)$, et l'adjoint va en sens
inverse. La condition sur $\underline{h}^{0}$ est écrite une première fois,
surchargée, sous un numéro qu'on lit « (iii)° ».}
\end{enumerate}

C'est la positivité de Weil, que la conjecture standard de type Hodge doit donner
pour les formes $\varphi_{i}(x, y) = \pm\,\langle L^{n-i}x, y\rangle$ corrigées sur
chaque morceau primitif.\footnote{« Standard conjectures on algebraic cycles »
(1969) déduit de la conjecture de type Hodge la positivité de
$\operatorname{Tr}(u\,{}^{t}u)$ au sens de Weil. En caractéristique nulle, la
conjecture de type Hodge est un théorème, par la théorie de Hodge ; en
caractéristique $p$, elle n'est connue que pour les surfaces (Segre, Grothendieck)
et quelques cas particuliers.} La page 158 en tire que les
$\operatorname{End}(\underline{h}^{i}(X))$ sont des \emph{algèbres de Riemann} ---
des algèbres semi-simples munies d'une involution positive ---, donc semi-simples, ce
qui permet de construire la catégorie des motifs comme catégorie abélienne. Elle
discute l'indétermination du système : remplacer $\varphi_{i}$ par
$\varepsilon_{i}\varphi_{i}$, $\varepsilon_{i} = \pm 1$, ne change pas les adjoints,
donc respecte (i) et (iii) ; (iii$_{0}$) impose $\varepsilon_{0} = 1$ et (ii)
$\varepsilon_{i}\varepsilon_{j} = \varepsilon_{i+j}$, donc $\varepsilon_{i} =
\varepsilon^{i}$ : il reste un signe, qu'on fixe en degré $1$ en demandant, en
géométrie, que les formes correspondent aux polarisations des variétés
abéliennes.\footnote{La fin de la page 158 est d'une écriture très rapide, et
l'enchaînement est celui que la transcription propose. Son premier mot, lu « Lui »,
suggère une note à l'usage d'un correspondant, que la page ne nomme pas.}

\emph{(iii bis).} L'axiome (iii) équivaut, moyennant (ii), (iii$_{0}$) et la dualité
qui fait des morphismes $X \to Y$ des éléments de $\Gamma\,\underline{h}^{2i}(X
\boxtimes Y)(i)$, à ceci : \emph{sur $\operatorname{Hom}(\Lambda(-i),
\underline{h}^{2i}(X)) \simeq \Gamma\,\underline{h}^{2i}(X)(i)$, la forme induite par
$\varphi_{2i}$ est définie positive.}\footnote{La page écrit $\varphi_{i}$ pour la
forme sur $\underline{h}^{2i}(X)$ ; dans son indexation, c'est $\varphi_{2i}$. La
forme induite sur deux classes $a$, $b$ est $\varphi_{2i} \circ (a \otimes b) \in
\operatorname{End}\Lambda(-2i) = \mathbb{Q}$.} C'est, pour les cycles algébriques de
codimension $i$, la forme même de la conjecture de type Hodge.

Soient maintenant $\varphi, \psi \in \Phi_{X}$, avec $\tilde{\psi} = \tilde{\varphi} A$.
Pour $u \in \operatorname{End}(\underline{h}^{i}(X))$, l'adjoint relatif à $(\varphi,
\psi)$ est $u'' = u'A$, et (iii) donne $\operatorname{Tr}(u u' A) \geqslant 0$ pour
tout $u$ : $A$ est dans le cône positif de l'algèbre de Riemann
$\operatorname{End}(\underline{h}^{i}(X))$ définie par $\varphi$, c'est-à-dire que
$\varphi$ et $\psi$ définissent la même structure riemannienne.\footnote{Le cône des
$\sum u'u$ étant autopolaire pour la forme trace, la condition dit que $A$ est dans
son adhérence, et à l'intérieur puisqu'il est inversible. Le dossier 10, pages 24 à
30, développe cette géométrie des cônes.}

Plus généralement (page 164), soit $M$ effectif de degré $i$, et deux plongements
$\lambda \colon M \to \underline{h}^{i}(X)$, $\mu \colon M \to \underline{h}^{i}(Y)$,
avec les formes induites $\varphi_{0}$, $\psi_{0}$ par $\varphi \in \Phi_{X}$, $\psi \in
\Phi_{Y}$. Alors $\varphi_{0}$ et $\psi_{0}$ sont définies et équivalentes. En effet
$\lambda'\lambda = \mathrm{id}_{M}$ et $\mu'\mu = \mathrm{id}_{M}$ par définition des
formes induites, et pour $u \in \operatorname{End}(M)$, d'adjoint $u'$ relatif à
$(\varphi_{0}, \psi_{0})$,
\[
\operatorname{Tr}(\mu u \lambda')'(\mu u \lambda') =
\operatorname{Tr}\lambda u' \mu' \mu u \lambda' =
\operatorname{Tr}\lambda'\lambda\, u' \mu'\mu\, u = \operatorname{Tr} u'u ,
\]
qui est donc $\geqslant 0$, nul seulement si $\mu u \lambda' = 0$, c'est-à-dire si $u =
0$.\footnote{La page écrit $\underline{h}^{j}(Y)$ sur le croquis, $\underline{h}^{i}(Y)$
dans le texte ; on lit $i$. Le passage de la nullité de $\mu u \lambda'$ à celle de
$u$ utilise que $\mu$ est injectif et que $\lambda'$ est surjectif, $\lambda'\lambda$
étant l'identité ; la page dit seulement « OK ! ».} On obtient ainsi, sur tout motif
effectif $M$, une classe $\Psi_{M}$ de formes riemanniennes, qui satisfait les
analogues de (i) à (iii bis).

Enfin (page 166) : on a $B(M(-1), M(-1)) \simeq B(M, M)$, et il faut vérifier que
$\Psi_{M(-1)}$ correspond à $\Psi_{M}$ --- « cela doit pouvoir se faire ». Par (ii), il
suffit que la forme canonique $1 \in B(\Lambda(-1), \Lambda(-1))$ soit dans
$\Psi_{\Lambda(-1)}$, vérification triviale dans le cas classique. Ceci permet
d'étendre la notion de forme définie positive aux motifs non nécessairement
effectifs, et c'est la dernière phrase du dossier.\footnote{Un paragraphe encadré de
la page 166, sur une explicitation directe par l'objet $\Lambda(-1)$, est barré de
trois diagonales. La notion obtenue est ce qu'on appellera une \emph{polarisation}
de la catégorie tannakienne des motifs (Saavedra, 1972 ; Deligne--Milne, 1982) ;
les feuillets ne la nomment pas ainsi.}

\end{document}
